\section{IBL and homotopy transfer}
\label{ch:ibl}

This chapter reads the cubic action in the language of algebras with operations of several inputs and several outputs, the \IBL{} algebras of Cieliebak, Fukaya and Latschev, and reads the elimination of the auxiliary fields as a homotopy transfer.
It expands the section of the paper on the \IBL{} perspective and the subsection of the all-order section on integrating out as a homotopy transfer.
Both are reproduced as they stood before the paper was condensed, followed by three earlier notes: a companion note on the bookkeeping of the raised products and of the \IBL{} relations, the section of the overview note on the \IBL{} structure, and the sections of a pedagogical note on the two presentations of the quartic structure and on its physical grading.
Only this chapter contains the two record-only sections at its end: the ghost numbers of the kinetic operators and the reading of the canonical-model theorem, which correct several wordings of the earlier texts, and the Maurer--Cartan form of the raised family, with its components, the coupling it requires and the components of its gauge transformation.

\paragraph{Assumptions and conventions.}

Everything in this chapter is at genus zero, in the physical ghost number, with the coupling $\gs$ of ghost number zero.
Input strands of a product sit at ghost number two and output strands at ghost number three, so that $\Ngh(\sprod{m}{n}) = 3 n - 2 m$.
The algebraic statements are statements about symmetric multilinear forms on a finite-dimensional graded space, which for the string vertices is a level truncation; the \IBL{} definitions and the theorems of~\cite{Cieliebak:2015:HomologicalAlgebraRelated} are quoted, not re-derived.
The products studied are those read from the fundamental vertices; from the quintic order on, the cubic action uses reduced vertices, whose readings with several outputs are not studied here.

\subsection*{The \IBL{} section of the paper}

The first block is the section of the paper on the \IBL{} perspective before its condensation.
It holds the following results.

\noindent\status{Proven} The physical ghost number of the operations, $\Ngh(\sprod{k}{\ell}) = 3 \ell - 2 k$ at genus zero, $3 \ell - 2 k + g$ at genus $g$, the genus term derived at $g \le 1$ and taken from the ghost-additive self-gluing at $g \ge 2$.\record{009-a3-ibl-pedagogical (R23)}

\noindent\status{mixed} The family raised with the inverse of the odd form $\Omega$ is the family of the paper twisted by the parity of the inputs, the two carrying the same information, although on inputs of mixed parity they are different maps; the \IBL{} generating function holds for each family in its own form, and the theorem below holds in the generating-function form under the coefficient rule of the paper.\record{088-ibl-generating-sign (R87)}

\noindent\status{Proven} The raised family with the operations with no input or no output included satisfies every genus-zero \IBL{} relation if and only if the $\ell_n$ form a cyclic \Loo{} algebra; the degree bookkeeping needs $d = \tfrac{1}{2}$, outside the integer values of~\cite{Cieliebak:2015:HomologicalAlgebraRelated}.\record{026-linfty-to-ibl-repackaging (R31)}

\noindent\status{mixed} Without those operations the relation at signature $(1, 2)$ holds, and so does every relation with one input or one output, but the relation at $(2, 2)$ is not implied by the \Loo{} relations, as a statement on general cyclic \Loo{} algebras; whether the closed-string products satisfy it is not evaluated.\record{026-linfty-to-ibl-repackaging (R16)}

\noindent\status{mixed} The braces of a parity-inhomogeneous differential operator $\mathcal D$ form a curved \Loo{} algebra if and only if the brace criterion~\eqref{p:eq:ibl:brace-criterion} holds; the family with two or more outputs is not the brace structure of any such operator within the hosts considered.\record{026-linfty-to-ibl-repackaging (R32)}

\noindent\status{mixed} The map $\Psi \mapsto (\Psi, \Psi^2)$ is not an \Loo{} morphism off the classical complex: the morphism identity with two inputs fails by $\hat{Q}_B (\Psi^2)$, within strict maps whose linear part is the inclusion, on the even-input sectors of a finite model.\record{029-morphism-check (R36)}

\noindent\status{Conditional(PH1) -- Verified (finite-dimensional model)} On the classical complex the elimination of $\Sigma_2$ is a homological-perturbation transfer, with transferred brackets $(Q_B, \gs \ell_2, \gs^2 \ell_3, 0)$ and transfer map $\Psi \mapsto (\Psi, \Psi^2)$, a cyclic \Loo{} quasi-isomorphism, under the invertibility of $\sprod{2}{2}$ on the quotient by its kernel, level by level.\record{037-hpt-transfer-sigma (R41)}

A later check of this section corrects several of its wordings, and the paragraph on homotopy transfer is replaced by an override in the block; the other corrections are the following.\record{114-ibl-remarks-check (R111)}
\begin{itemize}
  \setlength{\itemsep}{0.3ex}

  \item The loop homotopy algebras of~\cite{Markl:2001:LoopHomotopyAlgebras} are \Loo{} algebras with a second-order operator of Batalin--Vilkovisky type; they are not presented as \IBL{} algebras, and the inverse pairing alone does not make them one.
    For \IBL{} structures in string field theory the references are~\cite{Cieliebak:2015:HomologicalAlgebraRelated, Markl:2017:MVFormalismIBL} and Münster--Sachs~\cite{Muenster:2013:QuantumOpenClosedHomotopy}.

  \item The first of the two \enquote{features} is correct in its facts and misleading in its framing.
    The statement it should make is the ghost-number balance of \cref{ch11:sec:transfer}: the kinetic operator of the $n$-string component, of ghost number $2 n$, has ghost number $n$, and no operator acting strand by strand with ghost number one can appear; in an \IBL{} algebra, by contrast, the differential acts on every strand of a multi-string state, so the auxiliary fields of the cubic action are not the multi-string components of an \IBL{} field.
    The sentence on the inverse pairing as the only object with two outputs describes the loop homotopy algebra, not the \IBL{} framework.

  \item In the second feature, the auxiliary nature of $\Sigma_2$ and the parity count of the $2$-string pairing are correct, the count under the rule that makes the total multi-string field even.
    The Batalin--Vilkovisky structure that an \IBL{} algebra defines is a different notion, a higher-order differential operator on the symmetric algebra, which needs no symplectic form, so the paragraph is not contradicted by it.
    The stronger reason why the cubic action is not of that form is the Maurer--Cartan comparison of \cref{ch11:sec:mc}, within the class of local redefinitions considered there.

  \item The last sentences of that paragraph draw a contrast that does not exist.
    The quadratic form of $\Sigma_2$ is the quartic vertex itself, $\pair[2]{\Sigma}{\sprod{2}{2} \Sigma'} = \vertex{4}(\Sigma, \Sigma')$, and the self-adjointness of $\sprod{2}{2}$ for the $2$-string pairing is equivalent to the pair-exchange symmetry of $\vertex{4}$, exactly as the cyclicity of $\ell_n$ is the symmetry of $\vertex{n + 1}$ read through the pairing.

  \item The statement that the raised family is not an \IBL{} algebra in the sense of the definition of~\cite{Cieliebak:2015:HomologicalAlgebraRelated} means first that its relations include the operations with no input or no output, which that definition excludes; the degree $d = \tfrac{1}{2}$ is a second point.
    That value is forced: the strand degree $3 - \Ngh$ is the only one affine in the ghost number that reproduces the degrees of~\cite{Cieliebak:2015:HomologicalAlgebraRelated}, and it gives $d = \tfrac{1}{2}$, for which their genus-counting variable would be odd (the derivation is~\eqref{ped:eq:cflmatch} of the pedagogical note below).
\end{itemize}

\noindent\status{no ledger entry} The value $d = \tfrac{1}{2}$ is a property of the raised family, not of closed string field theory.\record{120-muenster-degree-dictionary}
Münster and Sachs~\cite{Muenster:2013:QuantumOpenClosedHomotopy} use no integer $d$ and give the loop parameter no degree: they shift the closed string space by two, so that the string field has degree zero and the odd symplectic form degree $-1$, and every structure map of their loop homotopy algebra, the products $l_n^g$ and the inverse symplectic form, has degree one.
In the bookkeeping of~\cite{Cieliebak:2015:HomologicalAlgebraRelated} this is $d = 0$ with the cohomological sign of the grading, and the operations with no input place their algebra in the extension of section~8 of that reference, not in its Definition~2.3.
They never read a vertex with several outputs: in their grading the product $\sprod{m}{n}$ has degree $n$, the degree of $l_{m + n - 1}$ composed with $n - 1$ copies of the inverse symplectic form.
No strand degree affine in the ghost number makes the degree of $\sprod{m}{n}$ independent of both $m$ and $n$, so a family that contains the vertices read with several outputs cannot have a uniform degree, and matching the form of~\cite{Cieliebak:2015:HomologicalAlgebraRelated}, independent of the outputs and linear in the inputs, forces $d = \tfrac{1}{2}$.
Their shift by two is even and brings no Koszul sign, their odd form is the form $\Omega$ of the first block up to an overall sign that their text does not fix, and their raising by its inverse is the family $p^{\Omega}_{m, n}$, equal to $\sprod{m}{n}$ on inputs of even parity.
The two sentences of the first block on $d = \tfrac{1}{2}$ are statements about~\cite{Cieliebak:2015:HomologicalAlgebraRelated} and stand; the degree count is machine-checked, including a check that no integer $d$ fits.

\subsection{\IBL{} perspective}
\label{p:sec:ibl}

The construction of \cref{p:sec:hs} reads the string vertices with several of their slots as outputs.
Operations with several inputs and several outputs are the defining data of an \IBL{} algebra, the homotopy version of an involutive Lie bialgebra: a bracket and a cobracket with their higher homotopies, involutive meaning that the bracket composed with the cobracket vanishes~\cite{Cieliebak:2015:HomologicalAlgebraRelated, Markl:2017:MVFormalismIBL}.
  In string field theory it appears in the loop homotopy algebra of the quantum theory~\cite{Markl:2001:LoopHomotopyAlgebras, Muenster:2013:QuantumOpenClosedHomotopy}, whose only operation with two outputs is the inverse pairing, as discussed below.

In this section, we write the cubic action in that language.
After studying the lowest product $p_{2,2}$ in detail to gain some intuition, we generalize to higher orders.
In particular, from the products $\ell_n$ alone, one builds a whole family $\sprod{m}{n}$ by applying the inverse pairing to some of their slots.
We ask whether the $\Loo$ relations of the $\ell_n$ imply the \IBL{} relations for that family, and, conversely, whether such a family defines an $\Loo$ structure on the space of multi-string fields.
Note that we write everything in terms of the CFT data of \cref{p:app:conventions}, with physical ghost numbers and without suspension.


\subsubsection{Operations with several outputs}
\label{p:sec:ibl:operations}

An $\Loo$ algebra is a set of operations $\ell_n$ with $n$ inputs and one output, the string products of the original theory (\cref{p:sec:hs:action}).
The cubic action contains instead the product $\sprod{2}{2} \colon \Hn{2} \to \Hn{2}$ of~\eqref{p:eq:hs1:p22}, with two inputs and two outputs, as the kinetic operator of $\Sigma_2$.
At higher orders the construction uses $\sprod{3}{2}$, $\sprod{3}{3}$ and the products built in the same way from the reduced vertices.
The structure whose operations $\sprod{m}{n}$ have $m$ inputs and $n$ outputs, with one master equation relating all of them, is an \IBL{} algebra~\cite{Cieliebak:2015:HomologicalAlgebraRelated}.
The single-output operations form an $\Loo$ algebra, the operation with one input and two outputs is a cobracket, and the operations with more outputs are their higher homotopies.

Such a structure also acts on multi-string states.
Each operation $\sprod{m}{n}$ is extended to the symmetric algebra $\bigoplus_{n \ge 1} \Hn{n}$ by acting on every subset of $m$ strands of a multi-string state and leaving the other strands untouched.
The master equation is an identity between operators on that algebra~\cite{Cieliebak:2015:HomologicalAlgebraRelated, Muenster:2013:QuantumOpenClosedHomotopy}.
The field of the theory is accordingly a sum of components with any number of strands~\cite[sec.~9]{Cieliebak:2015:HomologicalAlgebraRelated}.
This is the field content of the cubic action: one component per number of strands, $\Psi$ with one strand and the auxiliary field $\Sigma_n$ with $n$ strands.
We use the total field $\boldsymbol{\Psi} = \sum_{n \ge 1} \gs^{n-1} \, \Psi_n$ of~\eqref{p:eq:all:total-field}, with $\Psi_1 = \Psi$ and $\Psi_n = \Sigma_n$ for $n \ge 2$, in terms of which every cubic term of the action carries one power of $\gs$ and every kinetic term none (\cref{p:sec:all-orders}).
  The powers of $\gs$ are the rescaling of the auxiliary field met at quartic order, $\Sigma_2 \to \Sigma_2 / \gs$, extended to the whole tower (compare the expansion parameter of~\cite[sec.~IV.B]{Muenster:2013:QuantumOpenClosedHomotopy}).

\subsubsection{Cubic action on the symmetric algebra}
\label{p:sec:ibl:symmetric-algebra}

Each $n$-string space carries the pairing $\pair[n]{\cdot}{\cdot}$, with one $c_0^-$ per strand, which is graded-symmetric for every $n$ and non-degenerate.
The products $\sprod{m}{n}$ are defined from the vertices by reading $n$ slots through it, as in~\eqref{p:eq:conv:sprod}.
The sphere correlator saturates at ghost number six and each pairing at five.
Hence an output strand of a product sits at ghost number three and an input strand at two, and we get:
\begin{equation}
  \label{p:eq:ibl:ghost-numbers}
  \Ngh(\sprod{m}{n})
    = 3 n - 2 m,
\end{equation}
where in particular $\Ngh(\sprod{2}{2}) = 2$.
Reading one more slot of the same vertex as an output raises the ghost number by five: an input slot at ghost number two becomes an output slot at ghost number three.
Here the physical grading differs from the one of~\cite{Cieliebak:2015:HomologicalAlgebraRelated}, in which the degree of an operation does not depend on its number of outputs.
Assigning the degree $3 - \Ngh$ to every strand turns the ghost number of $\sprod{m}{n}$ into the degree of~\cite{Cieliebak:2015:HomologicalAlgebraRelated} with $d = \tfrac{1}{2}$, a value their definition excludes.
We use that dictionary only to compare with their definition below, and otherwise work with the self-contained ghost-number balance~\eqref{p:eq:ibl:ghost-numbers}.
At this stage, the question of whether these products satisfy the $\IBL$ relations is still open.

The cubic action is a function on the symmetric algebra.
At quartic order, with the components $\Psi_1 = \Psi$ and $\Psi_2 = \Sigma_2$ of~\eqref{p:eq:all:total-field}, the action~\eqref{p:eq:hs4:action-tilde} reads:
\begin{equation}
  \label{p:eq:ibl:action-components}
  \begin{aligned}
    \caction[\le 4]
      &
      = \frac{1}{2} \, \pair{\Psi_1}{Q_B \Psi_1}
        - \frac{1}{4!} \, \pair[2]{\gs \Psi_2}{\sprod{2}{2}(\gs \Psi_2)}
      \\ &
      \quad
      + \frac{\gs}{3!} \, \pair{\Psi_1}{\ell_2(\Psi_1^2)}
        + \frac{2 \gs}{4!} \, \pair[2]{\gs \Psi_2}{\sprod{2}{2}(\Psi_1^2)},
  \end{aligned}
\end{equation}
using the definition of $\sprod{2}{2}$ (\cref{p:sec:hs:action}).
Every term is the pairing, on some $\Hn{n}$, of one component $\gs^{n-1} \Psi_n$ of $\Psi_{\mathrm{tot}}$ with a product applied to other components: this is the form of an action built from the operations $\sprod{m}{n}$ and the multi-string pairings.
Two features distinguish it from the actions of the \IBL{} type found in the literature.

First, the kinetic operator of the $2$-string component is $\sprod{2}{2}$, of ghost number two, and not the BRST operator: a term $\pair[2]{\Sigma_2}{\hat Q_B \Sigma_2}$, with $\hat Q_B$ the BRST operator acting on either strand, vanishes by ghost-number conservation: the $2$-string pairing is nonzero only when the ghost numbers paired on each strand add up to five, ten in total, while $\Sigma_2$ and $\hat Q_B \Sigma_2$ carry $4 + 5 = 9$.
In the loop homotopy algebra of closed string field theory~\cite{Markl:2001:LoopHomotopyAlgebras, Muenster:2013:QuantumOpenClosedHomotopy, Muenster:2014:HomotopyClassificationBosonic}, by contrast, the differential on every component of the symmetric algebra is the extension $\hat Q_B$ of the BRST operator.
The only object with two outputs in that algebra is the inverse pairing itself, which enters at one loop and carries no vertex.
The genus-zero part of its master equation is the $\Loo$ relations together with the invariance of the inverse pairing, and no product with two outputs built from a vertex appears in it.
The cubic action is a classical action with such a product, and it is a different object.

Second, the $2$-string component is auxiliary.
Its equation of motion, $\sprod{2}{2}(\Sigma_2) = \sprod{2}{2}(\Psi^2)$, is a field equation built from $\sprod{2}{2}$ alone, and it fixes $\Sigma_2 = \Psi^2$ at this order on the quotient~\eqref{p:eq:hs4:quotient} because $\sprod{2}{2}$ is injective there.
The two sectors differ in their Batalin--Vilkovisky structure.
The single-string pairing has the odd counterpart $\Omega$ of~\eqref{p:eq:conv:odd-symplectic}, which defines the antibracket.
With the parity given by the ghost number modulo two, the $2$-string pairing between $\Sigma_2$ and the tensors of ghost number six is even, so that no single odd symplectic form covers both sectors (\cref{p:app:conventions}).
  The master action of \cref{p:sec:gauge} nevertheless pairs $\Sigma_2$ with an odd antifield $\Sigma_2^+$, because there the Grassmann parity is read from the target-space ghost number, not from the worldsheet ghost number modulo two.
The symmetry of the $\Sigma_2$ kinetic term under the exchange of its two arguments is the pair-exchange symmetry of the quartic vertex, a property of the vertex rather than of the pairing.
This symmetry makes the completion of the square of \cref{p:sec:hs:quartic} possible.

\subsubsection{General correspondence}
\label{p:sec:ibl:general}

We forget for now the auxiliary fields and ask what the products $\ell_n$ of the original theory determine by themselves.
Let
\begin{equation}
  \label{p:eq:ibl:inverse-pairing}
  \Delta
    := \pair{\cdot}{\cdot}^{-1}
\end{equation}
be the inverse pairing, an element of $\Hn{2}$ whose two strands have ghost numbers adding up to five.
Feeding one strand of $\Delta$ into an input of $\ell_n$ and keeping the other strand as a second output defines
\begin{equation}
  \label{p:eq:ibl:raising}
  \sprod{n-1}{2}
    = \Delta \circ \ell_n,
\end{equation}
and iterating produces every $\sprod{m}{n}$ from the $\ell_n$ alone.
This is the content of the definition~\eqref{p:eq:conv:sprod}: reading $n$ slots of $\vertex{m+n}$ through the $n$-string pairing is raising $n$ indices with $\Delta$.

The raising that closes the antibracket uses instead the inverse of the odd symplectic form $\Omega$ of~\eqref{p:eq:conv:odd-symplectic}, which differs from $\Delta$ by a sign $-1$ on every slot of the vertex that holds an odd state.
Raising every slot in this way defines the family $p^{\Omega}_{m, n}$:
\begin{equation}
  \label{p:eq:ibl:raised-family}
  \vertex{m+n}(A_1, \ldots, A_n, B_1, \ldots, B_m)
    =: \Omega_n \bigl( A_1 \otimes \cdots \otimes A_n, p^{\Omega}_{m, n}(B_1, \ldots, B_m) \bigr),
\end{equation}
where $\Omega_n(X, Y) := (-1)^{\abs{X}} \, \pair[n]{X}{Y}$ for $X$ of definite parity, extended by linearity.
We take $\Omega_n$ as the $n$-th tensor power of $\Omega$, with the Koszul sign of~\eqref{p:eq:conv:pairing-n}.
Here $\abs{X}$ is the parity of the states, not of their coefficients, so that on a string field of unrestricted ghost number the sign acts component by component.
We write the relations below for this family.
  Since the pairing joins states of opposite parity and the vertices are even, this family differs from the family of~\eqref{p:eq:conv:sprod} by the parity of the inputs:
\begin{equation}
  \label{p:eq:ibl:raised-vs-sprod}
  p^{\Omega}_{m, n}(B_1, \ldots, B_m)
    = (-1)^{\abs{B_1} + \cdots + \abs{B_m}} \, \sprod{m}{n}(B_1, \ldots, B_m).
\end{equation}
The two families coincide on inputs whose states have even total parity, whatever the number of outputs, in particular on the classical string field, of ghost number two, and on the auxiliary fields $\Sigma_k$, of ghost number $2k$.
On inputs of odd total parity, such as a single gauge parameter, they differ by a sign, already for one output.
On the string fields of unrestricted ghost number used below, whose components carry both parities, the sign changes from component to component, and the two families are different maps.
We include in the family the vertices themselves, $p^{\Omega}_{m, 0} = \sprod{m}{0} = \vertex{m}$, read with no output, and the fully raised vertices $p^{\Omega}_{0, n} = \sprod{0}{n} \in \Hn{n}$, with no input.
The whole family is encoded in the action evaluated at the sum of two string fields $\Psi$ and $\Psi'$ of unrestricted ghost number:
\begin{equation}
  \label{p:eq:ibl:generating}
  \begin{aligned}
    S(\Psi + \Psi')
      &
      = \sum_{\substack{m, n \ge 0 \\ m + n \ge 2}} \frac{\gs^{m+n-2}}{m! \, n!} \, \Omega_n \bigl( \Psi'^n, p^{\Omega}_{m, n}(\Psi^m) \bigr)
    \\
      &
      = \sum_{\substack{m, n \ge 0 \\ m + n \ge 2}} \frac{\gs^{m+n-2}}{m! \, n!} \, \pair[n]{\Psi'^n}{\sprod{m}{n}(\Psi^m)},
  \end{aligned}
\end{equation}
where each family is read through its own form and the two sums agree term by term, each term being $\tfrac{\gs^{m+n-2}}{m! \, n!} \, \vertex{m+n}(\Psi'^n, \Psi^m)$.
Here the terms with $n = 0$ are the vertices, $\tfrac{\gs^{m-2}}{m!} \, \vertex{m}(\Psi^m)$, those with $m = 0$ are $\tfrac{\gs^{n-2}}{n!} \, \vertex{n}(\Psi'^n)$, and the three terms with $m + n = 2$ are built from $\vertex{2}(A, B) = \pair{A}{Q_B B}$, with $\sprod{1}{1} = Q_B$.
Up to the factor $\gs^{m+n-2} / n!$, the product $\sprod{m}{n}$ is the $m$-th derivative at $\Psi = 0$ of the term of order $n$ in $\Psi'$, and the sum of these terms over $n$ is the $m$-point function in the background $\Psi'$.
In the framework of~\cite{Cieliebak:2015:HomologicalAlgebraRelated} an $\Loo$ algebra is an \IBL{} algebra whose operations with two or more outputs, and those of higher genus, vanish.
This embedding says nothing about the raised family, and the question is whether the raised family is an \IBL{} algebra as well.

The $\Loo$ relations of the $\ell_n$ are equivalent to the classical master equation $(S, S) = 0$ for the action $S(\Psi)$~\cite{Zwiebach:1993:ClosedStringField}, with the bracket on the string field defined by the inverse of $\Omega$.
Applied to $S(\Psi + \Psi')$, with the bracket contracting a derivative with respect to $\Psi$ against a derivative with respect to $\Psi'$, the same equation reads:
\begin{equation}
  \label{p:eq:ibl:polarized-master}
  \frac{\pd^r S(\Psi + \Psi')}{\pd \Psi^a} \, (\Omega^{-1})^{ab} \, \frac{\pd^l S(\Psi + \Psi')}{\pd \Psi'^b}
    = 0,
\end{equation}
where $\Psi = \Psi^a e_a$ and $\Psi' = \Psi'^a e_a$ in a basis $e_a$ of $\Hs$.
In this subsection $\Psi$ and $\Psi'$ are string fields of unrestricted ghost number, whose coefficients carry the Grassmann parity that makes each of them even.
Products and forms are evaluated on them with the coefficients moved to the left, each coefficient passing only the basis states to its left.
For an odd map, this differs from the standard Koszul extension, in which the map also passes each coefficient and picks up its parity.
On fields of ghost number two alone the contraction with $\Omega^{-1}$ vanishes identically, and~\eqref{p:eq:ibl:polarized-master} would hold for any products.
Since $S(\Psi + \Psi')$ depends on the sum only, \eqref{p:eq:ibl:polarized-master} is the master equation at the point $\Psi + \Psi'$, and it holds if and only if the $\ell_n$ form a cyclic $\Loo$ algebra.
Sorting~\eqref{p:eq:ibl:polarized-master} by the degrees $(m, n)$ in $(\Psi, \Psi')$ gives one relation per pair $(m, n)$.
In each of its terms the derivative with respect to $\Psi'$ frees an output strand of one product, $\Omega^{-1}$ carries it into an input of another product, freed by the derivative with respect to $\Psi$, and the remaining strands stay paired with $\Psi'$.
Read through the first form of~\eqref{p:eq:ibl:generating}, the strand enters the second product as the output of the first, since $\Omega^{-1}$ inverts $\Omega$.
The relation at $(m, n)$ is then, up to an overall sign common to all $(m, n)$, the sum of the compositions of two products of the family $p^{\Omega}_{m, n}$.
Read through the second form, the same relation joins the products $\sprod{m}{n}$ by $\Omega^{-1}$, and we write the configurations below in this form.
The relation at $(m, n)$ states that the sum over all ways of composing two products of the family through one contracted strand, with $m$ inputs and $n$ outputs in total, each ordered configuration counted once, vanishes.
These are the genus-zero \IBL{} relations in the form of~\cite[sec.~8]{Cieliebak:2015:HomologicalAlgebraRelated}, where operations with no input or no output are allowed, written in the generating-function language of their section~7.
With $\Delta$ in place of $\Omega^{-1}$, a configuration and its mirror image, in which the two products exchange their roles, would enter with opposite signs, and every relation would hold identically.

Equation~\eqref{p:eq:ibl:polarized-master} can therefore be read as follows: \emph{the family $p^{\Omega}_{m, n}$ obtained from the vertices by raising every index with $\Omega^{-1}$, with $m + n \ge 2$, satisfies the genus-zero \IBL{} relations at every $(m, n)$, in the generating-function form above, if and only if the products $\ell_n$ form a cyclic $\Loo$ algebra, in the standard Koszul convention for the odd maps $\ell_n$.}
The proof, in both directions, is the observation that~\eqref{p:eq:ibl:polarized-master} is the master equation at the point $\Psi + \Psi'$.
The argument is finite-dimensional, and for the string field it holds in the same formal sense as the master equation itself, since the inverse pairing $\Delta$ is a sum over a basis of $\Hs$.
The relations contain nothing beyond the master equation of the original theory, since the family is its Taylor expansion about a second field.

The family is not an \IBL{} algebra in the sense of~\cite[Def.~2.3]{Cieliebak:2015:HomologicalAlgebraRelated}, for two reasons.
That definition requires at least one input and one output for every operation.
Its degree bookkeeping, with an integer $d$, corresponds here to $d = \tfrac{1}{2}$ in the physical ghost number, with one suspension sign per slot when the slots of an operation are permuted.
The relations of~\cite[Def.~2.3]{Cieliebak:2015:HomologicalAlgebraRelated}, among the operations with at least one input and one output only, are a different statement.
For $\min(m, n) = 1$, read in the same generating-function form, they coincide with the ones above, and both hold.
At $(m, 1)$ they are the $\Loo$ relations.
At $(1, n)$ they are the relations among the operations with one input, which at $(1, 2)$ state that $p^{\Omega}_{1, 1}$ and the cobracket $p^{\Omega}_{1, 2}$, equal to $Q_B$ and $\sprod{1}{2}$ on inputs of even parity, are compatible.
At $(2, 2)$ they do not coincide.

The relation~\eqref{p:eq:ibl:polarized-master} at degree $(2, 2)$ is the relation $[Q_B, \ell_3] + \ell_2 \circ \ell_2 = 0$ of the $\ell_n$, written with two fields.
It involves four configurations, up to the exchange of the two products:
\begin{equation}
  \label{p:eq:ibl:configurations-22}
  \begin{aligned}
    T_1
      &
      = \tfrac{1}{2} \, \Omega(Q_B \Psi', \ell_3(\Psi, \Psi, \Psi')),
    &
    T_2
      &
      = \tfrac{1}{2} \, \Omega(\ell_3(\Psi, \Psi', \Psi'), Q_B \Psi),
    \\
    T_3
      &
      = \tfrac{1}{4} \, \Omega(\ell_2(\Psi', \Psi'), \ell_2(\Psi, \Psi)),
    &
    T_4
      &
      = \Omega(\ell_2(\Psi, \Psi'), \ell_2(\Psi, \Psi')),
  \end{aligned}
\end{equation}
where the numerical factors are those of the Taylor expansion of $S(\Psi + \Psi')$.
The part of degree $(2, 2)$ of~\eqref{p:eq:ibl:polarized-master}, divided by $2 \gs^2$, is:
\begin{equation}
  \label{p:eq:ibl:polarized-22}
  T_1 + T_2 + T_3 + \tfrac{1}{2} \, T_4
    = 0.
\end{equation}
Here the products and $\Omega$ are evaluated on $\Psi$ and $\Psi'$ with the coefficients moved to the left, each coefficient passing only the basis states to its left.
For two odd fields $L = L^a e_a$ and $K = K^b e_b$, this gives $\Omega(L, K) = \pair{e_a}{e_b} \, L^a K^b$, which is symmetric in $L$ and $K$.
The four terms are the configurations in which an output of $\sprod{2}{2}$ enters $Q_B$, an output of $Q_B$ enters $\sprod{2}{2}$, an output of $\ell_2$ enters $\sprod{1}{2}$ and an output of $\sprod{1}{2}$ enters $\ell_2$, the joined strand being raised with $\Omega^{-1}$, each written as the pairing of the two products joined by the contracted strand.
In~\eqref{p:eq:ibl:polarized-master} each of $T_1$, $T_2$, $T_3$ also appears with the roles of the two products exchanged, an output of $\sprod{0}{2}$ entering $\ell_3$, of $\sprod{1}{3}$ entering $\sprod{2}{0}$ and of $\sprod{0}{3}$ entering $\sprod{3}{0}$, with the same value, whereas $T_4$ is its own mirror image and appears once.
After this division, $T_4$ therefore enters~\eqref{p:eq:ibl:polarized-22} with half the weight of $T_1$, $T_2$, $T_3$.
The same contraction restricted to the operations with at least one input and one output contains none of the mirror terms, and the relation it imposes at degree $(2, 2)$ gives the four configurations the same weight:
\begin{equation}
  \label{p:eq:ibl:io-22}
  T_1 + T_2 + T_3 + T_4
    = 0.
\end{equation}
The difference of the left-hand sides of~\eqref{p:eq:ibl:io-22} and~\eqref{p:eq:ibl:polarized-22},
\begin{equation}
  \label{p:eq:ibl:residual-22}
  \frac{1}{2} \, \Omega(\ell_2(\Psi, \Psi'), \ell_2(\Psi, \Psi')),
\end{equation}
does not vanish for a general cyclic $\Loo$ algebra.
The cyclic Lie algebra $\alg{sl}_2 \oplus \alg{sl}_2^*$, placed at ghost numbers one and four, is a counterexample.
A cyclic differential graded Lie algebra satisfies the relation at $(2, 2)$ only when $\ell_2(\ell_2(A, B), C) = 0$ for all $A$, $B$, $C$, a condition much stronger than the Jacobi identity.

At every $(m, n)$ with $m, n \ge 2$ the relation without the operations $\sprod{k}{0}$ and $\sprod{0}{k}$ is likewise not implied by the $\Loo$ relations.
We have not evaluated whether the residual vanishes for the closed-string vertices.
It plays no role in the construction, which uses only the relations with $\min(m, n) = 1$, as we explain below.

For the converse direction, the question is which $\Loo$ structures on the symmetric algebra, with a field $\Psi_{\mathrm{tot}}$ as in~\eqref{p:eq:all:total-field}, a family of this kind produces.
Extending each $\sprod{m}{n}$ of~\eqref{p:eq:conv:sprod} to an operator $\hat{p}_{m, n}$ on the symmetric algebra, as for an \IBL{} algebra, and summing them with a formal parameter,
\begin{equation}
  \label{p:eq:ibl:D}
  \mathcal D
    := \sum_{m, n} \hbar^{m-1} \, \hat{p}_{m, n},
\end{equation}
we obtain an operator whose components have both parities, since $\hat Q_B$ is odd and $\sprod{2}{2}$ is even.
The higher derived brackets of a differential operator of finite order on a graded-commutative algebra are obtained by nesting graded commutators of the operator with the multiplications by the arguments~\cite{Akman:1997:GeneralizationsBatalinVilkoviskyAlgebras, Bering:1996:AlgebraHigherAntibrackets}.
They are the standard way of producing an $\Loo$ algebra on the symmetric algebra from a higher-order coderivation~\cite{Muenster:2013:QuantumOpenClosedHomotopy, Muenster:2014:HomotopyClassificationBosonic}.
For an odd operator they form an $\Loo$ algebra if and only if $\mathcal D^2 = 0$.

For the parity-inhomogeneous $\mathcal D$ at hand this criterion is not the right one: there are operators with $\mathcal D^2 \neq 0$ whose brackets form a curved $\Loo$ algebra.
The correct criterion requires $\mathcal D$ to act as a differential operator of finite order on a graded-commutative algebra.
Writing $F_{\mathcal D}(a)$ for its action on the exponential of an even element $a$, the brackets form a curved $\Loo$ algebra if and only if
\begin{equation}
  \label{p:eq:ibl:brace-criterion}
  \e^{-a} \mathcal D^2(\e^{a})
    = F_{\mathcal D}(a)^2,
  \qquad
  F_{\mathcal D}(a)
    := \e^{-a} \mathcal D(\e^{a}),
\end{equation}
for every even $a$, auxiliary Grassmann parameters allowed.
On single-string inputs the brackets reduce to sums of the $\sprod{n}{k}$ over the number $k$ of outputs.
For the single-output components, which we have checked with up to three inputs, the criterion then reduces at leading order in $\hbar$ to the statement above.
We have not tested the components with no output.

On multi-string inputs we do not know whether the statement above implies the criterion.
The family with two or more outputs admits no realization as a differential operator on a polynomial algebra in the strand variables, at least when the raised slots are treated alike, by maps of definite parity, and $\hbar$ is even: on a graded-commutative algebra the outputs and the inputs are symmetrized with the same parity, whereas the single contraction by the inverse pairing requires them to be symmetrized with opposite parities.
The brace construction in that form therefore cannot settle the question, and we have not tested other realizations.
We leave it open.
It is the general form of the question whether a multi-string field $\Psi_{\mathrm{tot}}$ carries an $\Loo$ structure determined by the $\ell_n$.

\subsubsection{Restricted correspondence}
\label{p:sec:ibl:restricted}


We return to the cubic action.
Its $2$-string product is the raised quartic vertex, $\sprod{2}{2} = \Delta \circ \ell_3$, and the question is which \IBL{} relations it satisfies.\footnote{The Maurer--Cartan equation of the raised family requires one more coupling, $\tfrac{\gs}{6} \, \vertex{3}(\Sigma_2, \Psi)$, which supplies the term $\gs \, \sprod{1}{2}(\Psi)$ that the original closed string field theory does not satisfy either. The factor $1/6$ is the normalization of the $2$-string component of the Maurer--Cartan element, $\Sigma_2/6$ rather than the $\Psi_2 = \Sigma_2$ of the total field: in the term with $\vertex{4}$, this component and two copies of $\Psi$, the exponential gives the coefficient $\tfrac12$ where the action has $\tfrac{2}{4!}$. Being a mixing of order $\gs$ with a field whose propagator is of order $\gs^{-2}$, this coupling changes the eliminated action already at order $\gs^0$ by the nonlocal quadratic term $\tfrac{1}{6} \pair[2]{\sprod{1}{2}(\Psi)}{\sprod{2}{2}^{-1} \sprod{1}{2}(\Psi)}$ whenever that term is finite and nonzero; we do not include it.}
At $(2, 2)$ the answer is the one found above: the relation holds with the operations with no input or no output, and without them it is not implied by the $\Loo$ relations.

In the gauge invariance of the cubic action only the relations with $\min(m, n) = 1$ enter.
The transformation of $\Sigma_2$ is that of $\Psi^2$, and the invariance to $O(\gs^3)$ is the $\Loo$ relations of the $\ell_n$ through quartic order (\cref{p:sec:hs:quartic}).
The $(2, 2)$ relation without the operations $\sprod{k}{0}$ and $\sprod{0}{k}$ never enters.
The relation with them enters only through its content, the $\Loo$ relation $[Q_B, \ell_3] + \ell_2 \circ \ell_2 = 0$ of the original theory, which the gauge invariance uses at order $\gs^2$.
Nor is any \IBL{} relation needed for $\Sigma_2$ to be a contact field.
As shown in \cref{p:sec:hs:quartic}, its auxiliary nature rests on the injectivity of $\sprod{2}{2}$ on the quotient, which is a property of the quartic form and not a relation among products.

These statements are limited to the quartic order and to the product read from the quartic vertex.
From quintic order on, the kinetic operator of $\Sigma_2$ is dressed by a background-dependent term:
  \begin{equation}
  \kin{2}{\Psi}{3}
    = \sprod{2}{2}
      - \frac{\gs}{5} \, \sprod{3}{2}(\Psi, \cdot).
\end{equation}
The quadratic terms of the higher-order auxiliary fields involve the reduced vertices $\dvertex{n}$ of \cref{p:sec:all-orders}, and the products of the construction are those vertices read with several outputs.
Whether the products of the cubic action satisfy the relations without the operations $\sprod{k}{0}$ and $\sprod{0}{k}$ is open already at $(2, 2)$.
Whether these products of the higher orders satisfy the relations, with or without these operations, is open as well.
We regard these as the main questions of the restricted correspondence.


\subparagraph{Homotopy transfer.}

Theorem~6.1 of~\cite{Cieliebak:2015:HomologicalAlgebraRelated} states that every \IBL{} algebra over a field of characteristic zero is homotopy equivalent to an \IBL{} structure on its homology with respect to the differential $\sprod{1}{1}$.
The transferred operations are given by an explicit recursion built from a splitting of the space into the homology, the image of the differential and a complement, together with a contracting homotopy.
The recursion is of the type of the homological perturbation lemma.
This is a transfer to a smaller space, the minimal model (the canonical model of~\cite{Cieliebak:2015:HomologicalAlgebraRelated}), in which every off-shell mode is integrated out, with the contracting homotopy as propagator; its genus-zero operations are the tree-level on-shell amplitudes (\cite{Kajiura:2007:NoncommutativeHomotopyAlgebras} for $A_\infty$ algebras; the construction is the same for the closed-string $\Loo$ algebra).
It is not the inverse of a strictification, which trades operations for a larger space: that inverse is the elimination of the auxiliary fields, a transfer back onto the original theory, and the minimal model is a further transfer, reached from either theory.
For $\Loo$ algebras strictification is a classical result~\cite{Loday:2012:AlgebraicOperads}, realized for closed string field theory by the construction of \cref{p:sec:bar-cobar}.
For \IBL{} algebras we know of no analogous theorem.

The map $\Psi \mapsto (\Psi, \Psi^2)$ relates the cubic action and the original theory as homotopy algebras, and its status depends on the complex on which it is considered.
Integrating $\Sigma_2$ out is a transfer of this kind for $\Loo$ algebras rather than an instance of Theorem~6.1.
The cubic action is not known to carry an \IBL{} structure, and the target, the quartic-truncated theory with its differential $Q_B$, is not a minimal model.
The transfer contracts the pair formed by $\Sigma_2$ and its equation of motion, closer to a transfer along a chain homotopy equivalence than to the passage to homology.

On the classical complex, the fields and their equations of motion, with $\Psi$ at ghost numbers two and three and $\Sigma_2$ at four and six, the elimination of $\Sigma_2$ is the homological-perturbation transfer of \cref{p:sec:all-orders}, with the transfer map $\Psi \mapsto (\Psi, \Psi^2)$, a cyclic $\Loo$ quasi-isomorphism from the quartic-truncated theory to the cubic one under the assumptions stated there.
  Off the classical complex, the same map fails the morphism identity with two inputs by~\eqref{p:eq:all:morphism-fails}.
  Both facts reflect the ghost number of $\sprod{2}{2}$.
  The map sends $\Psi$, at ghost number two, to $\Psi^2$, at ghost number four, and its linearization sends $Q_B \Psi$, at ghost number three, to $Q_B \Psi \otimes \Psi + \Psi \otimes Q_B \Psi$, at ghost number five.
  A differential on the $2$-string slot compatible with the map would therefore need ghost number one, whereas $\Ngh(\sprod{2}{2}) = 2$.
  On the classical complex there is no $2$-string state at ghost number five, so the failing term is absent and the transfer exists.
  With the brackets read off the cubic action and any grading by the total ghost number, the map does not extend to gauge parameters, at ghost number one, nor to antifields.
  It says nothing about the gauge algebra or the BRST spectrum, which are the subject of \cref{p:sec:gauge}.
  A factorization of $\sprod{2}{2}$ through an operator of ghost number $-1$ would evade the obstruction, and whether one exists is open.



\medskip
The second block is the subsection of the all-order section of the paper that states the same transfer in the terms of that section, before the condensation.
Its references to the all-order construction point to \cref{ch:all-orders}, and the finite-dimensional model it cites is in \cref{ch:finite-models}.
The status of its statement is that of the transfer on the classical complex listed above.

\subsubsection{Integrating out as a homotopy transfer}
\label{p:sec:all:homotopy-transfer}

The elimination of the auxiliary fields has an algebraic reading, which \cref{p:sec:ibl} develops and which we state here in the terms of this section.
The field $\Sigma_2$ and its equation of motion form a pair that the elimination contracts: the equation is algebraic, its linear part is $\sprod{2}{2}$, and the inverse of $\sprod{2}{2}$ on the quotient is the contracting homotopy.
Eliminating $\Sigma_2$ along this inverse is a homological-perturbation transfer, and the brackets it transfers to the string-field sector are those of the original theory truncated at the quartic order.
The three-bracket is the $\Sigma_2$-exchange tree between two quartic sources, which collapses to the contact vertex as in \cref{p:sec:hs:quartic}, and no higher bracket is generated at this order.
The transfer map is $\Psi \mapsto (\Psi, \Psi^2)$, the solution of the auxiliary equation, with no higher components.
On the classical complex, the fields and their equations of motion, it is a cyclic $\Loo$ quasi-isomorphism from the quartic-truncated theory to the cubic one, not a strict one since its quadratic component is not zero.
We verified this in a finite-dimensional model of the vertex data (\cref{p:app:finite-models}).
Beyond the model the statement assumes the quartic form invertible on the quotient by its kernel at each mass level of the level truncation and at generic total momentum, and that at such momenta the quotient and the propagator, the inverse of the form on the quotient, have a limit as the mass level grows.
We do not assume that a null vector at one mass level remains a null vector at all higher levels.

On a complex concentrated in two degrees the $\Loo$ relations and the morphism identities are vacuous beyond three statements.
These are the equality of the transferred brackets with those of the original theory, the isomorphism on cohomology, which is the invertibility of the quartic form restated, and cyclicity.
  The word ``quasi-isomorphism'' means these three statements here.
  The quasi-isomorphism controls only the solutions and their linearizations.
With the brackets read off the cubic action and any grading by the total ghost number, it does not extend to the gauge parameters or to the antifields, since there the ghost number of $\sprod{2}{2}$ obstructs it, as \cref{p:sec:ibl} explains.
The transfer was written for $\Sigma_2$ alone, and we have not examined whether the elimination of the whole tower, from the top down, is a sequence of such transfers.

Off the classical complex, the same map is not a morphism of $\Loo$ algebras, and we display the identity that fails.
Write $\tilde \ell_n$ for the brackets read off the cubic action on $\Hs \oplus \Hn{2}$, $\phi_1 = \iota$ for the inclusion of $\Hs$ and $\phi_2(\Psi, \Psi) = \Psi^2$ for the quadratic component of the map.
The single-string component of every morphism identity with even inputs holds exactly.
The $2$-string component of the identity with two inputs reads:
\begin{equation}
  \label{p:eq:all:morphism-fails}
  \Big[ \tilde \ell_1 \big( \phi_2(\Psi, \Psi) \big) + \tilde \ell_2(\Psi, \Psi) - \phi_1 \big( \ell_2(\Psi, \Psi) \big) - 2 \, \phi_2(Q_B \Psi, \Psi) \Big]_{\Hn{2}}
    = - \big( Q_B \Psi \otimes \Psi + \Psi \otimes Q_B \Psi \big),
\end{equation}
where the terms in $\sprod{2}{2}$ on the left cancel for the normalization of $\phi_2$ fixed by the auxiliary solution, and what remains is of order $\gs^0$.
The cause is that $\Sigma_2$ has no $Q_B$ kinetic term.
The differential of the cubic action on $\Hn{2}$ is proportional to $\gs^2 \sprod{2}{2}$, not to $Q_B \otimes 1 + 1 \otimes Q_B$, so nothing on the left is of order $\gs^0$ in the $2$-string slot except the last term.
No normalization of $\phi_2$, constant or $\gs$-dependent, repairs this, and no higher component of the map enters the identity with two inputs.
For generic vertices a counting argument indicates that a cubic component $\phi_3$ does not repair the identity with three inputs either, and this was not tested.
The defect is proportional to the equation of motion of $\Psi$: the map sends solutions to solutions and gauge transformations to gauge transformations through the composite law~\eqref{p:eq:hs4:gauge-sigma}, and that is its full content off-shell.
The failing term sits at ghost number five, and it is absent from the classical complex, which is why the transfer above exists there and not beyond.
This was checked in the finite-dimensional model for the even-input identities.
  We did not write the odd-input sectors.


\subsection*{The companion note on the bookkeeping}

The third block is a companion note written to fix the bookkeeping of the raised products: how one vertex defines all the maps $\sprod{k}{\ell}$, how a single contraction between two vertices composes them, which \IBL{} relations hold and why the augmented relation fails at $(2, 2)$, and three senses of the statement that \Loo{} implies \IBL{}.
Its results are the theorem and the obstruction listed with the first block, in the generating-function language.

\begin{quote}
  \small
  \textbf{Translation box: companion note on the bookkeeping $\to$ write-up.}
  \begin{itemize}
    \setlength{\itemsep}{0.3ex}
    \item $\mc V_n$, $p_{k, \ell}$, $\braket{\cdot, \cdot}_\ell$ are the paper's $\vertex{n}$, $\sprod{k}{\ell}$, $\pair[\ell]{\cdot}{\cdot}$.
    \item $\chi$ is a placeholder that marks an output leg in the generating function $H(\Psi, \chi) = S(\Psi + \chi)$; it is not a field, and not the partner field $\Upsilon_2$ of \cref{ch:partner-field}.
    \item Pairing sign: the note takes the one-string pairing graded \emph{antisymmetric} on its support, $\braket{A, B} = - \braket{B, A}$, and its inverse $\Delta$ antisymmetric.
      In the conventions of the write-up the one-string pairing is symmetric on its support, $s = + 1$ in~\eqref{p:eq:conv:pairing-symmetry}, and the odd, antisymmetric form is $\Omega$ of~\eqref{p:eq:conv:odd-symplectic}.
      The note's pairing has the exchange property of $\Omega$, so its intermediate signs are those of the family raised with an antisymmetric inverse pairing, which differs from the family of the paper by the parity of the inputs~\eqref{p:eq:ibl:raised-vs-sprod}.
      The two coincide on even inputs; the note's statements on the two-string field are stated on physical inputs, all at ghost number two, where nothing changes.
    \item Weights: the generating function weights $\sprod{k}{\ell}$ by $\gs^{k + \ell - 2}$; there is no $\hbar$.
    \item Reality: the note uses no reality condition; the write-up's convention is the star convention, $\Sigma_k^\star = (-1)^k \Sigma_k$, with the sign $-1$ for odd $k$.
  \end{itemize}
\end{quote}

\subsection{One vertex, many maps: \texorpdfstring{$\Psi$}{Psi} counts inputs, \texorpdfstring{$\chi$}{chi} counts outputs}
\label{ibk:sec:gen}

\subsubsection{The pairing, its inverse, and where the sign lives}
\label{ibk:sec:gen:pairing}

The pairing $\braket{A,B}$ is nonzero only when the ghost numbers of $A$ and $B$ add up to five, so it pairs states of opposite parity, and on that support it is antisymmetric,
\begin{equation}
  \label{ibk:eq:gen:omega}
  \braket{A,B}
    = -(-1)^{\varepsilon_A\varepsilon_B}\braket{B,A}
    = -\braket{B,A},
  \qquad
  \omega_{ab}
    := \braket{e_a,e_b}
    = -\omega_{ba}.
\end{equation}
As an operation it is odd (one insertion of $c_0^-$). The vertices $\mc V_n$ are totally graded symmetric in the physical parity, and $\mc V_n$ is nonzero only when the ghost numbers of its arguments add up to $2n$.

The inverse pairing is the bivector $\Delta$ with components $\Delta^{ab}$, $\omega_{ab}\Delta^{bc}=\delta_a{}^c$. It is antisymmetric, its components are nonzero only for $\varepsilon_a+\varepsilon_b=1$, and it carries ghost number $+5$. One consequence is used constantly: whenever $\Delta^{ab}\neq0$, the two indices have opposite parity, $(-1)^{\varepsilon_b}=-(-1)^{\varepsilon_a}$. Raising a leg of a vertex means contracting it with $\Delta$, and we do this with no extra sign, so that the string products are defined by the bare pairing identity $\braket{A_0,\ell_n(A_1,\dots,A_n)}=\mc V_{n+1}(A_0,A_1,\dots,A_n)$ and, in particular, $\ell_1=Q$ on the nose.

A sign nevertheless has to appear somewhere in this bookkeeping, and it is cleanest to isolate it before any product is written down. It is the weight carried by the classical BV antibracket. With $\Psi=\Psi^ae_a$ an even superfield and left derivatives $\pd_a=\pd/\pd\Psi^a$, the antibracket of two even functionals contains the factor $(-1)^{\varepsilon_a}$ that converts a right derivative into a left one,
\begin{equation}
  \label{ibk:eq:gen:antibracket}
  S(\Psi)
    = \sum_{n\ge2}\frac{1}{n!}\,\mc V_n(\Psi^n),
  \qquad
  (S,S)
    = (-1)^{\varepsilon_a}\,\pd_aS\,\Delta^{ab}\,\pd_bS
    = \pd_aS\,\tilde\Delta^{ab}\,\pd_bS,
\end{equation}
where the second form collects the weight into the sign-dressed bivector
\begin{equation}
  \label{ibk:eq:gen:tilde-delta}
  \tilde\Delta^{ab}
    := (-1)^{\varepsilon_a}\,\Delta^{ab},
\end{equation}
which, unlike $\Delta$, is graded \emph{symmetric} on its support. That symmetry is the whole point: without the weight, $(S,S)$ would vanish identically, because on the support of $\Delta$ the two factors $\pd_aS$ and $\pd_bS$ commute while $\Delta^{ab}$ is antisymmetric. With it, the classical master equation $(S,S)=0$ is equivalent to $(\mc H,\braket{\cdot,\cdot},\{\ell_n\})$ being a cyclic $L_\infty$ algebra; this is the standard statement of Zwiebach's construction~\cite{Zwiebach:1993:ClosedStringField}. Physically, $\pd_aS$ is the equation of motion, and \eqref{ibk:eq:gen:antibracket} says that the master equation is the statement that the equation-of-motion vector
\begin{equation}
  \label{ibk:eq:gen:eom}
  \mc E(\Psi)
    := \sum_{n\ge2}\frac{1}{(n-1)!}\,\ell_{n-1}(\Psi^{n-1})
    = Q\Psi+\tfrac12\ell_2(\Psi,\Psi)+\tfrac16\ell_3(\Psi^3)+\cdots
\end{equation}
has vanishing weighted square, $(S,S)=\tilde\omega(\mc E,\mc E)$ with $\tilde\omega(X,Y):=(-1)^{\varepsilon_c}X^c\omega_{cd}Y^d$ the pairing dressed by the same weight. The rule to carry along is short: \emph{every contraction that closes a bracket carries the sign of the strand it closes}. It is carried once per closed contraction, by $\tilde\Delta$ or equivalently by $\tilde\omega$, and never by a product.

\subsubsection{Raising legs: inputs and outputs of one vertex}
\label{ibk:sec:gen:raising}

Take one vertex $\mc V_{k+\ell}$ and decide that $k$ of its legs are inputs and $\ell$ are outputs. The inputs are fed the states $B_1,\dots,B_k$; the $\ell$ remaining legs are cut open, and what comes out is a state in $\mc H^{\otimes\ell}$. The cleanest way to say this is to say what the resulting object pairs to, exactly as $\ell_n$ is defined by what $\mc V_{n+1}$ pairs to,
\begin{equation}
  \label{ibk:eq:gen:pkl}
  \mc V_{k+\ell}(A_1,\dots,A_\ell,B_1,\dots,B_k)
    =: \braket{A_1\otimes\cdots\otimes A_\ell,\ p_{k,\ell}(B_1,\dots,B_k)}_\ell,
\end{equation}
with no sign anywhere. Here $\braket{\cdot,\cdot}_\ell$ is the $\ell$-fold pairing, which has to be fixed once because it is the one place where an ordering choice enters:
\begin{equation}
  \label{ibk:eq:gen:ell-pairing}
  \braket{A_1\otimes\cdots\otimes A_\ell,\ C_1\otimes\cdots\otimes C_\ell}_\ell
    := (-1)^{\sum_{i<j}\varepsilon_{C_i}\varepsilon_{A_j}}\;
      \prod_{i=1}^{\ell}\braket{A_i,C_i}.
\end{equation}
Strand $i$ on the left pairs with strand $i$ on the right, and the prefactor is the Koszul sign collected in interleaving the two blocks, $A_1\cdots A_\ell C_1\cdots C_\ell\mapsto A_1C_1A_2C_2\cdots$: this is $\braket{\cdot,\cdot}^{\otimes\ell}$ in the standard convention, and it is the pairing under which the multi-string field is tested strand by strand. Another choice --- reversed strand order, or no interleaving sign --- changes $p_{k,\ell}$ by a sign depending only on the output parities; the one place where that sign becomes visible is noted below.

Solving \eqref{ibk:eq:gen:pkl} for the components of $p_{k,\ell}$ is one contraction with $\Delta$ per output leg,
\begin{equation}
  \label{ibk:eq:gen:pkl-comp}
  p_{k,\ell}(B)^{c_1\dots c_\ell}
    = (-1)^{\sum_{i<j}\varepsilon_{c_i}\bar\varepsilon_{c_j}}\;
      \Delta^{c_1b_1}\cdots\Delta^{c_\ell b_\ell}\;
      \mc V_{k+\ell}\bigl(e_{b_1},\dots,e_{b_\ell},B_1,\dots,B_k\bigr),
\end{equation}
the prefactor being the interleaving sign of \eqref{ibk:eq:gen:ell-pairing} rewritten on the output indices, using $\varepsilon_{b_j}=\bar\varepsilon_{c_j}$ on the support of $\Delta$. For $\ell\le1$ the prefactor is empty and raising is plain contraction: $p_{k,1}=\ell_k$ are the string products, with $p_{1,1}=\ell_1=Q$, and $p_{k,0}=\mc V_k$ is the vertex itself read as a functional. For $\ell=2$ the prefactor is $(-1)^{\varepsilon_{c_1}\bar\varepsilon_{c_2}}$, which is $+1$ whenever the two outputs have the same parity, in particular on physical inputs; from $\ell=3$ on it is genuinely there, component by component. Each raising converts an input slot into an output slot and adds the ghost number of $\Delta$, so
\begin{equation}
  \label{ibk:eq:gen:ghost}
  N_{\rm gh}(p_{k,\ell})
    = 3-2(k+\ell-1)+5(\ell-1)
    = 3\ell-2k.
\end{equation}
The counting says which object lives where: $p_{k,1}$ has odd ghost number and is an operator, while $p_{2,2}$ has $N_{\rm gh}=+2$, the ghost number of a string field, which is why it can act as a kinetic operator on the two-string field.

The one Koszul sign that matters appears when one wants the raised leg somewhere other than in the first slot, as happens when an output of one product is fed into a non-first input of another. Graded symmetry of $\mc V$ moves the leg $e_b$ past the inputs to its left, and since $\varepsilon_b=\bar\varepsilon_a$ on the support of $\Delta$,
\begin{equation}
  \label{ibk:eq:gen:dressing}
  \begin{split}
    \Delta^{ab}\,\mc V(A_1,\dots,A_{j-1},e_b,A_{j+1},\dots)\,e_a
      &
      = (-1)^{\bar\varepsilon_a(\varepsilon_{A_1}+\cdots+\varepsilon_{A_{j-1}})}
    \\
      &
      \quad\times
        \Delta^{ab}\,\mc V(e_b,A_1,\dots,A_{j-1},A_{j+1},\dots)\,e_a.
  \end{split}
\end{equation}
The rule to remember is that an output strand commutes past inputs and other outputs with its \emph{shifted} parity $\bar\varepsilon_a$, while inputs commute among themselves with $\varepsilon$. In words: cutting a leg open flips the parity that strand carries, and the bookkeeping has to follow that flip from the moment the leg is cut.

The symmetry of the outputs follows from the same observation together with \eqref{ibk:eq:gen:ell-pairing}. The inputs are graded symmetric in the physical parity, inherited from $\mc V$. The outputs inherit the symmetry type of the pairing itself: at $\ell=2$, swapping the two output indices in \eqref{ibk:eq:gen:pkl-comp} costs the swap inside $\mc V$, $(-1)^{\varepsilon_{b_1}\varepsilon_{b_2}}=(-1)^{\bar\varepsilon_{c_1}\bar\varepsilon_{c_2}}$, times the ratio of the two interleaving prefactors, $(-1)^{\varepsilon_{c_1}+\varepsilon_{c_2}}$, and the two combine into
\begin{equation}
  \label{ibk:eq:gen:symmetry}
  p_{k,2}(B)^{c_1c_2}
    = -(-1)^{\varepsilon_{c_1}\varepsilon_{c_2}}\,p_{k,2}(B)^{c_2c_1}:
  \qquad
  p_{k,\ell}:\ \mathrm{Sym}^k_{\varepsilon}\mc H\to\Lambda^\ell_{\varepsilon}\mc H,
\end{equation}
graded symmetric in the physical parity on the inputs and graded antisymmetric in the same parity on the outputs --- the symmetry type of $\braket{\cdot,\cdot}$. This per-slot asymmetry is why no single global shift of $\mc H$ turns the bare $p_{k,\ell}$ into a map between symmetric powers of one graded space; it has to be handled slot by slot.

\subsubsection{Where the weight lives is a choice, and here is the dictionary}
\label{ibk:sec:gen:convention}

Neither of the two places where the sign of \eqref{ibk:eq:gen:tilde-delta} can be put is forced. What is forced is that the weight $(-1)^{\varepsilon_a}$ of the antibracket appear \emph{somewhere}. The convention used here puts it in the bracket, as in \cref{ibk:sec:gen:pairing}: the products are the sign-free ones of \eqref{ibk:eq:gen:pkl}, and every closed contraction carries one factor $(-1)^\varepsilon$ through $\tilde\Delta$ or $\tilde\omega$. An earlier draft of this note made the opposite choice and put the weight into the raising, defining products $\tilde\ell_n=\tilde\Delta\circ\mc V_{n+1}$ and, for several outputs, $\tilde p_{k,\ell}=\tilde\Delta^{\otimes\ell}\circ\mc V_{k+\ell}$ with the raised legs leftmost. The two families are related component by component,
\begin{equation}
  \label{ibk:eq:gen:dictionary}
  \tilde p_{k,\ell}(B)^{c_1\dots c_\ell}
    = \underbrace{(-1)^{\sum_i\varepsilon_{c_i}}}_{\text{one }-(-1)^{\varepsilon}\text{ per raised leg}}\;
      \underbrace{(-1)^{\sum_{i<j}\varepsilon_{c_i}\bar\varepsilon_{c_j}}}_{\text{interleaving sign of }\eqref{ibk:eq:gen:ell-pairing}}\;
      p_{k,\ell}(B)^{c_1\dots c_\ell}.
\end{equation}
The first factor is the expected one: a factor $-(-1)^{\varepsilon}$ per raised leg, with $\varepsilon$ the parity of the paired basis state, equivalently $(-1)^{\varepsilon}$ with $\varepsilon$ the parity of the output index. The second is not a product of one-leg signs and is an artefact of the ordering convention \eqref{ibk:eq:gen:ell-pairing}.

Two consequences matter for the physical inputs used in this note. On inputs at ghost number two, \eqref{ibk:eq:gen:ghost} forces the output ghost numbers to add up to $3\ell$ on every nonvanishing component, so $\sum_i\varepsilon_{c_i}\equiv\ell$ and the first factor of \eqref{ibk:eq:gen:dictionary} is $(-1)^\ell$ throughout. Hence $\tilde\ell_k=-\ell_k$ and $\tilde p_{k,2}=p_{k,2}$: the one-output products differ by an overall sign, the two-output ones do not. In particular the two-string kinetic operator $p_{2,2}$ is the same object in both conventions. What the earlier choice bought was that every contraction closing an antibracket becomes a bare index sum with no leftover weight, and that the outputs are uniformly graded symmetric in the shifted parity $\bar\varepsilon$ for all $\ell$; what it cost was a dressed $\tilde\ell_1=(-1)^{\varepsilon_{\rm out}}Q$ instead of $Q$, and dressed string products throughout. The present choice buys $\ell_1=Q$ and Zwiebach's products undressed, at the price of \eqref{ibk:eq:gen:symmetry} and of the interleaving sign from $\ell=3$ on. Nothing below depends on the choice: the two families differ by one uniform sign per closed contraction, so every relation of \cref{ibk:sec:compositions} is the same identity in either convention, with the names of the products exchanged by \eqref{ibk:eq:gen:dictionary}.

One remark to keep in view. Since the outputs of $p_{k,2}$ are antisymmetric in the physical parity by \eqref{ibk:eq:gen:symmetry}, while a two-string field is expected to live in a symmetric square, the two statements agree only on the components where the two output parities are equal. On physical inputs they always are; off the physical sector they are genuinely different subspaces. Every symmetric-power statement about the two-string field below is therefore a statement on physical inputs.

\begin{example}[$p_{2,2}$, the $\Sigma$ kinetic operator]
  \label{ibk:ex:p22}
  For $k=\ell=2$, \eqref{ibk:eq:gen:pkl} is the defining relation of the kinetic operator of the two-string field,
  \begin{equation}
    \label{ibk:eq:gen:p22}
    \mc V_4(A_1,A_2,A_3,A_4)
      = \braket{A_1\otimes A_2,\ p_{2,2}(A_3,A_4)}_2,
    \qquad
    N_{\rm gh}(p_{2,2})
      = +2,
  \end{equation}
  with no sign: the quartic vertex, with two of its legs held open, \emph{is} the object that pairs with a two-string state. In components,
  \begin{equation}
    \label{ibk:eq:gen:p22-comp}
    p_{2,2}(A_3,A_4)^{c_1c_2}
      = (-1)^{\varepsilon_{c_1}\bar\varepsilon_{c_2}}\,
        \Delta^{c_1b_1}\Delta^{c_2b_2}\,
        \mc V_4(e_{b_1},e_{b_2},A_3,A_4),
  \end{equation}
  and on physical inputs the prefactor is $+1$. There the two outputs sit at ghost number three each, so both are odd, and $\varepsilon$-antisymmetry in \eqref{ibk:eq:gen:symmetry} becomes symmetry under exchange: $p_{2,2}(\Psi,\Psi)$ is symmetric in its two output strands. This is the statement that $\Sigma$ lives in $\mathrm{Sym}^2\mc H_0$ and not in the antisymmetric part, now read off from the parities.
\end{example}

\subsubsection{The placeholder \texorpdfstring{$\chi$}{chi} and the generating function}
\label{ibk:sec:gen:genfun}

Rather than carrying the output vectors around, one may contract them with a second, even superfield $\chi=\chi^ae_a$ and record the result as a single number. This is exactly what \eqref{ibk:eq:gen:pkl} does when all the test states are taken equal to $\chi$,
\begin{equation}
  \label{ibk:eq:gen:one-output}
  \braket{\chi^{\otimes\ell},\ p_{k,\ell}(\Psi^k)}_\ell
    = \mc V_{k+\ell}(\chi^\ell,\Psi^k)
    = \mc V_{k+\ell}(\Psi^k,\chi^\ell),
\end{equation}
the last step because $\mc V$ is graded symmetric and $\Psi,\chi$ are even. So the number $\mc V_{k+\ell}(\Psi^k,\chi^\ell)$ is, with no sign at all, the matrix element of $p_{k,\ell}$ between the inputs $\Psi^k$ and the test state $\chi^{\otimes\ell}$: the power of $\Psi$ counts inputs, the power of $\chi$ counts outputs. Physically, $\chi$ is not a second dynamical field: it is a marker for an open leg, and nothing below depends on giving it a dynamical meaning. Collecting all the components with the symmetry factors,
\begin{equation}
  \label{ibk:eq:gen:H}
  H(\Psi,\chi)
    := \sum_{k,\ell\ge0,\ k+\ell\ge2}\frac{1}{k!\,\ell!}\,\mc V_{k+\ell}(\Psi^k,\chi^\ell)
    = S(\Psi+\chi),
\end{equation}
by the binomial expansion of $\mc V_n((\Psi+\chi)^n)$. So the generating function of the raised products is nothing but the action evaluated at a shifted field, and $p_{k,\ell}$ is the $k$-th derivative of the $\ell$-th-order vertex: the $k$-point function in the background $\chi$. This is the reason the whole family will turn out to contain no more information than $S$ itself does. We write $H_{k,\ell}$ for the bidegree-$(k,\ell)$ term of \eqref{ibk:eq:gen:H}.

This is the Weyl-algebra picture of Cieliebak, Fukaya and Latschev~\cite[\S7]{Cieliebak:2015:HomologicalAlgebraRelated}, with their variables $p\leftrightarrow\Psi$ for the inputs and $q\leftrightarrow\chi$ for the outputs, and with their formal weighting parameter odd, because our pairing is odd.

\subsection{Compositions and the IBL relations in the generating-function language}
\label{ibk:sec:compositions}

\subsubsection{One contraction between two vertices}
\label{ibk:sec:comp:one}

Feeding one output of $p_{k_1,\ell_1}$ into one input of $p_{k_2,\ell_2}$ is, in components, a single $\tilde\Delta$-contraction between two vertices,
\begin{equation}
  \label{ibk:eq:comp:contraction}
  \mc V_{k_2+\ell_2}(\dots,e_a,\dots)\,\tilde\Delta^{ab}\,\mc V_{k_1+\ell_1}(e_b,\dots),
\end{equation}
with the dressing \eqref{ibk:eq:gen:dressing} for moving the contracted legs into place. In the generating-function language the same operation is the one-sided bracket
\begin{equation}
  \label{ibk:eq:comp:bracket}
  \{F,G\}
    := \frac{\pd F}{\pd\Psi^a}\;\tilde\Delta^{ab}\;\frac{\pd G}{\pd\chi^b},
\end{equation}
which differentiates its first argument with respect to an input and its second with respect to an output. The indices here are component indices and nothing else: $\Psi=\Psi^ae_a$ and $\chi=\chi^ae_a$ are expanded on the basis, $\Psi^a$ and $\chi^a$ are the (Grassmann-even) coordinates, and $\pd F/\pd\Psi^a$ is the left derivative of the function $F$ with respect to the coordinate $\Psi^a$, a component of a covector that the bivector $\tilde\Delta^{ab}$ then contracts. The weight sitting in $\tilde\Delta$ is the one of \cref{ibk:sec:gen:pairing}: this contraction closes a bracket. Applied to two bidegree components,
\begin{equation}
  \label{ibk:eq:comp:bidegree}
  \{H_{k_1,\ell_1},H_{k_2,\ell_2}\}
    = \frac{1}{(k_1-1)!\,\ell_1!\,k_2!\,(\ell_2-1)!}\,
      \mc V(e_a,\Psi^{k_1-1},\chi^{\ell_1})\,\tilde\Delta^{ab}\,\mc V(e_b,\Psi^{k_2},\chi^{\ell_2-1}),
\end{equation}
which has bidegree $(k_1+k_2-1,\ell_1+\ell_2-1)$: one input and one output have been consumed. The factorials are the symmetry factors of \eqref{ibk:eq:gen:H} after one differentiation on each side. Differentiating $\frac{1}{k_1!\ell_1!}\mc V(\Psi^{k_1},\chi^{\ell_1})$ with respect to $\Psi^a$ removes one of the $k_1$ equal $\Psi$ slots, and since any one of them can be the one removed, the derivative produces a factor $k_1$ which cancels one factorial and leaves $1/(k_1-1)!\,\ell_1!$. The same happens on the right for the $\chi$-derivative, which removes one of the $\ell_2$ equal $\chi$ slots and leaves $1/k_2!\,(\ell_2-1)!$. So the factorials are not put in by hand: they are what survives of the $1/n!$ of the action once one leg on each vertex has been used up by the contraction. Cieliebak, Fukaya and Latschev's Proposition~7.2 says that the composition $\hat p_{k_2,\ell_2}\circ_1\hat p_{k_1,\ell_1}$ of the coderivation extensions, with spectators passing through, corresponds to the bracket \eqref{ibk:eq:comp:bidegree} with weight exactly one: the spectators are the untouched variables of the polynomial. This is the convenient feature of the generating function. The hat-extension in bare variables has no uniform sign rule on all parity sectors, because a spectator is an input (parity $\varepsilon$) which is also an output (parity $\bar\varepsilon$), whereas in \eqref{ibk:eq:comp:bracket} everything commutes correctly because $\Psi$ and $\chi$ are even superfields.

It is useful to name the individual configurations. Write $T(i,j\,|\,m,n)$ for the contraction \eqref{ibk:eq:comp:bidegree} in which the first vertex keeps $i$ free input legs and $j$ free output legs and the second keeps $m$ and $n$, with the symmetry factor $1/(i!j!m!n!)$. Relabelling $a\leftrightarrow b$ and using the graded symmetry of $\tilde\Delta$ shows that a configuration equals its mirror image,
\begin{equation}
  \label{ibk:eq:comp:mirror}
  T(m,n\,|\,i,j)
    = T(i,j\,|\,m,n).
\end{equation}
The content of \eqref{ibk:eq:comp:mirror} is that a diagram does not care which of its two vertices one calls the emitter and which the absorber: the same propagator joins them either way. With the plain $\Delta$ the same steps would instead give a minus sign, and every sum below would vanish identically; this is the sharp form of the remark after \eqref{ibk:eq:gen:antibracket}.

\subsubsection{The IBL relations and what augmentation removes}
\label{ibk:sec:comp:aug}

The genus-zero IBL relations in the sense of Cieliebak, Fukaya and Latschev's Definition~2.3 concern the family with $k,\ell\ge1$ only. In generating-function form they are the vanishing, bidegree by bidegree, of the bracket of
\begin{equation}
  \label{ibk:eq:comp:H0}
  H_0(\Psi,\chi)
    := S(\Psi+\chi)-S(\Psi)-S(\chi)
    = \sum_{k,\ell\ge1}H_{k,\ell}
\end{equation}
with itself. The subtraction removes the all-input components $H_{k,0}=\frac1{k!}\mc V_k(\Psi^k)$ and the all-output components $H_{0,\ell}$. In the IBL vocabulary this is the augmentation: the Hamiltonian is required to vanish when either set of variables is set to zero. Here $S(\Psi)$ and $S(\chi)$ are the same nonzero function, so the full generating function \eqref{ibk:eq:gen:H} is not augmented, and $H_0$ is what one gets by forcing the condition.

Now compare with the master equation. By the chain rule for $u=\Psi+\chi$, the bracket of the full generating function with itself is the antibracket \eqref{ibk:eq:gen:antibracket} evaluated at the shifted field,
\begin{equation}
  \label{ibk:eq:comp:master}
  \{S(u),S(u)\}
    = (\pd_aS)(u)\,\tilde\Delta^{ab}\,(\pd_bS)(u)
    = (S,S)\big|_{\Psi\to u},
\end{equation}
and its bidegree components $B_{k,\ell}$ vanish for all $k,\ell$ exactly when the algebra is cyclic $L_\infty$. Sorting the configurations in $B_{k,\ell}$ and in $R_{k,\ell}:=\{H_0,H_0\}|_{(k,\ell)}$,
\begin{equation}
  \label{ibk:eq:comp:sums}
  B_{k,\ell}
    = \sum_{\substack{i+m=k,\ j+n=\ell\\ (i,j),(m,n)\neq(0,0)}}T(i,j\,|\,m,n),
  \qquad
  R_{k,\ell}
    = \sum_{\substack{i+m=k,\ j+n=\ell\\ j\ge1,\ m\ge1}}T(i,j\,|\,m,n).
\end{equation}
The two sums run over the same diagrams but with different admissibility rules: $B_{k,\ell}$ keeps every way of splitting the free legs between the two vertices, while $R_{k,\ell}$ keeps only those in which the leg contracted on the left comes from an output slot and the one on the right from an input slot, both vertices being asked to have at least one leg of the relevant kind. Call a configuration \emph{mixed} when all four numbers $i,j,m,n$ are at least one, so that each vertex keeps both an input and an output, and \emph{pure} otherwise: in a pure configuration one vertex is the all-output component $p_{0,j}$ feeding one leg into an input of the other, or the all-input functional $p_{m,0}$ eating one output of the other. Let $M_{k,\ell}$ be the sum of the mixed configurations and $P_{k,\ell}:=R_{k,\ell}-M_{k,\ell}$ the pure ones that survive in $R$. Grouping the terms of $B$ into mirror pairs and using \eqref{ibk:eq:comp:mirror}, a pure pair has exactly one member in $R$ and a mixed pair has both, so
\begin{equation}
  \label{ibk:eq:comp:decomposition}
  B_{k,\ell}
    = 2P_{k,\ell}+M_{k,\ell},
  \qquad
  R_{k,\ell}
    = \tfrac12 B_{k,\ell}+\tfrac12 M_{k,\ell}
  \qquad (k,\ell\ge1).
\end{equation}

This is the key observation, in one line. The master equation contains, besides the output-to-input contractions between two components with $k,\ell\ge1$, the contractions involving an all-input or an all-output component. The augmented relations drop the latter. But the master equation constrains only the single combination $2P+M$; it does not make $P$ and $M$ vanish separately. Given a cyclic $L_\infty$ algebra, therefore,
\begin{equation}
  \label{ibk:eq:comp:half-M}
  R_{k,\ell}
    = \tfrac12 M_{k,\ell}
    = -P_{k,\ell},
\end{equation}
and the augmented relation at $(k,\ell)$ holds if and only if there is no mixed configuration or the mixed ones happen to cancel. A mixed configuration needs $i,m\ge1$, hence $k\ge2$, and $j,n\ge1$, hence $\ell\ge2$. So for $\min(k,\ell)=1$ there is none, $M_{k,\ell}=0$, and the augmented relations hold: at $(k,1)$ they are the $L_\infty$ relations themselves, and at $(1,\ell)$ they express the BRST compatibility of the tower of multi-string states emitted by a single vertex, written out in \cref{ibk:sec:comp:1ell}. For $k,\ell\ge2$ the augmented relation is the \emph{additional} quadratic condition $P_{k,\ell}=0$ on the products, which the $L_\infty$ relations do not imply; what it says at the first such signature, with the string fields put back in, is the subject of \cref{ibk:sec:comp:22}.

\subsubsection{The \texorpdfstring{$(1,\ell)$}{(1,l)} relations: the emitted multi-string states are BRST compatible}
\label{ibk:sec:comp:1ell}

The signature $(1,\ell)$ is the one with a single input, and the object it constrains has a direct physical reading. Feed the background field $\Psi$ into one leg of the vertex $\mc V_{\ell+1}$ and cut the remaining $\ell$ legs open: what comes out is $p_{1,\ell}(\Psi)$, the $\ell$-string state emitted by that vertex when one of its legs absorbs $\Psi$. It is a map $\mc H\to\mc H^{\otimes\ell}$ of ghost number $3\ell-2$, and the family $\{p_{1,\ell}\}_{\ell\ge1}$, which starts at $p_{1,1}=Q$, is a tower of such emissions of growing multiplicity. The $(1,\ell)$ relations say that this tower is compatible with the BRST charge: acting with $Q$ on an emitted state and emitting from the BRST variation of the background differ only by emissions that factor through a shorter one.

Written out, at $\ell=2$ this is a plain commutator,
\begin{equation}
  \label{ibk:eq:comp:1ell-2}
  p_{1,2}(Q\Psi)+\hat Q\,p_{1,2}(\Psi)
    = 0,
  \qquad\text{i.e.}\qquad
  [\hat Q,p_{1,2}]
    = 0,
\end{equation}
where $\hat Q$ is the extension of $Q$ to two strands, acting on each in turn with the Koszul sign of the strand it passes, as in \eqref{ibk:eq:gen:dressing}. So the two-string state emitted by the cubic vertex is BRST covariant on the nose. At $\ell=3$ the covariance is spoiled, and by a definite amount,
\begin{equation}
  \label{ibk:eq:comp:1ell-3}
  p_{1,3}(Q\Psi)
    + \bigl(\hat p_{1,2}\otimes1\bigr)\bigl(p_{1,2}(\Psi)\bigr)_{\rm sym}
    + \hat Q\,p_{1,3}(\Psi)
    = 0,
\end{equation}
where the middle term is $p_{1,2}$ acting on one of the two strands already emitted, the other strand looking on, symmetrised over the three resulting output strands. Reading \eqref{ibk:eq:comp:1ell-3} as a statement about $p_{1,3}$: the three-string emission of the quartic vertex fails to commute with $Q$ by exactly the two-step process in which the cubic vertex emits twice. This is the same pattern as the $L_\infty$ relations themselves, where the failure of $Q$ to be a derivation of $\ell_3$ is the Jacobiator of $\ell_2$. What \eqref{ibk:eq:comp:1ell-2}--\eqref{ibk:eq:comp:1ell-3} are, precisely, is the single-input part of the genus-zero relations in coderivation form: the one-leg compositions $\hat p_{k_2,\ell_2}\circ_1\hat p_{k_1,\ell_1}$ of \cref{ibk:sec:compositions}, summed at $k=1$. It is tempting to compress this into \enquote{the family $Q+\sum_{\ell\ge2}p_{1,\ell}$ squares to zero}, and one should resist: what the relations say is that the single-input components of the coderivation-form relations vanish, which constrains those compositions and is not the statement that an operator squares to zero. From two outputs on, no differential operator on a single graded-commutative polynomial host in the strand variables, with the same parity-homogeneous raising map in every slot of a given type and an even grading parameter, realises the family with its propagator contraction; other hosts --- the symbol calculus, separate generators for inputs and outputs, an odd parameter --- are untested, not excluded. See \cref{ibk:sec:cubic}.

At $k=1$ the augmented and the non-augmented relation say the same thing. Allowing the components with no input or no output adds, at each $\ell$, exactly as many terms as there already were, and each new term equals its partner by the mirror symmetry \eqref{ibk:eq:comp:mirror}; the sum is doubled and the relation is unchanged. This is the structural difference from $(2,2)$: there the extra terms are again mirrors, but only of the pure terms; the mixed term is its own mirror and is not doubled, so the relation changes weight.

Of this tower, the cubic action at order $g_s^2$ uses only \eqref{ibk:eq:comp:1ell-2}. The quartic vertex enters the action at that order through $\mc V_4(\Psi,\Psi,\cdot,\cdot)=p_{2,2}$ contracted with the auxiliary two-string field, and the gauge variation with the composite law $\delta\Sigma=\delta\Psi\otimes\Psi+\Psi\otimes\delta\Psi$ closes because of the two relations at $(2,1)$ and $(1,2)$. The first is $[Q,\ell_2]=0$; the second is \eqref{ibk:eq:comp:1ell-2}, which is equivalent to it, since $p_{1,2}$ is the raised $\ell_2$. So the only structural input at this order is that $Q$ is a derivation of the cubic product, used once on the one-output side and once on the two-output side. The relation \eqref{ibk:eq:comp:1ell-3}, and with it $p_{1,3}$, first matters at order $g_s^3$.

\subsubsection{The \texorpdfstring{$(2,2)$}{(2,2)} relation with the fields inserted}
\label{ibk:sec:comp:22}

At $(2,2)$ one would like to read the relation the way one reads $[Q,\ell_2]=0$ (\enquote{$Q$ is a derivation of the string product}) or $[Q,\ell_3]+\ell_2\circ_1\ell_2=0$ (\enquote{the failure of $Q$ to be a derivation of $\ell_3$ is the Jacobiator of $\ell_2$}). It does not have that shape, and seeing why is the point of this subsection.

Sorting the bidegree-$(2,2)$ part of $\{S(u),S(u)\}=0$ by configuration gives seven ordered terms. Each is one contraction between two vertices, and each has a name:
\begin{enumerate}
  \setlength{\itemsep}{0.6ex}

  \item $\tfrac12\,\hat Q\,p_{2,2}(\Psi,\Psi)$: the quartic vertex with two legs on the background and two open, one of the open strands then hitting the BRST charge. The configuration carries $Q$ on \emph{one} strand, while $\hat Q$ acts on both; the $\tfrac12$ is that difference.

  \item $\tilde p_{2,2}(Q\Psi,\Psi)$: the same vertex, with one of the two background legs replaced by $Q\Psi$.

  \item $\tfrac12\,\tilde p_{1,2}\bigl(\tilde\ell_2(\Psi,\Psi)\bigr)$: the two backgrounds first fuse through the cubic product, and the result is fed into the single input of the two-string emission.

  \item $2\,\bigl[\bigl(\ell_2\otimes1\bigr)\bigl(p_{1,2}(\Psi)\otimes\Psi\bigr)\bigr]_{\rm sym}$: the cubic vertex emits a two-string state, and one of its two strands is glued, through the cubic product, to the remaining background leg; either of the two emitted strands can be the glued one, and the factor two counts the two ways. Both vertices keep an input and an output --- this is the \emph{mixed} configuration of \cref{ibk:sec:comp:aug}, and at this signature it is the only one.

  \item $\bigl[\bigl(\ell_3(\Psi,\Psi,\cdot)\otimes1\bigr)\tilde p_{0,2}\bigr]_{\rm sym}$: the no-input component $\tilde p_{0,2}$, a two-string state produced by the propagator out of the quadratic vertex, with one of its strands absorbed by $\ell_3$ together with the two backgrounds.

  \item $\bigl(\mc V_2(\Psi,\cdot)\otimes1^{\otimes2}\bigr)\tilde p_{1,3}(\Psi)$: the three-string emission of the quartic vertex, one of whose strands is closed off against the background through the no-output component $p_{2,0}=\mc V_2$.

  \item $\tfrac12\bigl(\mc V_3(\Psi,\Psi,\cdot)\otimes1^{\otimes2}\bigr)\tilde p_{0,3}$: likewise with the three-string no-input component $\tilde p_{0,3}$ closed off against the cubic vertex carrying the two backgrounds.
\end{enumerate}
Each term is a map $\mathrm{Sym}^2\mc H\to\mc H^{\otimes2}$ evaluated at $(\Psi,\Psi)$ and symmetrised over its two output strands --- on physical inputs, ordinary symmetrisation, by \cref{ibk:ex:p22}. Three conventions have to be fixed before the seven can be added up, and all three are visible in the coefficients just written; the relative signs are those of the display below.

First, $\hat Q$ is the coderivation of $Q$ on the two output strands, $Q\otimes1+1\otimes Q$, the strand it passes contributing the Koszul sign of its \emph{shifted} parity $\bar\varepsilon$ as in \eqref{ibk:eq:gen:dressing}; so read, $\hat Q$ preserves the symmetry type of the outputs. Second, $[\,\cdot\,]_{\rm sym}$ is the graded symmetrisation over the two free output strands in $\bar\varepsilon$, normalised so that a map already symmetric is left unchanged, and $\otimes1$ means that the map on its left is fed strand one of the tensor on its right, the other strand looking on. Third, the products are the raised ones $\tilde p_{k,\ell}$, $\tilde\ell_n$ of \eqref{ibk:eq:gen:dictionary}, with two exceptions. In the fourth and fifth terms, where the output of a one-output product is itself a free strand of the result, the sign-free $\ell_2$ and $\ell_3$ of \eqref{ibk:eq:gen:pkl} appear, and they have to stay sign-free: the strand they act on is an open one, not a physical input, so the relation $\tilde\ell_n=-\ell_n$ of \cref{ibk:sec:gen} does not apply to them. And where a two-output product is evaluated on physical arguments the two conventions coincide and the tilde is dropped, which is why the two objects the cubic action uses, $p_{2,2}(\Psi,\Psi)$ and $p_{1,2}(\Psi)$, appear undressed. The third term is not an exception: there $\tilde\ell_2$ acts on the two backgrounds, so it may equally be written $-\ell_2(\Psi,\Psi)$, and the reader who prefers the sign-free product throughout should flip the sign of that term.

The coefficients are then a leg count: a term written on one output strand carries $\tfrac12\binom2i\binom2j$, where $i$ and $j$ are the numbers of backgrounds and of open strands sitting on one of the two vertices of the configuration $T(i,j\,|\,m,n)$ --- either one, since $\binom2i=\binom2m$. The two that are easy to get wrong are the first and the fourth. On the first, the leg count gives $\tfrac12\binom22\binom21=1$ per strand, and $\hat Q$ has two strands, so the coefficient in front of $\hat Q$ is $\tfrac12$, as said in the list. On the fourth it gives two, the two strands of $p_{1,2}(\Psi)$ that $\ell_2$ can absorb. With these coefficients each term written above \emph{is} the corresponding configuration, and the relation is that the seven add to zero,
\begin{equation}
  \label{ibk:eq:comp:22-map}
  \begin{split}
    &-\tfrac12\,\hat Q\,p_{2,2}(\Psi,\Psi)
      + \tilde p_{2,2}(Q\Psi,\Psi)
      - \tfrac12\,\tilde p_{1,2}\bigl(\tilde\ell_2(\Psi,\Psi)\bigr)
      - 2\,\bigl[\bigl(\ell_2\otimes1\bigr)\bigl(p_{1,2}(\Psi)\otimes\Psi\bigr)\bigr]_{\rm sym}
    \\
    &\quad
      - \bigl[\bigl(\ell_3(\Psi,\Psi,\cdot)\otimes1\bigr)\tilde p_{0,2}\bigr]_{\rm sym}
      + \bigl(\mc V_2(\Psi,\cdot)\otimes1^{\otimes2}\bigr)\tilde p_{1,3}(\Psi)
      + \tfrac12\bigl(\mc V_3(\Psi,\Psi,\cdot)\otimes1^{\otimes2}\bigr)\tilde p_{0,3}
      = 0.
  \end{split}
\end{equation}
This is the identity on physical backgrounds, $N_{\rm gh}(\Psi)=2$. Every coefficient and every relative sign in it has been measured row by row against the corresponding configuration, on all parity sectors and on two random models of the vertices; off the physical sector the same seven terms close, the rows other than the first carrying one Koszul sign each that is not written here. The physical-sector form above was also checked directly, as an identity assembled with the sign-free products and no sign tensors, rather than only read off the all-sector one.
Where those dressings sit inside the named compositions --- the piece of bookkeeping an earlier version of this note left open --- is exactly what the two sign-free $\ell$'s and the two dropped tildes record: in the two rows where the output of a one-output product is itself a free strand of the result, the product that appears is the sign-free one.

The last three terms are the mirror images of the first three in the sense of \eqref{ibk:eq:comp:mirror} --- exchanging which of the two vertices is differentiated in $\Psi$ and which in $\chi$ turns a $p_{k,\ell}$ into a $p_{k-1,\ell+1}$ --- so they are not new physical input: as test-field scalars the last three equal the first three, and the seven-term sum collapses to twice the first three plus the mixed one.

Now the reading. The first two terms of \eqref{ibk:eq:comp:22-map} are the failure of $Q$ to be a derivation of the two-string kinetic operator $p_{2,2}$. They are exactly that, in the form of a graded commutator: since $p_{2,2}$ and $\Psi$ are both even, letting the coderivation act on the two \emph{inputs} gives twice $\tilde p_{2,2}(Q\Psi,\Psi)$, so the two terms together are $-\tfrac12[\hat Q,p_{2,2}](\Psi,\Psi)$ with $[\hat Q,p_{2,2}]=\hat Q\,p_{2,2}-p_{2,2}\hat Q$. Add the third term, the chain in which the quartic vertex is replaced by two cubic ones, which is the exact analogue of the Jacobiator term at $(3,1)$. Were the pattern of $(2,1)$ and $(3,1)$ to persist, that combination would vanish. It does not; multiplying \eqref{ibk:eq:comp:22-map} by $-2$ and keeping the first three terms on the left,
\begin{equation}
  \label{ibk:eq:comp:22-reading}
  \begin{split}
    &\underbrace{[\hat Q,p_{2,2}](\Psi,\Psi)}_{\text{failure of }Q\text{ to be a derivation of }p_{2,2}}
      + \underbrace{\tilde p_{1,2}\bigl(\tilde\ell_2(\Psi,\Psi)\bigr)}_{\text{the analogue of the Jacobiator}}
    \\
    &\qquad
      = -4\,\bigl[\bigl(\ell_2\otimes1\bigr)\bigl(p_{1,2}(\Psi)\otimes\Psi\bigr)\bigr]_{\rm sym}
        - 2\,\bigl[\bigl(\ell_3(\Psi,\Psi,\cdot)\otimes1\bigr)\tilde p_{0,2}\bigr]_{\rm sym}
    \\
    &\qquad\quad
      + 2\,\bigl(\mc V_2(\Psi,\cdot)\otimes1^{\otimes2}\bigr)\tilde p_{1,3}(\Psi)
      + \bigl(\mc V_3(\Psi,\Psi,\cdot)\otimes1^{\otimes2}\bigr)\tilde p_{0,3}.
  \end{split}
\end{equation}
The first term on the right is the two-string state obtained by gluing two cubic vertices along a propagator with one leg of each on the background. In the language of \cref{ibk:sec:gen} it is the square of the cubic vertex in test-field form: contracting the two open strands with $\chi$ turns it into $\tilde\omega\bigl(\ell_2(\Psi,\chi),\ell_2(\Psi,\chi)\bigr)$, a vertex paired with itself. That object is not zero --- it would be, were the weight of \cref{ibk:sec:gen:pairing} absent, and this is precisely where the weight earns its place. The other three terms on the right are the mirror ones, built from the components with no output, $p_{2,0}=\mc V_2$, and with no input, $\tilde p_{0,2}$ and $\tilde p_{0,3}$.

So at $(2,2)$ the pattern breaks: what would be the analogue of the Jacobi identity does not close on its own. What restores it can be said in two ways, and they are one statement seen before and after the mirror symmetry is used. Before: the first four terms of \eqref{ibk:eq:comp:22-map} do not add to zero --- they add to half the mixed configuration --- and the sum closes only once the last three, the ones built from $p_{2,0}=\mc V_2$, $\tilde p_{0,2}$ and $\tilde p_{0,3}$, are put back. After: since each of those three equals one of the first three, putting them back is the same as counting the first three twice against the mixed term. Keeping or dropping them is the whole difference between the two readings of \cref{ibk:sec:three-senses}, and it is worth stating the two side by side.

\begin{description}
  \setlength{\itemsep}{0.8ex}

  \item[Augmented --- the first four terms alone.] \emph{False} in general. Dropping the three mirror terms leaves a combination that does not vanish; the defect is half the mixed configuration, and on the random algebra of \cref{ibk:sec:comp:obstruction} it is of order ten, the size of the individual terms.

  \item[Non-augmented --- all seven terms.] \emph{True}, on every parity sector. This is the bidegree-$(2,2)$ component of the polarised master equation, and it holds whenever the $\ell_n$ form a cyclic $L_\infty$ algebra.
\end{description}
The two differ only in the weight given to the mixed term: the true relation counts it once against twice the sum of the other three, the augmented one counts it once against that sum. No choice of signs in the bare variables changes this, the raising being forced as in \eqref{ibk:eq:comp:M22}: the sign dressings of \cref{ibk:sec:gen} move the raised legs between slots and never alter the relative weight of a configuration, so no convention makes the mixed term enter the augmented combination with the weight $\tfrac12$ that would turn it into the true relation.

\subsubsection{Why the augmented relation fails: the obstruction at \texorpdfstring{$(2,2)$}{(2,2)}}
\label{ibk:sec:comp:obstruction}

The quickest way to see that the mixed term cannot be made to vanish is to go back to the test-field form. Every configuration is a weighted pairing of two products with superfields in the remaining slots, $\mc V(e_a,\dots)\tilde\Delta^{ab}\mc V(e_b,\dots)=\tilde\omega(\ell(\dots),\ell(\dots))$ with the dressed pairing $\tilde\omega$ of \cref{ibk:sec:gen:pairing}. At $(2,2)$ there is a single mixed configuration, the self-contraction of $X_a:=\mc V_3(e_a,\Psi,\chi)$ with the bivector,
\begin{equation}
  \label{ibk:eq:comp:M22}
  M_{2,2}
    = X_a\,\tilde\Delta^{ab}\,X_b
    = \tilde\omega\bigl(\ell_2(\Psi,\chi),\ell_2(\Psi,\chi)\bigr),
\end{equation}
which does not vanish by antisymmetry precisely because $\tilde\Delta$ is graded symmetric; with the plain $\Delta$ it would be identically zero, which is the version that had been computed on the even sector in the earlier check. The pure part is
\begin{equation}
  \label{ibk:eq:comp:P22}
  P_{2,2}
    = \tfrac12\tilde\omega\bigl(Q\chi,\ell_3(\Psi,\Psi,\chi)\bigr)
      + \tfrac12\tilde\omega\bigl(\ell_3(\Psi,\chi,\chi),Q\Psi\bigr)
      + \tfrac14\tilde\omega\bigl(\ell_2(\chi,\chi),\ell_2(\Psi,\Psi)\bigr),
\end{equation}
and the $L_\infty$ relation at this bidegree is $2P_{2,2}+M_{2,2}=0$. The three terms of $P_{2,2}$ are the pure configurations: in the first two one vertex is the quadratic vertex, i.e.\ $Q$ acting on a strand that is not shared, and in the third the two cubic vertices split the background legs and the test legs between them, one keeping only inputs and the other only outputs. None of them has both a free input and a free output on each vertex, which is what distinguishes them from $M_{2,2}$.

For a cyclic dg Lie algebra, $\ell_{n\ge3}=0$, the $Q$-terms drop out and $M_{2,2}=-\frac12\tilde\omega(\ell_2(\chi,\chi),\ell_2(\Psi,\Psi))$. Reading off the coefficient of $\chi^\mu\chi^\nu\Psi^\rho\Psi^\sigma$ and moving one $\ell_2$ across the pairing by cyclicity,
\begin{equation}
  \label{ibk:eq:comp:dg-lie}
  M_{2,2}\equiv0
  \iff
  \braket{\ell_2(\ell_2(e_\mu,e_\nu),e_\rho),e_\sigma}=0\ \ \forall\mu,\nu,\rho,\sigma
  \iff
  \ell_2\circ_1\ell_2=0.
\end{equation}
The augmented $(2,2)$ relation holds for a cyclic dg Lie algebra exactly when the bracket is two-step nilpotent. The smallest exactly solvable instance is the Chevalley--Eilenberg algebra of a Lie algebra $\alg g$ in BV form, $\mc H=\alg g_{(1)}\oplus\alg g^*_{(4)}$ with $\braket{c,c^*}=\langle c^*,c\rangle$, $S=\frac12\langle c^*,[c,c]\rangle$ and $Q=0$: there $M_{2,2}\neq0$ if and only if $[\alg g,[\alg g,\alg g]]\neq0$, so $\alg{sl}_2$ violates the relation and the Heisenberg algebra satisfies it. The obstruction needs at least three ghost-number levels; on a space whose ghost numbers occupy only $\{2,3\}$, $\ell_2\circ_1\ell_2$ vanishes by degree, which is why an earlier random search over such toys found only trivially satisfied relations.

On a randomly generated cyclic $L_\infty$ algebra with four ghost-number levels the same failure was seen numerically: at $(2,2)$, $(2,3)$ and $(3,2)$ the augmented residual $R_{k,\ell}$ equals half the mixed sum to machine precision and is of order $10$--$10^2$, and at $(2,2)$ it stays of order $10$ when the inputs are restricted to ghost number $2$, so the failure is not an artefact of unphysical sectors. Whether $P_{k,\ell}\neq0$ for \emph{every} $k,\ell\ge2$ on generic data was not established; what is established is that the condition is independent of the $L_\infty$ relations and fails at the three signatures tested.

\subsection{Three senses of \enquote{\texorpdfstring{$L_\infty$}{L-infinity} implies IBL}}
\label{ibk:sec:three-senses}

The slogan \enquote{a cyclic $L_\infty$ algebra implies an IBL structure on its raised products} can mean three things, and they have different truth values. The one that is true is the one to state first. Take the family of \emph{all} raised products, the components with no input and the components with no output included. That family satisfies the genus-zero relations precisely when the string products form a cyclic $L_\infty$ algebra; this is \cref{ibk:thm:non-augmented} below. It is the expected statement, and the reason is one line long: those relations are the polarised master equation $(S,S)(\Psi+\chi)=0$ sorted by the number of free legs, so they say neither more nor less than $(S,S)=0$ does. The two other readings ask for something stronger --- the relations of the IBL definition proper, which requires discarding the all-input and all-output components, and a small closed sub-family containing $p_{2,2}$ --- and both fail. The discussion in \cref{ibk:sec:compositions} decides each of the three.

\subsubsection{All raised products, augmented: false for \texorpdfstring{$k,\ell\ge2$}{k,l >= 2}}
\label{ibk:sec:senses:aug}

This is the reading in which the family $\{p_{k,\ell}\}_{k,\ell\ge1}$ is asked to satisfy the genus-zero relations of Definition~2.3 of Cieliebak, Fukaya and Latschev, i.e.\ $\{H_0,H_0\}=0$.

\begin{theorem}[Augmented relations]
  \label{ibk:thm:augmented}
  Let $\{\ell_n\}$ be a cyclic $L_\infty$ algebra on $(\mc H,\braket{\cdot,\cdot})$. The augmented family $\{p_{k,\ell}\}_{k,\ell\ge1}$ satisfies the genus-zero relation at every signature with $\min(k,\ell)=1$. At $k,\ell\ge2$ the relation is equivalent to the extra condition $M_{k,\ell}=0$ on the mixed configurations, which the $L_\infty$ relations do not imply; it fails at $(2,2)$ already for the Chevalley--Eilenberg algebra of $\alg{sl}_2$.
\end{theorem}

Both halves are read off \eqref{ibk:eq:comp:half-M}: the augmented residual is $R_{k,\ell}=\frac12M_{k,\ell}$, which vanishes for lack of mixed configurations when $\min(k,\ell)=1$ and is an independent condition otherwise. Physically, the augmented reading asks the products with at least one input and at least one output to close among themselves; what \cref{ibk:thm:augmented} says is that they do not, because a diagram in which each of the two vertices keeps a leg on both sides has no partner to cancel against once the pure components are gone. Since the two augmentation conditions fail by exactly the pure-input functionals and the pure-output elements, and expanding around a solution of the equations of motion changes $S$ but leaves both $S(\Psi)$ and $S(\chi)$ nonzero, no such twist restores augmentation; a genuine augmentation would need $S(\chi)=0$ identically, which we have not pursued.

\subsubsection{Non-augmented, M\"unster--Sachs sense: true, and nothing beyond \texorpdfstring{$L_\infty$}{L-infinity}}
\label{ibk:sec:senses:non-aug}

Keep every component, including the two-input zero-output pairing $p_{2,0}=\mc V_2$ and more generally $p_{k,0}=\mc V_k$ and the no-input elements $p_{0,\ell}\in\mc H^{\otimes\ell}$. Then the statement is exact and goes both ways.

\begin{theorem}[Non-augmented relations]
  \label{ibk:thm:non-augmented}
  Let $\{\mc V_n\}_{n\ge2}$ be graded-symmetric forms on $(\mc H,\braket{\cdot,\cdot})$ and let $p_{k,\ell}$ be the raised products \eqref{ibk:eq:gen:pkl}. The full family $\{p_{k,\ell}\}_{k,\ell\ge0,\,k+\ell\ge2}$ satisfies the genus-zero relations
  \begin{equation}
    \label{ibk:eq:senses:non-augmented}
    B_{k,\ell}
      = \sum_{\substack{k_1+k_2=k+1\\ \ell_1+\ell_2=\ell+1\\ k_1,\ell_2\ge1;\ \ell_1,k_2\ge0}}
        \hat p_{k_2,\ell_2}\circ_1\hat p_{k_1,\ell_1}\Big|_{\mathrm{Sym}^k}
      = 0
    \quad\text{for all }(k,\ell)
  \end{equation}
  if and only if $\{\ell_n=p_{n,1}\}$ is a cyclic $L_\infty$ algebra with cyclic forms $\mc V_n$.
\end{theorem}

The proof is \eqref{ibk:eq:comp:master}: the relations are the bidegree components of the polarised master equation $(S,S)(\Psi+\chi)=0$, and nothing else is used. Each $B_{k,\ell}$ collects the ways of gluing two vertices along one propagator so that $k$ input legs and $\ell$ output legs are left free, and asking them to cancel among themselves, signature by signature, is the same as asking the equations of motion to have vanishing weighted square at the shifted field $\Psi+\chi$. The sum in \eqref{ibk:eq:senses:non-augmented} is the genus-zero relation of Cieliebak, Fukaya and Latschev extended to $\ell_1=0$ (a zero-output factor eating one output of the other product) and to $k_2=0$ (a zero-input factor feeding one leg into the other). Structures with constants of this kind are the ones they admit in their \S8. Two caveats keep the identification from being literal: the degree bookkeeping for an odd pairing requires their $d=\frac12$, outside their integer-$d$ hypotheses, and the slot symmetries \eqref{ibk:eq:gen:symmetry} need the per-slot dictionary rather than one grading on a single space. So the family satisfies the genus-zero IBL relations in generating-function form, of the non-augmented type, and is not an IBL$_\infty$ algebra in the Definition~2.3 sense.

Two comments on what this does and does not give. First, the master equation organised leg by leg gains no structure beyond $L_\infty$: by \eqref{ibk:eq:gen:H} the family is the Taylor expansion of $S$ about a second field, and the relations are $(S,S)=0$ at the shifted point. Second, the bracket \eqref{ibk:eq:comp:bracket} is a single contraction between two distinct vertices, and this is what genus zero means here; the components with $k=0$ or $\ell=0$ enter through exactly such single contractions and do not introduce loops. Terms with two contractions between the two vertices belong to the next order of the loop expansion and are one class of the one-loop terms; the other is the self-contraction of a single vertex (the BV Laplacian), also outside the present statement.

The algebraic structure that governs a family of maps with several inputs and several outputs of this kind is the one studied by Markl and Voronov~\cite{Markl:2015:MVFormalism} and, in the BV setting, by Bashkirov and Voronov~\cite{Bashkirov:2017:BVFormalismLinftyAlgebras}. The family of \cref{ibk:thm:non-augmented} is the genus-zero part of such a structure, but only up to the parity caveat already stated: the generator assembled from the $p_{k,\ell}$ is not parity-homogeneous here, so the odd square-zero generator those formalisms assume is not available, and we do not claim the identification beyond the level of form.

The loop algebra of M\"unster and Sachs~\cite{Muenster:2014:HomotopyClassificationBosonic} is the other genus-zero packaging of the same cyclic $L_\infty$ data: at genus zero it consists of the single-output coderivation $L^0=\widehat{\sum_n\ell_n}$, which in our variables is the slice of $H$ linear in $\Psi$, together with a second-order term $\Omega_c^{-1}$ with zero inputs and two outputs, which in our variables is $\frac12\tilde\omega^{-1}(\chi,\chi)$ and sits at the first order of their loop expansion. Their family is augmented in the output direction, since every term has at least one output, but not in the input direction, and it has neither zero-output components nor multi-input products $p_{k\ge2,\ell}$. Their $\mathfrak L_c^2=0$ at genus zero is the $L_\infty$ relation together with the invariance of the bivector. The two packagings share the $L^0$ slice.

\subsubsection{Sub-algebra sense: what a sub-family containing \texorpdfstring{$p_{2,2}$}{p22} must carry}
\label{ibk:sec:senses:sub}

The reading intended for the cubic action is different again. Witten's cubic open string field theory is an $A_2\subset A_\infty$ structure, and the quantum $L_\infty$ algebra with the BV Laplacian sits inside IBL$_\infty$: a chosen subset of components is nonzero, the others are set to zero, and \emph{every} relation must still hold for the reduced family. The question is which sub-family of raised products on $\mc H$ can contain $p_{2,2}$ in this sense.

The count goes as follows. A candidate $\{Q,\ell_2,\ell_3,p_{1,2},p_{2,2}\}$ is not closed: the relation at $(1,3)$,
\begin{equation}
  \label{ibk:eq:senses:13}
  [\hat Q,\hat p_{1,3}]+\hat p_{1,2}\circ_1\hat p_{1,2}
    = 0,
\end{equation}
is of the same order $g_s^2$ as $(2,2)$ and its second term is nonzero (it equals $-[\hat Q,\hat p_{1,3}]\neq0$ by the $\min(k,\ell)=1$ result), so the three-string emission $p_{1,3}$ of the quartic vertex is forced, exactly as $\ell_3$ is forced by $(3,1)$, which reads $[Q,\ell_3]+\ell_2\circ_1\ell_2=0$ and has a nonzero second term, and as $p_{1,2}$ is forced by $(2,2)$: without it the $(2,2)$ relation would reduce to $[\hat Q,p_{2,2}]=0$ and lose the term $p_{1,2}\circ\ell_2$, which is nonzero on the toy algebra. The smallest augmented sub-family containing $p_{2,2}$ is therefore all augmented components of $\mc V_{\le4}$,
\begin{equation}
  \label{ibk:eq:senses:F4}
  F_4
    = \{Q,\ \ell_2,\ \ell_3,\ p_{1,2},\ p_{1,3},\ p_{2,2}\},
\end{equation}
whose relations at $k+\ell\le4$ are $(1,1),(2,1),(3,1),(1,2),(1,3)$, all of which hold, and $(2,2)$, which fails by $\frac12M_{2,2}\neq0$. The $(2,2)$ relation involves only components of $\mc V_2,\mc V_3,\mc V_4$, all present in $F_4$, so the sub-algebra reading changes nothing about it. Relations with $k+\ell\ge5$ fail for $F_4$ by terms that need $\mc V_{\ge5}$, in the same way as the quartic-truncated $L_\infty$ algebra fails at $(4,1)$: closure is to the order kept, never exact for string field theory, unlike $A_2\subset A_\infty$, which is exact because the product is associative, and unlike quantum $L_\infty$ plus $\Delta$, which is exact because the bivector is $Q$-invariant.

The status of the $(2,2)$ relation is thus settled, and it corrects an earlier conclusion of this project: for the raised products of a cyclic $L_\infty$ algebra the augmented $(2,2)$ relation is \emph{false}. The vanishing step of the earlier check had never been computed; the tensor it asserted to vanish, the raised $[Q,\ell_3]+\ell_2\circ_1\ell_2$, is not the cyclic $L_\infty$ relation, and once the weight of \cref{ibk:sec:gen:pairing} is put where it belongs the augmented relation fails by $\frac12M_{2,2}$, with the $\alg{sl}_2$ Chevalley--Eilenberg algebra as counterexample. The relation that does hold at $(2,2)$ is the non-augmented one, $2P_{2,2}+M_{2,2}=0$, with the no-output component $p_{2,0}=\mc V_2$ and the no-input component $p_{0,2}$ present. Consequently no augmented IBL sub-family on $\mc H$ contains $p_{2,2}$; the sub-structure reading survives only in the non-augmented sense, where the components of $\mc V_{\le4}$ close exactly under every relation with $k+\ell\le4$, including $(2,2)$, because the $(k,\ell)$ relation only involves $\mc V_n$ with $n\le k+\ell$. That is the non-augmented theorem truncated by vertex rank.

\subsection{What this means for the cubic action}
\label{ibk:sec:cubic}

Nothing changes for the cubic action at order $g_s^2$. The action uses only $Q$, $\ell_2$, $p_{1,2}$ and $p_{2,2}$, and its consistency statements, the composite gauge law of $\Sigma$ in particular, use only the relations with $\min(k,\ell)=1$; the $(2,2)$ relation never enters, and the cyclicity of the two-string kinetic term is a property of the action, not of the brace products discussed below. One earlier inference is withdrawn along with the $(2,2)$ relation: that $p_{2,2}$ is $Q$-exact, which had been read off from that relation, no longer follows.

The organising structure is the master equation, and the raised products are its Taylor coefficients about a background. The direction that carries new content is the $L_\infty$ structure on the multi-string space $W=\bigoplus_{n\ge0}\mathrm{Sym}^n\mc H$, and there the result is as follows. Let $\hat{\mathfrak p}=\sum_{k,\ell}\hat p_{k,\ell}$ be any generator of this kind, a finite-order differential operator on a graded-commutative host algebra $A$ (the criterion below is host-independent; whether the family of raised products is realized this way is the subject of the last paragraphs), graded by a formal parameter $t$ that counts inputs, each component $\hat p_{k,\ell}$ carrying the weight $t^{k-1}$ (set $t=1$ at the end), and let $\ell^W_n=\Phi^n_{\hat{\mathfrak p}}$ be its Koszul braces, defined by $F(a):=\e^{-a}\hat{\mathfrak p}\,\e^{a}=\sum_n\frac1{n!}\ell^W_n(a^n)$ for even $a$. The generator is parity-inhomogeneous in every grading of $W$ ($\hat Q$ is odd, $\hat p_{2,2}$ is even), so the classical brace theorem for an odd square-zero operator does not apply, and its conclusion is false here: the braces form a curved $L_\infty$ algebra, in every input-parity sector, if and only if
\begin{equation}
  \label{ibk:eq:cubic:criterion}
  F_{\hat{\mathfrak p}^2}(a)
    = F_{\hat{\mathfrak p}}(a)^2
  \qquad\text{for all even }a\in A\otimes\Lambda,
\end{equation}
with $\Lambda$ an arbitrary auxiliary Grassmann algebra; equivalently, the braces of $\hat{\mathfrak p}^2$ are the shuffle square of those of $\hat{\mathfrak p}$. The auxiliary odd parameters are needed: over even $a\in A$ alone only the relations on even inputs follow, and the derivation $\theta\pd_\theta$ on $\Lambda(\theta)$ satisfies the criterion vacuously while $(\ell^W_1)^2\theta\neq0$. This is not $\hat{\mathfrak p}^2=0$: the free theory has curved-$L_\infty$ braces with $\hat{\mathfrak p}^2\neq0$, and the single component $\hat p_{2,0}$ has non-$L_\infty$ braces although its symbol square has no contraction. For an odd homogeneous operator the criterion reduces to $D^2=0$, which recovers the classical theorem with its hypothesis now stated. Read physically, \eqref{ibk:eq:cubic:criterion} asks that squaring the generator produce nothing beyond the square of what it produces in a background: the multi-string products are then a genuine homotopy algebra rather than a mere list of maps, and it is that, not the nilpotency of the generator, which is the natural condition to test for the multi-string products. The equivalence is proven by a Wick-type lemma separating $\e^{-a}D_1D_2\e^{a}$ from $F_{D_1}F_{D_2}$; the underlying identity and the counterexamples were also checked on the toy algebra, the equivalence itself by hand.

Its single-string content is \cref{ibk:thm:non-augmented}: on single-string inputs the braces reduce to $\ell^W_n(\Psi_1,\dots,\Psi_n)=\sum_\ell p_{n,\ell}$, and the lowest contraction order of \eqref{ibk:eq:cubic:criterion} there is exactly $B_{k,\ell}=0$. No suspension sign enters the one-output components: on the host built from the physical strand variables they are the ordinary coderivation extension, the physical space being already the shifted one, and this was checked at arity up to three. For the components with two or more outputs there is no single graded-commutative polynomial host in the strand variables, with the same parity-homogeneous raising map in every slot of a given type and an even grading parameter, against which a convention could be fixed --- that is the next point, which also lists the hosts not yet tried --- and their ordering convention is bookkeeping.

Whether the leading-order relations on \emph{multi-string} inputs also follow from the single-string ones is open, and it is open in a specific way. At that order the criterion contains, besides the single contraction of \cref{ibk:sec:compositions}, terms in which one derivative of the left factor lands on a derivative of $a$ produced by the right factor; these are absent on single strings and cancel for odd homogeneous operators. Whether they cancel here cannot be decided on an operator host of the kind used above. Take any graded-commutative polynomial algebra with one generator per strand, of whatever parities, with the same parity-homogeneous raising map in every slot of a given type and an even grading parameter: the operator realising the family then either projects one of the two slot symmetries to zero --- the outputs or the inputs, depending on which parity the generators carry --- or replaces the propagator contraction of the products by a bare index sum, because two raising maps in a row compose to the identity. The argument is a parity count on the single contraction, and the three hosts on which it could be run fail in exactly the way it predicts. What is untested, and is where the question survives, is the symbol calculus with its own contraction, a host with separate generators for inputs and outputs, a non-polynomial product, and an odd parameter in a host that is not graded commutative.

Two facts about pairings survive at the level of the brace products, independently of which relations the braces satisfy. No degree-diagonal pairing $\bigoplus_nc_n\omega_n$ on $W$ makes $\omega_W(a_0,\ell^W_n(a_1,\dots,a_n))$ cyclic: already at arity two the term coming from $\hat p_{2,0}$ has no partner. This is a hand computation of one component, not machine-checked, and the component it turns on, $\hat p_{2,0}$, is one the sign checks above did not exercise. Cyclicity is restored by the cotangent doubling $T^*W$, whose action $\mc S(\Phi,\Phi^\vee)=\langle\Phi^\vee,\e^{-\Phi}\hat{\mathfrak p}\,\e^{\Phi}\rangle$ solves the classical master equation exactly when the braces satisfy the curved $L_\infty$ relations, i.e.\ under \eqref{ibk:eq:cubic:criterion}. The two-string kinetic term of the cubic action is, in this picture, a re-weighted slice of the brace products (again up to the ordering convention of the components with two outputs, which nothing here fixes), and it is not the brace structure of any operator: the arity-one brace of such an operator would have to be its own arity-two brace's generator, and a one-line check on $\Psi$ shows it is not.

One question is left open, and it is the one worth recording here because nothing above addresses it. In the basis used throughout this note the gauge transformations of the auxiliary multi-string fields $\Sigma_k$ are simple. It is then natural to ask whether some gauge can be \emph{unfixed} for the $\Sigma_k$, in the same sense in which unfixing the gauge of $\Psi$ replaces the gauge-fixed kinetic operator $c_0^+L_0^+$ by $Q$, and whether the products one would reach that way are more general than the raised products of a single $L_\infty$ algebra. We have no construction to propose, and the question is not settled here.



\subsection*{The overview note}

The fourth block is the section of the overview note on the \IBL{} structure: the definition and the master equation in the conventions of~\cite{Cieliebak:2015:HomologicalAlgebraRelated}, the relations at low signature, the embedding and the projection between \Loo{} and \IBL{} algebras, homotopy transfer and strictification, and a list of convention traps.

\noindent\status{Verified} The degree and Euler-characteristic bookkeeping $\abs{\mathfrak{p}_{k, \ell, g}} = - 2 d (k + g - 1) - 1$, $\chi = 2 - 2 g - k - \ell$, and the \Loo{} relations as the truncation to one output and genus zero, in the degree convention of~\cite{Cieliebak:2015:HomologicalAlgebraRelated}, not in the physical ghost number.\record{001-cubic-sft-overview-note (R5)}

The displayed relation at signature $(2, 2, 0)$ is the form of the relation of~\cite{Cieliebak:2015:HomologicalAlgebraRelated}, which stands; for the raised products it is not implied by the \Loo{} relations without the operations with no input or no output, as stated with the first block.
The paragraph on homotopy transfer read the canonical-model theorem as the elimination of the auxiliary fields and called the minimal model quasi-inverse to strictification; it is corrected by an override along the lines of \cref{ch11:sec:transfer}.
The note identifies the formal variable $\hbar$ of~\cite{Cieliebak:2015:HomologicalAlgebraRelated} with the coupling; the grading analysis of \cref{ch11:sec:mc} identifies the coupling with their second variable $\tau$ at genus zero and keeps $\hbar$ as a separate formal parameter.

\begin{quote}
  \small
  \textbf{Translation box: overview note $\to$ write-up.}
  \begin{itemize}
    \setlength{\itemsep}{0.3ex}
    \item $\mtwo = \mathfrak{p}_{2, 2, 0}$ $\to$ $\sprod{2}{2}$; $\ell_k = \mathfrak{p}_{k, 1, 0}$, $Q = \mathfrak{p}_{1, 1, 0}$ as in the paper.
    \item Degrees $\abs{\cdot}$ are the suspended degrees of~\cite{Cieliebak:2015:HomologicalAlgebraRelated} on $C[1]$, with $\abs{\hbar} = 2 d$; at genus zero they match the physical ghost number through the strand degree $3 - \Ngh$, at $d = \tfrac{1}{2}$.
    \item $\hbar$ $\to$ a formal variable, identified with $\gs$ in the note and not in \cref{ch11:sec:mc}; $\tau = 1$ in the note, $\tau = \gs$ at genus zero in \cref{ch11:sec:mc}.
    \item $\omega_2 = \omega_1^{\otimes 2}$ $\to$ the $2$-string pairing $\pair[2]{\cdot}{\cdot}$.
    \item Reality: none used in the block; the write-up's convention is $\Sigma_k^\star = (-1)^k \Sigma_k$, with the sign $-1$ for odd $k$.
  \end{itemize}
\end{quote}

\subsection{The \texorpdfstring{$\IBL$}{IBL-infinity} structure}
\label{ov:sec:IBL}

\subsubsection{Definition and the master equation}

An $\IBL$-algebra \cite{Cieliebak:2015:HomologicalAlgebraRelated} of degree $d$ is a graded module $C$ with operations
\begin{equation}
	\pklg : E_kC \to E_\ell C,\qquad k,\ell\ge1,\ g\ge0,\qquad
	|\pklg| = -2d(k+g-1)-1,
	\label{ov:eq:IBLdeg}
\end{equation}
where $E_kC$ is the degree-$k$ piece of the reduced symmetric algebra on the suspension $C[1]$,
and $g$ is the genus, the associated surface having Euler characteristic $\chi=2-2g-k-\ell$.
The structure relations are packaged in the master equation
\begin{equation}
	\boxed{\;\hat{\mathfrak{p}}\circ\hat{\mathfrak{p}}=0,\qquad
	\hat{\mathfrak{p}}=\sum_{k,\ell\ge1,\,g\ge0}
	\hat{\mathfrak{p}}_{k,\ell,g}\,\hbar^{\,k+g-1}\,\tau^{\,k+\ell+2g-2},\;}
	\label{ov:eq:IBLmaster}
\end{equation}
with $\hbar,\tau$ formal bookkeeping variables ($|\hbar|=2d$, $|\tau|=0$); under the
finiteness condition satisfied in SFT one sets $\tau=1$, and $\hbar$ plays the role of the
genus-counting $\gs$. The hat denotes the extension of an operation to all of the symmetric
algebra by applying it to every $k$-subset and symmetrising. The product $\mtwo$ of
\cref{p:sec:hs:quartic} is $\mathfrak{p}_{2,2,0}$; the $\Loo$ brackets are
$\ell_k=\mathfrak{p}_{k,1,0}$, the BRST differential $Q=\mathfrak{p}_{1,1,0}$.
(Degree/Euler bookkeeping checked symbolically: C5a--c.)

Spelled out at low signature in the classical ($g=0$) sector, the master equation gives the
familiar $\Loo$ relations on the single-output ($\ell=1$) operations,
\begin{equation}
	\hat\ell_1\circ\hat\ell_1=0,\qquad
	\hat\ell_1\circ\hat\ell_2+\hat\ell_2\circ\hat\ell_1=0,\qquad
	\hat\ell_1\circ\hat\ell_3+\hat\ell_2\circ\hat\ell_2+\hat\ell_3\circ\hat\ell_1=0,
	\label{ov:eq:Loorels}
\end{equation}
together with a two-output ($\ell=2$) sector. The genuine low-order CFL relations involving
two-output operations are, in order of signature (\cite{Cieliebak:2015:HomologicalAlgebraRelated}, Lemma 2.6 table), the chain-map
relation for the cobracket $\mathfrak{p}_{1,2,0}$ at signature $(1,2,0)$,
\begin{equation}
	\hat{\mathfrak{p}}_{1,1,0}\circ_1\mathfrak{p}_{1,2,0}
	+\mathfrak{p}_{1,2,0}\circ_1\hat{\mathfrak{p}}_{1,1,0}=0,
	\label{ov:eq:IBLchain}
\end{equation}
and the structure equation for $\mtwo=\mathfrak{p}_{2,2,0}$ at signature $(2,2,0)$,
\begin{equation}
	\boxed{\;
	\hat{\mathfrak{p}}_{1,1,0}\circ_1\mathfrak{p}_{2,2,0}
	+\mathfrak{p}_{2,2,0}\circ_1\hat{\mathfrak{p}}_{1,1,0}
	+\mathfrak{p}_{1,2,0}\circ_1\mathfrak{p}_{2,1,0}
	+\hat{\mathfrak{p}}_{2,1,0}\circ_1\hat{\mathfrak{p}}_{1,2,0}=0,
	\;}
	\label{ov:eq:IBLrels}
\end{equation}
where $\circ_1$ is partial composition feeding one output into one input. With
$Q=\mathfrak{p}_{1,1,0}$, $\ell_2=\mathfrak{p}_{2,1,0}$, and the cobracket
$\mathfrak{p}_{1,2,0}$, relation \eqref{ov:eq:IBLrels} says that $\mtwo=\mathfrak{p}_{2,2,0}$ is a
\emph{$Q$-homotopy} for the cobracket--bracket composition
$\mathfrak{p}_{1,2,0}\circ_1\ell_2+\hat{\mathfrak{p}}_{2,1,0}\circ_1\hat{\mathfrak{p}}_{1,2,0}$,
i.e. that composition is $Q$-exact (equals
$-[\,Q,\mathfrak{p}_{2,2,0}\,]$), \emph{not} zero. There is in particular \emph{no} CFL relation
of the form ``$\ell_3$ commutes with $\mtwo$''; the only objects coupling to $\mtwo$ at this
order are $Q$, $\ell_2$, and the cobracket $\mathfrak{p}_{1,2,0}$. (The signature ordering
$(1,1,0)\prec(1,2,0)\prec(2,1,0)\prec(1,1,1)\prec(1,3,0)\prec(2,2,0)\prec\cdots$ places
$\ell_3=\mathfrak{p}_{3,1,0}$ \emph{after} $(2,2,0)$, so $\ell_3$ does not enter the
$\mtwo$ structure equation.) (Verified against \cite{Cieliebak:2015:HomologicalAlgebraRelated}: C6.) The schematic $\IBL$ action
is $S=\omega\circ\hat{\mathfrak{p}}(\e^\Psi)$ with a total symplectic form
$\omega=\omega_1+\omega_2$, $\omega_2:\Hilb^{\otimes2}\otimes\Hilb^{\otimes2}\to\C$.

\subsubsection{\texorpdfstring{$\Loo\Rightarrow\IBL$}{L to IBL} and
\texorpdfstring{$\IBL\Rightarrow\Loo$}{IBL to L}}

The two directions are clean and recorded in CFL \cite{Cieliebak:2015:HomologicalAlgebraRelated}, §2:
\begin{itemize}[leftmargin=*]
	\item \textbf{$\Loo\Rightarrow\IBL$ (embed).} Any $\Loo$-algebra $(C,\{\ell_k\})$ is an
	$\IBL$-algebra with $\mathfrak{p}_{k,1,0}=\ell_k$ and all other operations zero. This is
	trivially true.
	\item \textbf{$\IBL\Rightarrow\Loo$ (project).} Setting $\pklg=0$ for $\ell\ge2$ or $g>0$
	reduces the master equation \eqref{ov:eq:IBLmaster} to the standard $\Loo$ axioms; equivalently
	one projects any $\IBL$-algebra onto its $\ell=1,g=0$ operations. This always yields an
	$\Loo$-algebra but \emph{discards} the cobracket ($\ell=2$) and higher-genus data.
\end{itemize}
For the present programme the relevant statement is the on-shell equivalence: the HS/$\IBL$
action reduces to the $\Loo$ action on integrating out $\Sig$, and the $\Loo$ data are recovered
by the projection above. The conjecture that the full non-polynomial $\Loo$ closed SFT is
equivalent on-shell to an $\IBL$ SFT is precisely a statement
about this pair of maps being mutually inverse in an appropriate homotopy sense; the precise
condition is open (\cref{ch:open-problems}).

\subsubsection{Homotopy transfer (out) and strictification (in)}

Strictification and the elimination of the auxiliary fields are inverse to one another.

\subparagraph{Homotopy transfer --- the canonical model.} CFL's
Theorem~1.3/6.1 \cite{Cieliebak:2015:HomologicalAlgebraRelated} states that every $\IBL$-algebra over a field of characteristic
zero is homotopy equivalent to an $\IBL$-structure on its homology
$H=H(C,\mathfrak{p}_{1,1,0})$, with the transferred operations $\mathfrak{q}_{K,L,G}$ given by
an explicit recursive formula built from a projection $\pi:C\to H$, an inclusion
$\mathfrak{f}_{1,1,0}:H\hookrightarrow C$, and a contracting homotopy. It
is the $\IBL$ counterpart of the minimal-model theorem for $\Loo$: every off-shell mode is integrated out and only the cohomology remains.
Integrating out the auxiliary fields alone is a different transfer, along a contraction of the auxiliary fields and their equations of motion, and the minimal model is a further transfer reached from either theory (\cref{ch11:sec:transfer}). Münster--Sachs
\cite{Muenster:2014:HomotopyClassificationBosonic} make this concrete for closed SFT: choosing a (harmonious) Hodge
decomposition is the same as choosing a gauge, and the transferred minimal model on the BRST
cohomology has multilinear maps equal to the S-matrix amplitudes (a sum over Feynman graphs
built from the vertices, the propagator homotopy $h$, and the loop bivector
$g^{-1}=h\circ\omega^{-1}$).

\subparagraph{Strictification --- integrating auxiliary fields in.} The inverse direction is the
strictification theorem \cite{Borsten:2021:DoubleCopyHomotopy}: every $\Loo$-algebra is quasi-isomorphic to a
strict one (a differential graded Lie algebra, with $\ell_i=0$ for $i\ge3$) on a larger space
obtained by adjoining an \emph{infinite tower of auxiliary fields} that eliminate all
interaction vertices of arity $\ge4$. At the Lagrangian level this trades each higher vertex for
cubic vertices joined by auxiliary-field propagators (the ``blow-up'' of an $n$-point contact
vertex into $s/t/u$-channel trees). This is exactly the abstract content of the HS step:
$\ell_3$ is dualised into a cubic coupling of $\Sig$ to $\Psi^{\otimes2}$ plus a quadratic
kinetic term for $\Sig$. Strictification and the elimination of the auxiliary fields are
quasi-inverse: original $\Loo$ $\hookrightarrow$ strict $\Loo$ (in), and strict $\Loo$
$\twoheadrightarrow$ original $\Loo$ (out); the minimal model, original $\Loo$ $\twoheadrightarrow$ its cohomology, is a further transfer, whose genus-zero operations are the tree-level on-shell amplitudes (\cite{Kajiura:2007:NoncommutativeHomotopyAlgebras} for $\Aoo$ algebras).

\subparagraph{Where strictification stops.} The published strictification theorem
\cite{Borsten:2021:DoubleCopyHomotopy} applies to $\Loo$- (and $\Aoo$-, $C_\infty$-) algebras, i.e. to the
single-output sector $\mathfrak{p}_{k,1,0}$. The closed-string $\ell_3$, once HS-dualised,
produces an auxiliary field with a \emph{two-output} kinetic structure, i.e. an $\IBL$ vertex
$\mathfrak{p}_{1,2,0}$ (cobracket); the strict theory is then a differential involutive bi-Lie
(``dIBL'') algebra, not merely a dgLa. Whether an $\IBL$ analogue of the strictification theorem
exists --- a strict-dIBL counterpart for the cobracket sector --- is not established
\cite{Borsten:2021:DoubleCopyHomotopy}; this is the precise sense in which the $\Loo$ strictification covers only
part of the closed-string structure.

\subsubsection{Convention traps}

Several conventions must be tracked when importing CFL/Münster--Sachs formulae into SFT.
\begin{itemize}[leftmargin=*]
	\item \textbf{Suspension shift.} CFL build $E_kC$ on $C[1]$, so their degree $|c|$ is
	$\deg c-1$. The Münster--Sachs papers \cite{MunsterSachs2013,Muenster:2014:HomotopyClassificationBosonic} likewise
	place fields in degree zero after $A_c=s^{-2}\tilde A_c$. Translating into the ghost-number
	grading requires the inverse shift.
	\item \textbf{Formal $\hbar$ of degree $2d$.} The CFL $\hbar$ in \eqref{ov:eq:IBLmaster} is a
	formal variable carrying the genus weight; it is \emph{not} the physical $\hbar$ and should
	be identified with $\gs$ in the genus expansion.
	\item \textbf{$\omega_1$, $\Delta$ live only in the cyclic-DGA model.} CFL do not use a
	symplectic form or BV Laplacian in the abstract $\IBL$ definition; the pairing
	$\Mean{\cdot,\cdot}$ and its inverse appear only in the cyclic DGA construction (their
	§10--12). Identifying $\omega_1$ with that pairing and $\Delta=\omega_1^{-1}$ is part of
	\emph{building} the $\IBL$ structure from the string-field algebra, not a feature of the
	abstract axioms.
	\item \textbf{One $\omega$ per sector in QOCHA.} In Münster--Sachs \cite{MunsterSachs2013}
	the quantum open-closed homotopy algebra carries a \emph{single} odd symplectic form per
	sector ($\omega_o$ for open, $\omega_c$ for closed); there is \emph{no} factorised form
	$\omega_2=\omega_1^{\otimes2}$ anywhere. 
\end{itemize}



\subsection*{The pedagogical note}

The fifth and sixth blocks are three sections of a pedagogical note: the two presentations of the quartic structure, as a genus-zero \IBL{} algebra on the one-string space and as a cyclic \Loo{} algebra on the space enlarged by the two-string field; their physical ghost-number grading without suspension; and the count of suspension units for the pairing of $\Sigma_2$.
The section of the note on the reading of the $(2, 2, 0)$ relation, which carries a withdrawn argument, is not reproduced (\cref{ch:dead-ends}).

\noindent\status{Verified} The \IBL{} action reproduces the cubic action at order $\gs^2$, and the equation of $\Sigma$ follows by stripping the non-degenerate $2$-string pairing; it uses the full graded symmetry of $\vertex{4}$, a property of the vertex, and is checked on a finite model.\record{003-i2-ibl-action (R8)}

\noindent\status{Proven} The physical ghost number~\eqref{ped:eq:ngformula} of the structure maps, the same result as in the first block.\record{009-a3-ibl-pedagogical (R23)}

\noindent\status{Verified} The $2$-string pairing is graded symmetric in the physical statistics, hence not an odd symplectic form; demanding graded antisymmetry forces one shift of the whole field $\Sigma_2$.\record{002-omega2-construction (R7)}

Several statements of the note were superseded by later results, and the block carries overrides for them: the $(2, 2, 0)$ relation of the raised products, the attribution of the elimination of $\Sigma$ to a transfer of \IBL{} structures, the status of the morphism $(\phi_1, \phi_2)$, and the inverse needed by the propagator route.
Two more points are corrected here.
First, presentation (a) reads $\Sigma$ as the two-string component $\Psi_2$ of one Maurer--Cartan element; with the weight of~\cite{Cieliebak:2015:HomologicalAlgebraRelated} that component has ghost number five, and $\Sigma_2$ at ghost number four is the two-string component only after the conjugation~\eqref{ch11:eq:conjugation}, with the genus label $-1$ (\cref{ch11:sec:mc}).
Second, the shifted form~\eqref{ped:eq:omega2symp} is written with one suspension per tensor factor, $s \otimes s$; the result is one shift of the whole field $\Sigma_2$, as the sixth block explains, and the antisymmetric form so obtained is even in every uniform grading.

\begin{quote}
  \small
  \textbf{Translation box: pedagogical note $\to$ write-up.}
  \begin{itemize}
    \setlength{\itemsep}{0.3ex}
    \item $\mtwo = \Delta \circ \ell_3$ $\to$ $\sprod{2}{2}$; $V_n$ $\to$ $\vertex{n}$; $V_4(\Sig, \Sig)$ as the kinetic term of $\Sig$ $\to$ $\vertex{4}(\Sigma_2, \Sigma_2) = \pair[2]{\Sigma_2}{\sprod{2}{2} \Sigma_2}$.
    \item $\Sig$ $\to$ $\Sigma_2$; $\Sym^2 \Hilb_0$ $\to$ the level-matched $2$-string states $\Hn{2}$, quotiented by $\ker \sprod{2}{2}$.
    \item $\omega_1$, $G_2 = \omega_1^{\otimes 2}$ $\to$ $\pair{\cdot}{\cdot}$, $\pair[2]{\cdot}{\cdot}$; $\Delta = \omega_1^{-1}$ the inverse pairing.
    \item Degrees: the note's $d = \tfrac{1}{2}$ dictionary, strand degree $3 - \Ngh$, is the one of the first block.
    \item Exchange sign: the sixth block, in its source, took the one-string pairing antisymmetric ($s = -1$); the override restores $s = + 1$, and the symmetry of $G_2$, which depends on $s^2$ only, is unchanged.
    \item Reality: none used in the blocks; the write-up's convention is $\Sigma_k^\star = (-1)^k \Sigma_k$, with the sign $-1$ for odd $k$.
  \end{itemize}
\end{quote}

\subsection{Two presentations of one classical structure}
\label{ped:sec:two}

\subsubsection{Presentation (a): the genus-zero \texorpdfstring{$\IBL$}{IBL} algebra on
\texorpdfstring{$\Hilb_0$}{H}}
\label{ped:sec:two:ibl}

An $\IBL$-algebra of degree $d$ \cite{Cieliebak:2015:HomologicalAlgebraRelated} is a graded module $C$ with operations
$\pklg:E_kC\to E_\ell C$ ($k,\ell\ge1$, $g\ge0$), where $E_kC$ is the degree-$k$ part of the
reduced symmetric algebra, satisfying the master equation $\hat{\mathfrak p}\circ
\hat{\mathfrak p}=0$. The operations relevant at genus zero are
\begin{equation}
	\underbrace{\mathfrak p_{1,1,0}=Q}_{\text{differential}},\qquad
	\underbrace{\mathfrak p_{k,1,0}=\ell_k}_{\Loo\ \text{brackets}},\qquad
	\underbrace{\mathfrak p_{1,2,0}}_{\text{cobracket}},\qquad
	\underbrace{\mathfrak p_{2,2,0}=\mtwo}_{\text{HS product}}.
	\label{ped:eq:ops}
\end{equation}
Although the individual operations act on tensor powers of $C=\Hilb_0$, the $\IBL$-algebra
itself is a structure on the (reduced) \emph{symmetric algebra} $S\Hilb_0=\bigoplus_k E_k
\Hilb_0$: the $\pklg$ assemble into a single coderivation $\hat{\mathfrak p}$ on
$S\Hilb_0$, and the master field is accordingly a sum over multi-string components,
\begin{equation}
	\Psi=\sum_{n\ge1}\Psi_n,\qquad \Psi_n\in\Sym^n\Hilb_0,
	\label{ped:eq:masterfield}
\end{equation}
exactly as in the Münster--Sachs coderivation expansion $\mathfrak D=\sum_n x^{n-1}D^{(n)}$
\cite[\S IV.B]{MunsterSachs2013}. In this presentation the HS auxiliary is \emph{not} a new
field: $\Sig\equiv\Psi_2$ is the two-string component of the one master field
\eqref{ped:eq:masterfield}.

The string bracket $\ell_3:\Hilb_0^{\otimes3}\to\Hilb_0$ of the quartic
vertex, $V_4(A_1,A_2,A_3,A_4)=\Mean{A_1|\ell_3(A_2,A_3,A_4)}$ (the $c_0^-$-inserted
pairing, cf.\ Conventions), gives rise to the
\emph{two-output} product by raising one index with the inverse symplectic form
$\Delta=\omega_1^{-1}$:
\begin{equation}
	\mtwo\equiv\mathfrak p_{2,2,0}=\Delta\circ\ell_3:\ \Hilb_0^{\otimes2}\to\Hilb_0^{\otimes2},
	\qquad \Ngh(\mtwo)=+2,
	\label{ped:eq:mtwo}
\end{equation}
equivalently $V_4(A_1,A_2,A_3,A_4)=G_2\big(A_1\otimes A_2,\,\mtwo(A_3\otimes A_4)\big)$ with
the product pairing $G_2=\omega_1\otimes\omega_1$, which is the definition~\eqref{p:eq:conv:sprod} of $p_{2,2}$. (The ghost number $\Ngh(\mtwo)=+2$ is derived in \cref{ped:sec:grading:ngh}.) The point of presentation (a) is that $\mtwo$ has \emph{two} inputs and
\emph{two} outputs; it is not an $\Loo$ bracket, and no \emph{independent} auxiliary field
is in sight. The schematic $\IBL$ action pairs the
codifferential with the symplectic data,
\begin{equation}
	S=\omega\circ\hat{\mathfrak p}(\e^\Psi)
	\qquad\text{(schematic: action/EOM consistency established at $O(\gs^2)$;
	the abstract form is open)},
	\label{ped:eq:iblaction}
\end{equation}
and the cubic action at quartic order, \eqref{p:eq:hs4:action-tilde}, is what \eqref{ped:eq:iblaction} produces at $O(\gs^2)$ once the $\ell_3$
quartic is traded by the Hubbard--Stratonovich transformation.

\subsubsection{Presentation (b): the cyclic \texorpdfstring{$\Loo$}{L} algebra on
\texorpdfstring{$\Hilb_0\oplus\Sym^2\Hilb_0$}{the enlarged space}}
\label{ped:sec:two:loo}

Enlarge the field space to $\Hilb_0\oplus\Sym^2\Hilb_0$, with fields $(\Psi,\Sig)$. Now every
product is \emph{single-output}: the two outputs of $\mtwo$ are re-read as one output valued
in the new summand $\Sym^2\Hilb_0$, i.e. the second output becomes the $\Sig$-slot. The
products are read off the vertices of \eqref{p:eq:hs4:action-tilde}:
\begin{equation}
	W_3(A_2,B_1,C_1):=\Mean{A_2\,|\,\mtwo(B_1,C_1)},\qquad
	W_2(A_2,B_2):=\Mean{A_2\,|\,\mtwo(B_2)},
	\label{ped:eq:W3W2}
\end{equation}
with $A_2,B_2\in\Sym^2\Hilb_0$ and $B_1,C_1\in\Hilb_0$: $W_3$ is the asymmetric cubic vertex
coupling one $\Sig$ leg to two $\Psi$ legs, and $W_2$ is the $\Sig$ kinetic vertex. Together
with the ordinary $\Psi$-sector data $(Q,\ell_2)$ these are the brackets of an
$\Loo$-structure on $\Hilb_0\oplus\Sym^2\Hilb_0$; \eqref{p:eq:hs4:action-tilde} is its cyclic action
(cyclicity holding with respect to the shifted pairing, \cref{ped:sec:two:caveat}). This
is the theory that is manifestly cubic and that one quantises.

\subsubsection{The dictionary: what is established at \texorpdfstring{$O(\gs^2)$}{O(gs2)},
and two candidate \texorpdfstring{$\Loo\Leftrightarrow\IBL$}{L-IBL} correspondences}
\label{ped:sec:two:dict}

It is tempting to summarise the relation between (a) and (b) as \enquote{the classical
$\Loo\Leftrightarrow\IBL$ correspondence}. This note deliberately does \emph{not} do so:
no such correspondence exists as a theorem --- for $\Loo$-algebras strictification is a
theorem \cite{Borsten:2021:DoubleCopyHomotopy}, but its $\IBL$ counterpart, in either direction, is precisely
what the programme wants to establish. What is actually available is (i) a set of
$O(\gs^2)$ facts, listed below, and (ii) two candidate general correspondences,
each with a sharply stated open core.

\subparagraph{Established at $O(\gs^2)$.}
\begin{itemize}[leftmargin=*]
	\setlength{\itemsep}{0.6ex}

	\item \textbf{(a)$\to$(b): integrating $\Sig$ in.} Trading the two-output
	$\mathfrak p_{2,2,0}$ for a single-output product on the enlarged space, at the cost of
	promoting $\Psi_2$ to an independent field, \emph{is} the HS step itself
	(\cref{p:sec:hs:quartic}). This direction is unconditional and explicit.

	\item \textbf{(b)$\to$(a): integrating $\Sig$ out.} Eliminating $\Sig$ by its algebraic
	equation of motion (below) collapses the enlarged theory back to the data on $\Hilb_0$
	with the $\ell_3$ quartic restored: at $O(\gs^2)$ this is the Gaussian elimination
	(saddle substitution $=$ propagator transfer) and, at the amplitude
	level, the telescoping $V_4\circ V_4^{-1}\circ V_4=V_4$. It is
	\emph{conditional}: it needs $\mtwo$ injective on $\Sym^2\Hilb_0$ (open
	at CFT level; defused by the kernel quotient $\Sig\in\Sym^2\Hilb_0/\ker V_4$), and the diagrammatic route additionally needs
	the inverse of $V_4$ on the quotient, level by level in a truncation with stable kernels; a bounded inverse is not needed.
\end{itemize}

\subparagraph{What CFL actually provide (and what they do not).} CFL Theorem~1.3/6.1
\cite{Cieliebak:2015:HomologicalAlgebraRelated} is the \emph{minimal-model} theorem: every $\IBL$-algebra
$(C,\{\pklg\})$ over a field of characteristic $0$ carries an $\IBL$-structure on its
\emph{homology} $H(C,\mathfrak p_{1,1,0})$, homotopy equivalent to the original, with an
explicit recursive transfer formula built from a Hodge splitting $C=H\oplus B\oplus A$.
It is a statement about passing to $Q$-homology --- it is \emph{not} a strictification
theorem, and it is \emph{not} the $\Sig$-elimination: integrating $\Sig$ out keeps $\Psi$
fully off-shell. The analogue for the latter is Corollary~6.2 (transfer of an
$\IBL$-structure along any \emph{chain homotopy equivalence}), applied to $\Loo$ algebras, since the cubic action is not known to carry an $\IBL$ structure: when $\mtwo$ is invertible
on the quotient of $\Sym^2\Hilb_0$ by its kernel, the pair formed by $\Sig$ and its equation of motion is contractible on the classical complex, and
the $\Loo$ transfer along that contraction is the algebraic home of \enquote{integrating out the auxiliary
field} (\cref{ch11:sec:transfer}). The hypothesis is the invertibility of $\mtwo$ on that quotient, level by level; it does not follow from Theorem~6.1.

\subparagraph{Candidate correspondence 1 (restricted: auxiliary field via HS).} Integrating
a single auxiliary field in/out through the HS transformation produces a specific,
constrained form of $\IBL$ data: for example, the gauge transformation of $\Sig$ is
completely determined by that of $\Psi$ (the composite law of
\cref{p:sec:hs:quartic}), and $\Sig$ couples only through $V_4$. This direction is operationally
straightforward. Its core question --- whether the product $\mtwo=\Delta\circ\ell_3$
obtained from HS actually \emph{satisfies the $\IBL$ relations}, in particular the
$(2,2,0)$ structure equation \eqref{ped:eq:rel220}, as a consequence of the $\Loo$ relations
of the $\ell_n$ --- has a two-part answer at genus zero.
With the operations of CFL's definition only, at least one input and one output, the $(1,2,0)$ relation holds, $(1,2,0)\Leftrightarrow[Q,\ell_2]=0$, and so does every relation with one input or one output, but the $(2,2,0)$ relation is \emph{not implied} by the $\Loo$ relations: a cyclic $\Loo$ algebra built on $\mathfrak{sl}_2$ violates it.
With the operations with no input or no output included, the raised family satisfies every genus-zero relation if and only if the $\ell_n$ form a cyclic $\Loo$ algebra (\cref{p:sec:ibl:general}).
An earlier version of this paragraph stated that the $(2,2,0)$ relation follows from the $\Loo$ relations; that statement is withdrawn (\cref{ch:dead-ends}).

\subparagraph{Candidate correspondence 2 (general: repackaging, no auxiliary field).} Given
$\Loo$ products $\ell_n=:\mathfrak p_{n,1}$, applying $\Delta$ once produces two-output
products $\mathfrak p_{n-1,2}:=\Delta\circ\ell_n$; iterating $\Delta$ produces an entire
family $\{\mathfrak p_{m,n}\}$ built from the $\ell_n$ \emph{alone}. Nothing is
integrated in or out --- this is pure repackaging, and the question applies whether there
is a single field or several. The \textbf{open core}: do the $\Loo$ relations of the
$\ell_n$ imply the $\IBL$ relations for $\{\mathfrak p_{m,n}\}$? One may then enlarge:
introduce fields $\Psi_n\in\Sym^n\Hilb_0$ and ask the \emph{reverse} --- given a family
of $\IBL$ products, can one write $\Loo$ products on the enlarged vector space
$\bigoplus_n\Sym^n\Hilb_0$, and do the $\IBL$ relations imply the $\Loo$ relations there?
This general correspondence is related to the HS mechanism but plausibly more general;
its goal is not necessarily a cubic action but more general classes of SFT actions.

\subparagraph{The morphism data at $O(\gs^2)$.} The expected comparison morphism has leading
components $\phi_1=$ the inclusion $\Hilb_0\hookrightarrow\Hilb_0\oplus\Sym^2\Hilb_0$ and
$\phi_2=\cdot^{\otimes2}$ mapping into the $\Sym^2$ summand, $\Psi\mapsto\Psi^{\otimes2}$
--- the Yang--Mills-toy pattern (\cref{ov:sec:explicit:YM}). Its gauge-map component is the composite law (\cref{p:sec:hs:quartic}).
On the classical complex, the fields and their equations of motion, $(\phi_1,\phi_2)$ is the transfer map of a cyclic $\Loo$ quasi-isomorphism, under the invertibility of $\mtwo$ on the quotient (\cref{p:sec:all:homotopy-transfer}); off the classical complex it fails the morphism identity with two inputs, so it is not an $\Loo$ morphism of the full algebras.

\subparagraph{The equation of motion and the auxiliarity criterion.} Varying \eqref{p:eq:hs4:action-tilde} in
$\Sig$ gives, using that the kinetic bilinear $V_4$ is graded-symmetric on the even sector
(\cref{p:eq:conv:sprod}),
\begin{equation}
	\mtwo(\Sig)=\mtwo(\Psi^{\otimes2}),
	\label{ped:eq:Sigeom}
\end{equation}
which is solved by $\Sig=\Psi^{\otimes2}$. The solution is \emph{unique} iff
$\ker\mtwo=0$ on the level-matched $\Sym^2\Hilb_0$; any symmetric kernel
component of $V_4$ decouples --- it appears in neither the kinetic nor the source term --- so
defining $\Sig\in\Sym^2\Hilb_0/\ker V_4$ makes $\mtwo$ injective by construction. Then $\Sig$ is purely auxiliary and (b)$\to$(a) is exact.

\subsubsection{The caveat that must be carried: cyclicity needs the shifted
\texorpdfstring{$\omega_2^{\rm symp}$}{omega-2}}
\label{ped:sec:two:caveat}

The word \enquote{cyclic} in presentation (b) is not free. A cyclic $\Loo$-algebra requires
an \emph{odd, graded-antisymmetric} (symplectic) form on the enlarged space. On the
$\Psi$-sector this is the usual $\omega_1=(\cdot,c_0^-\cdot)_{\rm bpz}$. On the $\Sig$-sector,
however, the naive candidate $\omega_1^{\otimes2}$ is graded \emph{symmetric} in physical
statistics (\cref{ped:sec:omega2}), so it is \emph{not} a symplectic form on
$\Sym^2\Hilb_0$. The antisymmetric symplectic partner is the one-unit-shifted object
\begin{equation}
	\omega_2^{\rm symp}=\omega_1^{\otimes2}\circ(s\otimes s)
	\qquad\text{on }\big(\Sym^2\Hilb_0\big)[1],
	\label{ped:eq:omega2symp}
\end{equation}
in the Batalin--Vilkovisky (BV) presentation. Thus presentation (b), read as a genuinely
\emph{cyclic} $\Loo$-algebra, takes the shifted form \eqref{ped:eq:omega2symp} as a
prerequisite: it does not obviate the open question of the two-output symplectic form,
it presupposes its resolution in the BV presentation.

\subparagraph{Everything is CFT data; the shift is bookkeeping, not input.} A construction
that needed the suspension as an \emph{ingredient} would be suspect: if the theory is
well defined, every object must be expressible in worldsheet CFT data alone. It is. The
complete list of CFT inputs of this note is: the $c_0^-$-inserted overlap
$\omega_1(A,B)=\bra{A}c_0^-\ket{B}_{\rm bpz}$ with the normalisation
$\mean{c_{-1}\bar c_{-1}c_0^-c_0^+c_1\bar c_1}=-\tfrac12$ (Conventions); its two-strand
product $G_2=\omega_1^{\otimes2}$, i.e.\ the double overlap with one $c_0^-$ per strand;
the geometric vertices $V_3,V_4$ (surface states integrated over their vertex regions);
and the level-matched projection defining $\Hilb_0$. Every operation is built from these:
$\ell_3$ from $V_4$ by \eqref{ped:eq:mtwo}'s dual, $\Delta$ as the inverse overlap, $\mtwo=
\Delta\circ\ell_3$. The shifted $\omega_2^{\rm symp}$ of \eqref{ped:eq:omega2symp} is
\emph{not a new CFT object}: it is the \emph{same} double overlap $G_2$, evaluated on the
same states --- the suspension $s$ changes no correlator, no state and no operation. What
it changes is a single declaration of \emph{bookkeeping parity}: in the BV complex, $\Sig$
is assigned Grassmann parity $\epsilon(\Sig)+1$ instead of the physical
$\epsilon(\Sig)=\Ngh(\Sig)\bmod2$, and with that assignment the \emph{same} $G_2$
is graded antisymmetric, though even as an operation (the exchange \emph{sign convention} flips,
\cref{ped:eq:cs}). The BV presentation therefore buys a uniform language (cyclicity, master
equation) at the price of a bookkeeping parity that differs from the physical one by one
unit on the $\Sig$-sector; the physics --- correlators, action values, EOM --- is
presentation-independent. \Cref{ped:sec:omega2} spells out the mechanism (why the two
$c_0^-$'s make $G_2$ symmetric, and why exactly one shift unit flips it).

\subsection{Physical ghost-number grading: undoing the suspension}
\label{ped:sec:grading}

A standing objection to the suspension is that it hides the physics. This section
determines whether the two source frameworks suspend (they do), and then states every
$\IBL$ component relation in the physical ghost-number grading, which is unambiguous and is
the grading of all final objects.

\subsubsection{Do CFL and Münster--Sachs suspend? Yes --- fields and structure maps}
\label{ped:sec:grading:susp}

\begin{itemize}[leftmargin=*]
	\setlength{\itemsep}{0.6ex}

	\item \textbf{Cieliebak--Fukaya--Latschev \cite{Cieliebak:2015:HomologicalAlgebraRelated}.} The symmetric algebra $E_kC$
	is built on the \emph{suspension} $C[1]$, so their degree $|c|$ is the physical degree
	minus one: $|c|=\deg c-1$. The structure maps carry the suspended degree
	\begin{equation}
		|\pklg|=-2d(k+g-1)-1,
		\label{ped:eq:cfldeg}
	\end{equation}
	with $d$ a fixed integer (for SFT on a $(2n-1)$-dimensional contact target,
	$d=n-3$). This is a \emph{homological} grading: the differential $\mathfrak p_{1,1,0}$
	has $|\mathfrak p_{1,1,0}|=-1$, i.e. it \emph{lowers} the CFL degree, opposite to the
	BRST operator which \emph{raises} ghost number.

	\item \textbf{Münster--Sachs \cite{MunsterSachs2013,Muenster:2014:HomotopyClassificationBosonic}.} The closed-string
	field is placed in degree zero after the shift $A_c=s^{-2}\tilde A_c$; the open field
	after $A_o=s^{-1}\tilde A_o$. The odd symplectic form is
	$\omega_c=(\cdot,c_0^-\cdot)_{\rm bpz}\circ(s^2\otimes s^2)$: it is odd and graded
	antisymmetric \emph{by construction of the degree shift} --- \enquote{the apparent
	symmetry of the BPZ product is removed by the suspension}. The physical
	ghost-number grading of $\tilde A_c$ is separate and kept implicit.
\end{itemize}
So both frameworks suspend, and both fold the load-bearing sign of the odd symplectic form
into the shift. This is exactly the bookkeeping to be removed for the final
statements.

\subsubsection{Physical ghost numbers of the fields and structure maps}
\label{ped:sec:grading:ngh}

In the physical grading the fields carry $\Ngh(\Psi)=2$ and $\Ngh(\Sig)=4$, fixed so that
the couplings saturate the sphere ghost anomaly. The ghost numbers of the structure maps
follow from the sphere anomaly ($\Ngh=6$; one $c_0^-$ makes $\omega_1$ pair states with
$\Ngh$ sum $5$; the genus-zero $n$-point vertex saturates at $\sum_i\Ngh(A_i)=2n$):
\begin{equation}
	V_n(A_1,\dots,A_n)=\omega_1\big(A_1,\ell_{n-1}(A_2,\dots,A_n)\big)
	\ \Longrightarrow\
	\Ngh(\ell_n)=3-2n.
	\label{ped:eq:ellngh}
\end{equation}
Raising one input of $\ell_3$ into an output with $\Delta=\omega_1^{-1}$ shifts the ghost
number by $+5$, \emph{not} by $+1$: since $\omega_1$ saturates at $\Ngh$ sum $5$, its
inverse is a bivector $\Delta=\sum_\alpha u_\alpha\otimes v_\alpha$ with
$\Ngh(u_\alpha)+\Ngh(v_\alpha)=5$, so
\begin{equation}
	\mtwo(A\otimes B)=\sum_\alpha \ell_3(u_\alpha,A,B)\otimes v_\alpha
	\quad\Longrightarrow\quad
	\Ngh(\mtwo)=\Ngh(\ell_3)+5=+2.
	\label{ped:eq:Deltashift}
\end{equation}
Equivalently, by \eqref{ped:eq:mtwo}'s defining relation $V_4=G_2(\,\cdot\,,\mtwo\,\cdot\,)$:
the two outputs of $\mtwo$ each pair via one $\omega_1$ of $G_2$ with an external
$\Ngh=2$ state, so each output sits at $\Ngh=3$, and
$\Ngh(\mtwo)=(3+3)-(2+2)=+2$ --- the same out-minus-in rule that gives
$\Ngh(Q)=3-2=+1$ and $\Ngh(\ell_n)=3-2n$. The same raising applied to $\ell_2$ gives
$\Ngh(\mathfrak p_{1,2,0})=\Ngh(\ell_2)+5=+4$, which we confirm independently below by
ghost-number balance. The physical grading is collected in \cref{ped:tab:ngh}.

\begin{table}[htp]
	\centering
	\begin{tabular}{lccc}
		\toprule
		object & signature $(k,\ell,g)$ & role & $\Ngh$ \\
		\midrule
		$\Psi$ & --- & dynamical field & $+2$ \\
		$\Sig$ & --- & auxiliary field & $+4$ \\
		$Q=\mathfrak p_{1,1,0}$ & $(1,1,0)$ & BRST differential & $+1$ \\
		$\ell_2=\mathfrak p_{2,1,0}$ & $(2,1,0)$ & bracket & $-1$ \\
		$\mathfrak p_{1,2,0}$ & $(1,2,0)$ & cobracket & $+4$ \\
		$\ell_3=\mathfrak p_{3,1,0}$ & $(3,1,0)$ & higher bracket & $-3$ \\
		$\mtwo=\mathfrak p_{2,2,0}$ & $(2,2,0)$ & HS two-output product & $+2$ \\
		$\mathfrak p_{1,1,1}$ & $(1,1,1)$ & genus-1 involutivity homotopy & $+2$ \\
		\bottomrule
	\end{tabular}
	\caption{Physical ghost numbers of the fields and low-signature $\IBL$ structure maps.
		The map $\Ngh(\ell_n)=3-2n$ is \eqref{ped:eq:ellngh}; every $\ell\ge2$ (multi-output)
		map picks up $+5$ per $\Delta$-raising \eqref{ped:eq:Deltashift}.
		$\Ngh(\mathfrak p_{1,2,0})=+4$ and $\Ngh(\mathfrak p_{1,1,1})=+2$ are also forced by
		ghost-number balance of the $(2,2,0)$ and $(1,1,1)$ relations
		(\cref{ped:sec:grading:rels}; checked, \texttt{Hq4}).}
	\label{ped:tab:ngh}
\end{table}

\subparagraph{The two-parameter physical formula.} Collecting the two rules above --- every
output at $\Ngh=3$, every input at $\Ngh=2$ (out-minus-in), and $+1$ per handle --- the
physical ghost number of the general $\IBL$ structure map is the two-parameter (three, with
genus) count
\begin{equation}
  \boxed{\;\Ngh(\pklg)=3\ell-2k+g\;.}
  \label{ped:eq:ngformula}
\end{equation}
The genus-zero backbone $\Ngh(\mathfrak p_{k,\ell,0})=3\ell-2k$ is out-minus-in; its
$\ell=1$ slice is $\Ngh(\ell_k)=3-2k$ \eqref{ped:eq:ellngh}, and each extra output is a
$\Delta$-raising, $+5=(+3\text{ output})-(-2\text{ input})$ \eqref{ped:eq:Deltashift}. The genus
coefficient is fixed by $\Ngh$-homogeneity of the loop relations: the first relation
carrying genus, \eqref{ped:eq:rel111}, forces $\Ngh(\mathfrak p_{1,1,1})=+2=3-2+1$, so the
per-handle shift is $+1$ --- each handle is an operadic self-gluing that identifies one
output leg ($\Ngh=3$) with one input leg ($\Ngh=2$), and the ghost excess $(3-2)=+1$ is what
it costs.\footnote{The \emph{geometric} vertex $V_{g,k+\ell}$ saturates at
$\sum_i\Ngh(A_i)=2(k+\ell)$ for \emph{all} $g$ (the anomaly $6-6g$ exactly cancels the
$6g-6+2(k+\ell)$ moduli $b$-insertions), so the genus dependence resides in the algebraic
loop-operation normalisation; the loop relations pin its coefficient to $+1$. See
\cref{ped:tab:ngh} for every value used here.} Formula \eqref{ped:eq:ngformula} reproduces the whole
of \cref{ped:tab:ngh}: $\ell_2\!=\!-1$, $\ell_3\!=\!-3$, $\mathfrak p_{1,2,0}\!=\!+4$,
$\mtwo\!=\!+2$, $\mathfrak p_{1,1,1}\!=\!+2$.

\subparagraph{The suspended \texorpdfstring{$\to$}{->} physical bridge is \emph{per slot}.}
The CFL degree \eqref{ped:eq:cfldeg} is \emph{independent of the output number} $\ell$, whereas the physical grading \eqref{ped:eq:ngformula} \emph{grows} with $\ell$.
There is therefore no uniform per-map offset $|\pklg|=\pm\Ngh(\pklg)+\text{const}$ valid across $(k,\ell)$ at fixed $d$ --- and this absence is \emph{expected, not a defect}.
It is the ordinary consequence of the maps being multilinear ($k$ inputs) \emph{and} multi-output ($\ell$ outputs), each slot carrying its own suspension.
The mechanism is exactly the $L_\infty$ one \cite{Hohm:2017:LinftyAlgebrasField}: in the $\ell$-picture $\deg\ell_n=n-2$ grows with the number of inputs, while in the suspended $b$-picture $\deg b_n=-1$ is \emph{constant}, the two related by a per-slot suspension $s$ that puts one unit on the output and one on each input (\cite{Hohm:2017:LinftyAlgebrasField}, \S 3.1--3.2).
For the $\IBL$ maps the same per-slot shift acts on \emph{every} strand, the $k$ inputs as well as the $\ell$ outputs, as CFL's suspension $C[1]$ does.
Concretely, give a strand of ghost number $n$ the CFL degree $3-n$, with the sign reversed because the CFL grading is homological (\cref{ped:sec:grading:susp}).
The degree of a map is the degree of its outputs minus that of its inputs, so
\begin{equation}
  |\pklg|
    = \sum_{\text{out}}(3-n_{\rm out})-\sum_{\text{in}}(3-n_{\rm in})
    = 3\ell-3k-\Ngh(\pklg)
    = 3\ell-3k-(3\ell-2k+g)
    = -(k+g),
  \label{ped:eq:cflmatch}
\end{equation}
where the second step uses $\Ngh(\pklg)=\sum n_{\rm out}-\sum n_{\rm in}$ and the third \eqref{ped:eq:ngformula}.
Comparing with \eqref{ped:eq:cfldeg}, $-2d(k+g-1)-1=-(k+g)$ holds for every $(k,\ell,g)$ exactly when $d=\tfrac12$.
For example, $\ell_2$ has $\Ngh=-1$ and gets degree $3-6+1=-2$, and $\mtwo$ has $\Ngh=+2$ and gets $6-6-2=-2$, both equal to $-(k+g)$ with $(k,g)=(2,0)$.
The value is unique: allowing a general affine strand degree $\sigma n+\alpha$, with separate constants $\alpha$ on inputs and outputs, and matching \eqref{ped:eq:cfldeg} for all $k,\ell\ge1$, $g\ge0$ forces $\sigma=-1$, both constants equal to $3$, and $d=\tfrac12$; the $k$-slope, the $g$-slope and the constant term each give $d=\tfrac12$ separately, and without the sign reversal no $d$ works.
The match therefore holds for the degree, at $d=\tfrac12$, with three qualifications.
First, $d=\tfrac12$ lies outside CFL's hypotheses, which take $d$ integer ($|\hbar|=2d=1$ is odd here); it is not CFL's $d=n-3$ for any contact target, and CFL's integer-$d$ results (homotopy theory, Weyl-algebra correspondence) do not apply as stated.
Second, the regrading matches integer degrees, not slot symmetries: the strand degree $3-\Ngh$ has parity $\epsilon+1$ on every strand, so CFL's $E_kC$ would be graded symmetric in $\epsilon+1$ on the inputs, whereas the physical $\pklg$ are graded symmetric in $\epsilon$ on the inputs and in $\epsilon+1$ on the outputs (the physical $\Psi$, with $\Ngh=2$, gets the odd CFL degree $1$); the symmetries still need the per-slot parity dictionary.
Third, the $g$-slope of the match uses the $+g$ of \eqref{ped:eq:ngformula}, which is the ghost number of the self-gluing operation and is anchored by an independent computation only at $g=1$ (the $(1,1,1)$ relation); for $g\ge2$ it is the uniform ghost-additive self-gluing convention, consistent but not independently re-derived.
At genus zero the $k$-slope and the constant already fix $d=\tfrac12$ without it.
A non-integer $d$ is thus the correct outcome of the per-slot dictionary, not a symptom of a wrong one.
The upshot for the reader: the physical statements need no $d$; they are the $\Ngh$-balanced
operator equations of the next subsection.

\subsubsection{The component relations in physical ghost numbers}
\label{ped:sec:grading:rels}

Each $\IBL$ component relation is a homogeneous operator equation of a definite ghost number;
every term must carry the same $\Ngh$. We list the genus-zero chain up to $(2,2,0)$ and
verify homogeneity (all checked in \texttt{Hq4}).
\begin{subequations}
	\label{ped:eq:rels}
	\begin{align}
		\label{ped:eq:rel110}
		\text{(1,1,0):}\qquad
		&
		Q^2=0,
		&&\Ngh=+2;
		\\
		\label{ped:eq:rel210}
		\text{(2,1,0):}\qquad
		&
		Q\circ_1\ell_2+\ell_2\circ_1\hat Q=0,
		&&\Ngh=0;
		\\
		\label{ped:eq:rel120}
		\text{(1,2,0):}\qquad
		&
		\hat Q\circ_1\mathfrak p_{1,2,0}+\mathfrak p_{1,2,0}\circ_1\hat Q=0,
		&&\Ngh=+5;
		\\
		\label{ped:eq:rel310}
		\text{(3,1,0):}\qquad
		&
		Q\circ_1\ell_3+\ell_2\circ_1\ell_2+\ell_3\circ_1\hat Q=0,
		&&\Ngh=-2.
	\end{align}
\end{subequations}
The $(2,2,0)$ structure equation --- the one that governs $\mtwo$ --- is, in physical ghost
numbers,
\begin{equation}
	\boxed{\;
	\hat Q\circ_1\mathfrak p_{2,2,0}
	+\mathfrak p_{2,2,0}\circ_1\hat Q
	+\mathfrak p_{1,2,0}\circ_1\ell_2
	+\hat\ell_2\circ_1\hat{\mathfrak p}_{1,2,0}=0,
	\qquad \Ngh=+3,
	\;}
	\label{ped:eq:rel220}
\end{equation}
where $\circ_1$ feeds one output into one input. Ghost-number balance of \eqref{ped:eq:rel220}
is not automatic: the first two terms have $\Ngh=1+2=+3$, and demanding the last two match
\emph{forces} $\Ngh(\mathfrak p_{1,2,0})=+4$ (then $\mathfrak p_{1,2,0}\circ\ell_2$ and
$\ell_2\circ\mathfrak p_{1,2,0}$ both have $\Ngh=4+(-1)=+3$) --- in agreement with the
direct values $\Ngh(\ell_2)+5$ from \eqref{ped:eq:Deltashift} and out-minus-in
$(3+3)-2=+4$. This is a genuine, physically
transparent determination of the cobracket's ghost number, reproduced by \texttt{Hq4}
without any suspension input. The genus-1 relation
\begin{equation}
	\text{(1,1,1):}\qquad
	Q\circ_1\mathfrak p_{1,1,1}+\mathfrak p_{1,1,1}\circ_1 Q
	+\ell_2\circ_2\mathfrak p_{1,2,0}=0,
	\qquad \Ngh=+3,
	\label{ped:eq:rel111}
\end{equation}
then forces $\Ngh(\mathfrak p_{1,1,1})=+2$ (the entry of \cref{ped:tab:ngh}).



\medskip
The sixth block is the section of the same note on the count of suspension units for the pairing of $\Sigma_2$: two insertions of $c_0^-$ make the $2$-string pairing even as an operation, and one shift of the field makes it antisymmetric.

\subsection{The grading of \texorpdfstring{$\omega_2$}{omega-2}: one unit, not two}
\label{ped:sec:omega2}

A persistent grading puzzle is this: $\Sig\in\Sym^2\Hilb_0$ has
\emph{two} tensor factors, so one expects the symplectic partner $\omega_2^{\rm symp}$ to
carry \emph{two} degree-shift units, one per factor. It carries \emph{one}. The resolution
is that the \enquote{two} and the \enquote{one} are counting two
different $\Z_2$ data, and only one of them is the suspension.

\subparagraph{Two $c_0^-$ make the operation even.} The single-string form
$\omega_1(A,B)=\Mean{A|c_0^-|B}_{\rm bpz}$ has the exchange sign $\omega_1(B,A)=s\,(-1)^{\epsilon(A)\epsilon(B)}\,\omega_1(A,B)$ on its support, with $s=+1$ in the conventions of the paper, \eqref{p:eq:conv:pairing-symmetry}. The product form
$G_2=\omega_1\otimes\omega_1$ inserts $c_0^-$ \emph{twice}, hence is even as an operation, and
the two factors $s$ multiply to $s^2=1$, whichever sign the one-string form carries, so that $G_2$ is a graded-\emph{symmetric} bilinear:
\begin{equation}
	G_2(Y,X)=(-1)^{\epsilon(X)\epsilon(Y)}\,G_2(X,Y)
	\quad\text{on }\operatorname{supp}(G_2),
	\qquad \epsilon=\Ngh\bmod2
	\label{ped:eq:G2sym}
\end{equation}
(the two signs $s$ from the two $\omega_1$'s multiply to $+1$). \emph{This is
where the expected \enquote{two} lives}: the two tensor factors each supply a $c_0^-$, and the two
$c_0^-$'s make $G_2$ even-as-an-operation. The \enquote{two} is a
statement about the $c_0^-$ count, i.e. about the \emph{operation parity} --- not about the
suspension.

\subparagraph{One suspension unit makes it a symplectic form.} A graded-symmetric $G_2$ is not a
symplectic form. To obtain the graded-antisymmetric form asked for by a cyclic
$\Loo$/BV structure on the $\Sig$-sector, one shifts the \emph{field} $\Sig$ by a single
suspension unit, \eqref{ped:eq:omega2symp}. The shift is \emph{per field}, not per tensor factor,
because $\Sig$ is \emph{one} element of the enlarged complex $C=\Sym^2\Hilb_0$ --- one entry
of the form's complex, hence one $C\to C[1]$ --- irrespective of how many $\Hilb_0$-factors
build it. Concretely, the exchange sign flips iff
\begin{equation}
	G_2(Y,X)=-(-1)^{(\epsilon(X)+c_s)(\epsilon(Y)+c_s)}\,G_2(X,Y)
	\quad\text{on support}
	\ \Longrightarrow\ c_s=1,
	\label{ped:eq:cs}
\end{equation}
so exactly one odd shift unit is required and consistent.

\subparagraph{Reconciling the two counts.} The expectation \enquote{two factors $\to$ two units} conflates
the two $\Z_2$ data. The correct bookkeeping is:
\begin{center}
	\begin{tabular}{lcc}
		\toprule
		datum & source & count for $\Sig\in\Sym^2\Hilb_0$ \\
		\midrule
		operation parity of $G_2$ & number of $c_0^-$ insertions & \textbf{two} ($\Rightarrow$ even $\Rightarrow$ symmetric) \\
		suspension of the form's complex & $C\to C[1]$, per field & \textbf{one} ($\Rightarrow$ antisymmetric after shift) \\
		\bottomrule
	\end{tabular}
\end{center}
The expected \enquote{two} is real, but it is the number of $c_0^-$'s (which makes
$G_2$ even as an operation), while the suspension that
converts $G_2$ into a symplectic form is one unit because $\Sig$ is a single field. The even
$s^{-2}$ integer shift of Münster--Sachs' $A_c=s^{-2}\tilde A_c$ is parity-irrelevant; the
single $[1]$ that flips the symmetry is the suspension of the form's complex, a
BV-presentation choice, not a theorem about $\Sig$.



\subsection{Kinetic operators, the canonical model, and the auxiliary field}
\label{ch11:sec:transfer}

\noindent\status{Conditional(PH1, PH5, PH14, PH27) -- mixed}\record{114-ibl-remarks-check (R111)}
This section answers a question raised by the canonical-model theorem of~\cite{Cieliebak:2015:HomologicalAlgebraRelated}: does the theorem say that only the BRST operator carries physical states, and that every other quadratic term defines an auxiliary field?
We find that the first half is the content of the theorem, and that the second half holds for $\Sigma_2$ in the cubic action but does not follow from the theorem.
The setting is genus zero with a coupling of ghost number zero.
The statement on \IBL{} algebras assumes their axioms, with integer $d$, which the raised family does not satisfy, so it describes \IBL{} algebras and not the cubic action.
The statement on $\Sigma_2$ holds on the classical complex at quartic order assuming that $\sprod{2}{2}$ is invertible on the quotient by its kernel, level by level with stable kernels; the statement on the BRST spectrum assumes in addition that the kernel of $\vertex{4}$ vanishes off shell at generic momentum or is annihilated by $\vertex{5}$ and $\vertex{6}$, that the feedback of the auxiliary exchanges is not expressible through the diagonal values $\vertex{n}(\Psi^n)$ and boundary terms, at order $\gs^4$, and that the vertices kept form the truncation of a consistent genus-zero \Loo{} family.

\paragraph{Ghost number of a kinetic operator.}

A quadratic term $\pair[n]{\Phi}{K \Phi}$ for an $n$-strand field of total ghost number $N$ is nonzero only when the ghost numbers saturate the $n$-string pairing, five per strand:
\begin{equation}
  \label{ch11:eq:kinetic-balance}
  2 N + \Ngh(K)
    = 5 n,
  \qquad
  \Ngh(K)
    = 5 n - 2 N.
\end{equation}
The auxiliary field $\Sigma_n$ has its strands at ghost number two, $N = 2 n$, so its kinetic operator has ghost number
\begin{equation}
  \label{ch11:eq:kinetic-n}
  \Ngh(K_n)
    = n
    = 3 n - 2 n
    = \Ngh(\sprod{n}{n}).
\end{equation}
The extension of the BRST operator to $n$ strands,
\begin{equation}
  \hat{Q}_B
    = \sum_{i = 1}^n 1^{\otimes (i - 1)} \otimes Q_B \otimes 1^{\otimes (n - i)},
\end{equation}
with the Koszul signs, has ghost number one, so it can be a kinetic operator only for $n = 1$.
The same holds for any operator of ghost number one acting strand by strand, and for any distribution of the total ghost number $2 n$ among the strands.
For $\Sigma_2$ the kinetic operator is $\sprod{2}{2}$ at quartic order, dressed by $- \tfrac{\gs}{5} \sprod{3}{2}(\Psi, \cdot)$ beyond; the higher $\Sigma_n$ have kinetic operators built from the reduced vertices.

The contrast with an \IBL{} algebra is direct.
There the extension $\hat{\mathfrak{p}}_{1, 1, 0}$ of the differential to the symmetric algebra is a derivation acting on every strand~\cite{Cieliebak:2015:HomologicalAlgebraRelated}, and in the two-string component of the Maurer--Cartan equation it appears at the leading order in $\hbar$, with $\mathfrak{p}_{2, 2, 0}$ only at the next order.
The cubic action has the opposite pattern: no $\hat{Q}_B \Sigma_2$ term, by~\eqref{ch11:eq:kinetic-n}, and $\sprod{2}{2} \Sigma_2$ at leading order.
The difference rests on the coupling: in the \IBL{} algebra the two terms on one component are balanced by $\hbar$, of ghost number $-1$ in the dictionary of \cref{ch11:sec:mc}, while $\gs$ carries no ghost number.

\paragraph{The canonical model.}

Theorem~6.1 of~\cite{Cieliebak:2015:HomologicalAlgebraRelated} states that every \IBL{} algebra $(C, \{\mathfrak{p}_{k, \ell, g}\})$ over a field of characteristic zero is homotopy equivalent to an \IBL{} structure $\{\mathfrak{q}_{k, \ell, g}\}$ on the homology $H = H(C, \mathfrak{p}_{1, 1, 0})$, with $\mathfrak{q}_{1, 1, 0} = 0$, the $\mathfrak{q}$ being built recursively from a splitting $C = H \oplus B \oplus A$ and a contracting homotopy.
Its Corollary~6.2 transfers the structure along any chain homotopy equivalence.
Take $\mathfrak{p}_{1, 1, 0} = Q_B$.
The structure lives on the symmetric algebra, and in characteristic zero the homology of the derivation extension $\hat{Q}_B$ on the $k$-th symmetric power is the symmetric power of the homology (Künneth):
\begin{equation}
  \label{ch11:eq:kunneth}
  H \bigl( \mathrm{Sym}^k C, \hat{Q}_B \bigr)
    = \mathrm{Sym}^k H(Q_B).
\end{equation}
The canonical model therefore lives on the Fock space $\mathrm{Sym} \, H(Q_B)$ built on the one-particle states $H(Q_B)$.
This is the precise sense in which only $\mathfrak{p}_{1, 1} = Q_B$ determines the physical states: no other operation, whatever its numbers of inputs and outputs, adds or removes a state, and all of them become operations $\mathfrak{q}_{k, \ell, g}$ on that space.

The multi-string sectors of the canonical model are not empty: they are $\mathrm{Sym}^k H(Q_B)$, read physically as multi-particle states.
In an \IBL{} algebra no multi-string direction has an operation with $k$ inputs and $k$ outputs as its only quadratic term, because $\hat{\mathfrak{p}}_{1, 1, 0}$ acts on every strand; such operations enter weighted by $\hbar^{k - 1}$, as perturbations of $\hat{\mathfrak{p}}_{1, 1, 0}$, and the theorem transfers them rather than integrating out a field because of them.
So the statement that every other quadratic term defines an auxiliary field is not what the theorem says, and inside the \IBL{} framework it is false: the multi-string directions carry the multi-particle states.

\paragraph{Why $\Sigma_2$ is auxiliary.}

In the cubic action the situation is different.
Where the kinetic operator $\sprod{2}{2}$ is injective, on the quotient by its kernel, the pair formed by $\Sigma_2$ and its equation of motion is acyclic.
Integrating out that pair is a homotopy transfer along a chain homotopy equivalence, the analogue of Corollary~6.2 for \Loo{} algebras, not a passage to homology.
On the classical complex, the fields and their equations of motion, at quartic order, it is the homological-perturbation transfer of \cref{p:sec:all:homotopy-transfer}, with transferred brackets $(Q_B, \gs \ell_2, \gs^2 \ell_3, 0)$ and transfer map $\Psi \mapsto (\Psi, \Psi^2)$; it controls the solutions and their linearization, not the BRST spectrum.
The BRST spectrum is a separate statement: around a background $\Psi_0$, the linearized BRST operator of the master action of the cubic theory is similar to the linearized operator $Q^{(k)}_{\Psi_0}$ of the original theory truncated at order $k$, plus a contractible pair, so that their cohomologies agree in all ghost numbers to the truncation order ($\Sigma_2$ alone, $k \in \{2, 3\}$, with the kernel of $\vertex{4}$ quotiented; \cref{p:sec:gauge:spectrum}).
Under the assumptions stated at the head of this section, $\Sigma_2$ therefore carries no physical state, and only $Q_B$ does.
For the higher fields $\Sigma_n$, $n \ge 3$, the same is expected, since the ghost number~\eqref{ch11:eq:kinetic-n} does not depend on the order, but their dressed kinetic operators and the transfer of the whole tower have not been examined.

The physical meaning is a difference, not a similarity, with \IBL{} algebras.
At quartic order the cubic theory and the original one have the same physical states, as Theorem~6.1 would suggest, but by a different mechanism: in an \IBL{} algebra the multi-string directions carry $\hat{Q}_B$ and the states $\mathrm{Sym}^k H(Q_B)$; in the cubic action $\Sigma_2$ carries $\sprod{2}{2}$ and no state, and the multi-particle states are carried by products of $\Psi$.

\paragraph{Strictification, elimination and the minimal model.}

Three transfers are involved, and the earlier texts conflated two of them.
Strictification goes from the original theory to a larger quasi-isomorphic model with simpler operations: the cubic action, or the bar--cobar construction of \cref{ch:bar-cobar}.
The elimination of the auxiliary fields goes back from the cubic action to the original theory, a homotopy transfer along a contraction whose kernel is acyclic; for $\Sigma_2$ this is the statement above, on the classical complex and under the invertibility of $\sprod{2}{2}$, and it inverts the strictification up to homotopy.
The minimal model goes from the original theory to $H(Q_B)$, a transfer onto the cohomology in which every off-shell mode is integrated out with the contracting homotopy as propagator; its genus-zero brackets are the tree-level on-shell amplitudes (\cite{Kajiura:2007:NoncommutativeHomotopyAlgebras} for $\Aoo$ algebras; the construction is the same for the closed-string \Loo{} algebra).
All three are quasi-isomorphisms, and the minimal model is the end point of the transfer, reached from either theory.

The facts used here were reproduced in finite models, the balance~\eqref{ch11:eq:kinetic-n} by enumeration of the strand distributions and the equivalence of the self-adjointness of $\sprod{2}{2}$ with the symmetry of $\vertex{4}$ on exact data, and~\eqref{ch11:eq:kunneth} was checked on a random graded complex, the last with a control in which a dropped Koszul sign makes $\hat{Q}_B^2 \ne 0$.

\subsection{The Maurer--Cartan form and the gauge components}
\label{ch11:sec:mc}

This section asks whether the cubic theory is a Maurer--Cartan element of the raised family, read as an \IBL{} structure in the sense of~\cite{Cieliebak:2015:HomologicalAlgebraRelated}.
It first fixes the gradings, then computes the components of the Maurer--Cartan expression through order $\gs^2$ and compares them with the equations of motion of the cubic action, then identifies the coupling that the components require and what it does to the theory, and finally lists the components of the gauge transformation.
The setting, for the whole section, is the following.
The operations are the genus-zero raised family $\sprod{k}{\ell}$ with $k, \ell \ge 1$, truncated at $k + \ell \le 4$, which is the order $\gs^2$.
The extension $\hat{\phi}$ of an operation $\phi$ with $k$ inputs to a word in the symmetric algebra applies $\phi$ to every $k$-subset of the letters, as in~\cite[(2.1)]{Cieliebak:2015:HomologicalAlgebraRelated}.
Input strands (fields, gauge parameters, unfed strands of a field) carry the physical parity and output strands of the raised products the shifted parity $\Ngh + 1$, the parity rule of the raised products; the fields have even strands, $\Psi$ at ghost number two and $\Sigma_2$ with strands at $(2, 2)$.
The variable $\hbar$ is an even, central formal parameter with free powers, whose degree and ghost number are integer labels not tied to the Grassmann parity.
The comparisons below are identities between explicit expressions and use no hypothesis.
Reading the components as a gauge-covariant Maurer--Cartan equation needs more: that the single-contraction part of $\hat{\mathfrak{p}} \circ \hat{\mathfrak{p}}$, extended to words with odd letters, is the square of one operator reproducing the antibracket term by term, which is open; and the relations with two contractions between two operations, which are not implied by the genus-zero relations and are expected to fail for the genus-zero truncation unless genus-one operations are added.

\subsubsection{Grading of the Maurer--Cartan equation}

\noindent\status{Verified}\record{103-ibl-mc-conventions (R103)}
In the weighting of~\cite{Cieliebak:2015:HomologicalAlgebraRelated} the operator is $\hat{\mathfrak{p}} = \sum \hat{\mathfrak{p}}_{k, \ell, g} \hbar^{k + g - 1} \tau^{k + \ell + 2 g - 2}$, and at genus zero the coupling is $\tau$, which weights $\sprod{k}{\ell}$ by $\gs^{k + \ell - 2}$.
A Maurer--Cartan element with one- and two-string components is $\mathfrak{m} = \hbar^{-1} (\Psi_1 + \Psi_2)$, and the equation is the vanishing of the components of
\begin{equation}
  \label{ch11:eq:mc-expansion}
  e^{- \mathfrak{m}} \hat{\mathfrak{p}} (e^{\mathfrak{m}})
    = \sum_{k, \ell} \hbar^{k - 1} \gs^{k + \ell - 2} \sum_{a, b} \frac{\hbar^{- a - b}}{a! \, b!} \bigl[ \hat{\mathfrak{p}}_{k, \ell} (\Psi_1^a \Psi_2^b) \bigr]_{\mathrm{conn}},
\end{equation}
where the connected part keeps the $k$-subsets that take at least one strand from each factor.
A term in which $\sprod{k}{\ell}$ takes $t_i \in \{1, 2\}$ strands from the $i$-th factor $\Psi_2$ has the power $(k - 1) - (a + b)$ of $\hbar$, the power $k + \ell - 2$ of $\gs$, and $L = \ell + \sum_i (2 - t_i)$ output strands.
When the inputs leave one strand of a $\Psi_2$ unfed, that strand stays as an output: this produces the lifts of the one-output operations, among them the lift $\hat{Q} \Psi_2 = (Q_B \otimes 1 + 1 \otimes Q_B) \Psi_2$ of $\sprod{1}{1}$.
Write $y = \Ngh(\Psi_2)$.
The connected terms with $k + \ell \le 4$ landing in one or two strands are the following.
\begin{center}
  \small
  \begin{tabular}{llccc}
    \toprule
    component & term & $\hbar$ & $\gs$ & $\Ngh$ \\
    \midrule
    $E_1$ & $Q_B \Psi_1$, $\tfrac{1}{2} \ell_2(\Psi_1^2)$, $\tfrac{1}{6} \ell_3(\Psi_1^3)$ & $\hbar^{-1}$ & $\gs^{0, 1, 2}$ & $3$ \\
    & $\ell_2(\Psi_2)$, $\ell_3(\Psi_1, \Psi_2)$ & $\hbar^0$ & $\gs^{1, 2}$ & $y - 1$ \\
    $E_2$ & $\hat{Q} \Psi_2$, lifts of $\ell_2(\Psi_1, \cdot)$ and $\tfrac{1}{2} \ell_3(\Psi_1^2, \cdot)$ & $\hbar^{-1}$ & $\gs^{0, 1, 2}$ & $y + 1$ \\
    & $\sprod{1}{2}(\Psi_1)$, $\tfrac{1}{2} \sprod{2}{2}(\Psi_1^2)$ & $\hbar^{-1}$ & $\gs^{1, 2}$ & $6$ \\
    & $\sprod{2}{2}(\Psi_2)$ & $\hbar^0$ & $\gs^2$ & $y + 2$ \\
    & lift of $\ell_3(\Psi_2, \cdot)$ on a second $\Psi_2$ & $\hbar^0$ & $\gs^2$ & $2 y - 3$ \\
    \bottomrule
  \end{tabular}
\end{center}
At $y = 4$ every term of the $L$-strand component has ghost number
\begin{equation}
  \label{ch11:eq:ghost-rule}
  \Ngh
    = 2 L + \ell_p,
\end{equation}
with $\ell_p$ the number of outputs of the product used: the factors carry two per strand, $2 (k + L - \ell_p)$ in all, and the product adds $3 \ell_p - 2 k$.

Two criteria of homogeneity can be asked: all the terms of a component, every order in $\hbar$ together, at one grading, or each order in $\hbar$ separately, since for a formal $\hbar$ each order is a separate equation.
In the degree of~\cite{Cieliebak:2015:HomologicalAlgebraRelated}, with $\abs{\hbar} = 2 d = 1$, both components are homogeneous if and only if $y = 5$, which is their requirement $\abs{\mathfrak{m}_{2, 0}} = 1$ read through the strand degree $3 - \Ngh$.
In the physical ghost number with $\hbar$ and $\gs$ ghost-neutral, the one-string component is homogeneous only for $y = 4$ and the two-string component for no $y$, since its ghost numbers $y + 1$ and $y + 2$ never agree; at $y = 4$ it splits into a sector at ghost number five, the one-output lifts, and a sector at ghost number six, the two-output products.
Giving $\hbar$ the ghost number $x$, the components are homogeneous across the orders in $\hbar$ if and only if
\begin{equation}
  (x, y)
    = (-1, 5),
\end{equation}
so the degree of~\cite{Cieliebak:2015:HomologicalAlgebraRelated} is the physical ghost number with $\Ngh(\hbar) = -1$, component by component.
Order by order, the weighting of~\cite{Cieliebak:2015:HomologicalAlgebraRelated} with $\Psi_2 = \mathfrak{m}_{2, 0}$ requires $y = 5$ for every $x$.
The order-by-order criterion depends on the power of $\hbar$ that dresses the two-string component; across orders the result is $\Ngh(\mathfrak{m}_{2, g}) = 5 + g$ with $\Ngh(\hbar) = -1$, for every dressing, which is the degree $\abs{\mathfrak{m}_{2, g}} = 1 - g$ of their definition.
The gauge parameter $\lambda = \hbar^{-1} (\Lambda_1 + \Lambda_2)$, entering through $\delta e^{\mathfrak{m}} = [\hat{\mathfrak{p}}, \lambda] e^{\mathfrak{m}}$, has $\Ngh(\Lambda_1) = 1$ and, in their degree, $\Ngh(\Lambda_2) = 4$.

The value four for $\Sigma_2$ is recovered by a change of weighting.
Let $\varphi$ multiply the $L$-strand part of the symmetric algebra by $\hbar^{-L}$.
Since $\hat{\mathfrak{p}}_{k, \ell, g}$ takes $k$ strands and returns $\ell$,
\begin{equation}
  \label{ch11:eq:conjugation}
  \varphi^{-1} \, \hat{\mathfrak{p}}_{k, \ell, g} \hbar^{k + g - 1} \, \varphi
    = \hat{\mathfrak{p}}_{k, \ell, g} \hbar^{\ell + g - 1},
  \qquad
  \hat{\mathfrak{p}}(e^{\varphi(\mathfrak{m}')})
    = \varphi \bigl( \hat{\mathfrak{p}}'(e^{\mathfrak{m}'}) \bigr),
\end{equation}
with $\hat{\mathfrak{p}}' = \sum \hat{\mathfrak{p}}_{k, \ell, g} \hbar^{\ell + g - 1}$: the weight with one $\hbar$ per input beyond the first and the weight with one $\hbar$ per output beyond the first are conjugate.
Take $\mathfrak{m}' = \Psi + \Sigma_2$, independent of $\hbar$, that is $\mathfrak{m} = \hbar^{-1} \Psi + \hbar^{-2} \Sigma_2$.
A term of the $L$-strand component then carries $\hbar^{\ell_p - 1}$ in $\hat{\mathfrak{p}}'(e^{\mathfrak{m}'})$, and by~\eqref{ch11:eq:ghost-rule} its total ghost number is $2 L + \ell_p + (\ell_p - 1) x$, the same for all terms if and only if $x = -1$.
In this output-counting form, $(\Psi, \Sigma_2; \Lambda_1, \Lambda_2)$ at ghost numbers $(2, 4; 1, 3)$ is homogeneous, across orders for $\Ngh(\hbar) = -1$ and order by order for any value, and at genus zero the order in $\hbar$ of a term is the number of outputs of its product minus one.
The price is that $\Sigma_2$ receives the genus label $-1$ in the expansion of~\cite{Cieliebak:2015:HomologicalAlgebraRelated}, outside their range $g \ge 0$ and their filtration condition.
In either weighting the two-string component at ghost number four placed as $\mathfrak{m}_{2, 0}$ is excluded.
The obstruction belongs to the raised family: in the structure of Münster and Sachs~\cite{MunsterSachs2013}, whose only two-output datum is the inverse of the odd form, a two-string element at ghost number four is homogeneous with a neutral $\hbar$.
The gradings were measured term by term in a finite model with the ghost numbers of the vertices imposed, with control runs at the integer values $d = 0$ and $d = 1$, which give no solution.

\subsubsection{The components, and the comparison with the equations of motion}

\noindent\status{mixed}\record{103-ibl-mc-conventions (R104)}
We work in the output-counting form, $\mathfrak{m}' = \Psi + \Psi_2$ with the weight $\hbar^{\ell - 1} \gs^{k + \ell - 2}$ and $\Ngh(\Psi, \Psi_2) = (2, 4)$.
On words of even letters, the extension of an operation $\phi$ with $k$ inputs is the differential operator
\begin{equation}
  \hat{\phi}
    = \frac{1}{k!} \, \phi(e_{a_1}, \ldots, e_{a_k}) \, \pd^{a_k} \cdots \pd^{a_1},
\end{equation}
the outputs standing to the left.
With $\Psi = \psi^a e_a$ and $\Psi_2 = s^{a b} e_a e_b$, $s$ symmetric, the derivatives of $\mathfrak{m}'$ are $\pd^a \mathfrak{m}' = \psi^a + 2 s^{a b} e_b$, $\pd^a \pd^b \mathfrak{m}' = 2 s^{a b}$, and the third vanishes.
Conjugating by the exponential and summing over the set partitions of the derivatives gives, with $\hat{\phi} X = (\phi \otimes 1 + 1 \otimes \phi) X$ the extension to two strands of a map with one free input,
\begin{equation}
  \begin{aligned}
    e^{- \mathfrak{m}'} \hat{p}_{1, \ell} e^{\mathfrak{m}'}
      &
      = \sprod{1}{\ell}(\Psi) + \widehat{\sprod{1}{\ell}} \, \Psi_2,
    \\
    e^{- \mathfrak{m}'} \hat{p}_{2, \ell} e^{\mathfrak{m}'}
      &
      = \tfrac{1}{2} \sprod{2}{\ell}(\Psi^2) + \sprod{2}{\ell}(\Psi_2) + \widehat{\sprod{2}{\ell}(\Psi, \cdot)} \, \Psi_2 + \ldots,
    \\
    e^{- \mathfrak{m}'} \hat{p}_{3, \ell} e^{\mathfrak{m}'}
      &
      = \tfrac{1}{6} \sprod{3}{\ell}(\Psi^3) + \sprod{3}{\ell}(\Psi, \Psi_2) + \tfrac{1}{2} \widehat{\sprod{3}{\ell}(\Psi, \Psi, \cdot)} \, \Psi_2 + \widehat{\sprod{3}{\ell}(\Psi_2, \cdot)} \, \Psi_2 + \ldots,
  \end{aligned}
\end{equation}
where the dots have more strands.
Sorting by the number of strands and by the order in $\hbar$, the components through order $\gs^2$ are
\begin{subequations}
  \label{ch11:eq:mc-components}
  \begin{align}
    \label{ch11:eq:M1}
    \mathcal{M}_1
      &
      = Q_B \Psi + \frac{\gs}{2} \, \ell_2(\Psi^2) + \frac{\gs^2}{3!} \, \ell_3(\Psi^3) + \gs \, \ell_2(\Psi_2) + \gs^2 \, \ell_3(\Psi, \Psi_2),
    \\
    \label{ch11:eq:M2-5}
    \mathcal{M}_2^{(5)}
      &
      = \hat{Q} \Psi_2 + \gs \, \widehat{\ell_2(\Psi, \cdot)} \, \Psi_2 + \frac{\gs^2}{2} \, \widehat{\ell_3(\Psi, \Psi, \cdot)} \, \Psi_2 + \gs^2 \, \widehat{\ell_3(\Psi_2, \cdot)} \, \Psi_2,
    \\
    \label{ch11:eq:M2-6}
    \mathcal{M}_2^{(6)}
      &
      = \gs \, \sprod{1}{2}(\Psi) + \frac{\gs^2}{2} \, \sprod{2}{2}(\Psi^2) + \gs^2 \, \sprod{2}{2}(\Psi_2),
  \end{align}
\end{subequations}
at ghost numbers three, five and six, the last at order $\hbar$, together with three- and four-strand components, headed by $\gs^2 \hbar^2 \sprod{1}{3}(\Psi)$ at ghost number nine.
Every relative coefficient is $+ 1$: all the letters are even, so no Koszul sign arises, and the signs of the operations are those of the raised family.
The three equations are separate both by order in $\hbar$ and by ghost number, so their sum vanishes if and only if each does, whatever value $\hbar$ is given, $\hbar = 1$ included.
The terms of the three- and four-strand components and their ghost numbers were checked; their coefficients were derived by hand.
In the native grading of~\cite{Cieliebak:2015:HomologicalAlgebraRelated}, one parity per strand, the symmetric $\Sigma_2$ of the paper vanishes in the two-string part of the symmetric algebra, so the comparison below exists only with the parity rule stated at the head of the section.

The equations of motion of the cubic action at quartic order are~\eqref{p:eq:hs4:eom},
\begin{equation}
  \mathcal{E}_\Psi
    = Q_B \Psi + \frac{\gs}{2} \, \ell_2(\Psi^2) + \frac{\gs^2}{3!} \, \ell_3(\Psi, \Sigma_2),
  \qquad
  \mathcal{E}_\Sigma
    = \sprod{2}{2}(\Psi^2) - \sprod{2}{2}(\Sigma_2),
\end{equation}
with $\Sigma_2$ in the quotient by $\ker \sprod{2}{2}$.
The question is whether, for some identification $\Psi_2 = f(\Sigma_2, \Psi)$, the components are these equations.
The class of $f$ is that of the local redefinitions built from the data of the theory at the order of the truncation: polynomials linear in $\Sigma_2$ and of degree at most two in $\Psi$, built from the fields, the tensor product and the operations of the family, with coefficients power series in $\gs$ and independent of $\hbar$.
A product adds $\ell_p$ to $\Ngh - 2 L$, and $f$ has two strands at ghost number four, so it contains no product:
\begin{equation}
  f
    = \alpha \, \Sigma_2 + \beta \, \Psi^2,
  \qquad
  \alpha
    = \sum_j \alpha_j \gs^j,
  \quad
  \beta
    = \sum_j \beta_j \gs^j.
\end{equation}
In the one-string component, using $\ell_3(\Psi, \Psi^2) = \ell_3(\Psi^3)$,
\begin{equation}
  \mathcal{M}_1 - \mathcal{E}_\Psi
    = \gs \, \alpha \, \ell_2(\Sigma_2) + \gs^2 \bigl( \alpha - \tfrac{1}{6} \bigr) \ell_3(\Psi, \Sigma_2) + \gs \, \beta \, \ell_2(\Psi^2) + \gs^2 \bigl( \beta + \tfrac{1}{6} \bigr) \ell_3(\Psi^3).
\end{equation}
The four maps have different field degrees, so each coefficient must vanish: the first gives $\alpha_0 = 0$ and the second $\alpha_0 = \tfrac{1}{6}$.
The two $\Sigma_2$-terms of $\mathcal{M}_1$ come from the same factor $e^{\Psi_2}$ and carry the same $\alpha$, whereas $\mathcal{E}_\Psi$ has $\ell_3(\Psi, \Sigma_2)$ without $\ell_2(\Sigma_2)$.
In the two-string component at ghost number six,
\begin{equation}
  \mathcal{M}_2^{(6)}
    = \gs \, \sprod{1}{2}(\Psi) + \gs^2 \bigl( \tfrac{1}{2} + \beta \bigr) \sprod{2}{2}(\Psi^2) + \gs^2 \, \alpha \, \sprod{2}{2}(\Sigma_2).
\end{equation}
The cobracket term $\gs \sprod{1}{2}(\Psi)$ is independent of $f$ and linear in $\Psi$, and $\mathcal{E}_\Sigma$ has no such term; without it, the proportionality to $\mathcal{E}_\Sigma$ would need $\alpha + \beta = - \tfrac{1}{2}$, against $\alpha + \beta = 0$ from the one-string component.
The sector at ghost number five has no counterpart among the equations of motion, which have ghost numbers three and six, so it must vanish identically; its parts of degree one and two, $\alpha \hat{Q} \Sigma_2$ and $2 \beta \, Q_B \Psi \cdot \Psi$, vanish only for $\alpha = \beta = 0$.
Finally, the equations of motion are invariant under $\Sigma_2 \to \Sigma_2 + N$ with $N \in \ker \sprod{2}{2}$, while the components change by $\alpha \gs \ell_2(N)$ and $\alpha \hat{Q} N$, so for $\alpha \ne 0$ they are not defined on the quotient.
Hence no $f$ of the class makes the quartic-order equations of motion the components.
The conclusion does not depend on the reading of \enquote{are}: if each component is only required to be a combination of the equations of motion, $f = 0$ makes $\mathcal{M}_1$ such a combination and removes the ghost-five sector, but the term $\gs \sprod{1}{2}(\Psi)$ survives for every $f$, by an argument that counts the ghost number and the field degree of the combinations (not a check); and on a free solution $\Psi = \varepsilon \Psi_0$, $Q_B \Psi_0 = 0$, $\Sigma_2 = \Psi^2$, the component~\eqref{ch11:eq:M2-6} at first order in $\varepsilon$ is $\gs \sprod{1}{2}(\Psi_0) = 0$, false whenever $\vertex{3}(\cdot, \cdot, \Psi_0) \ne 0$.

The components~\eqref{ch11:eq:M1} and~\eqref{ch11:eq:M2-6} are the Euler--Lagrange equations of
\begin{equation}
  \label{ch11:eq:S-MC}
  S_{\mathrm{MC}}(\Psi, \Psi_2)
    = S_{\le 4}(\Psi) + \gs \, \vertex{3}(\Psi_2, \Psi) + \frac{\gs^2}{2} \, \vertex{4}(\Psi_2, \Psi^2) + \frac{\gs^2}{2} \, \vertex{4}(\Psi_2, \Psi_2) + O(\gs^3),
\end{equation}
the vertex functionals evaluated on $e^{\Psi + \Psi_2}$ without the term $\vertex{2}(\Psi_2)$, which has no input; the ghost-five sector is not the variation of any action and comes from the lifts alone.
The cubic action is instead $S_{\le 4}(\Psi) - \tfrac{\gs^2}{4!} \vertex{4}(\Sigma_2 - \Psi^2, \Sigma_2 - \Psi^2)$.
The coupling $\gs \vertex{3}(\Psi_2, \Psi)$ is the common origin of $\gs \ell_2(\Psi_2)$ in~\eqref{ch11:eq:M1} and of $\gs \sprod{1}{2}(\Psi)$ in~\eqref{ch11:eq:M2-6}.

Part of the mismatch is not specific to the cubic action.
The terms $\gs \sprod{1}{2}(\Psi)$ and $\gs^2 \sprod{1}{3}(\Psi)$ are independent of $f$ and present at $\Psi_2 = 0$: with $\Psi_2 = 0$, or any $\Psi_2(\Psi)$ built locally from the family, the original quartic closed string field theory is not a Maurer--Cartan element of the raised family either.
A $\Psi_2(\Psi)$ containing the inverse of $\sprod{2}{2}$ removes the cobracket term at the price of a modified kinetic operator of $\Psi$.
These baseline terms measure the distance between the raised family and any tree-level closed string field theory; removing them means switching off the cobracket on the classical field, the step back to the \Loo{} truncation.
The mismatches specific to the Hubbard--Stratonovich form are the term $\alpha \gs \ell_2(\Sigma_2)$, the coefficient of the source, the ghost-five sector, the $f$-dependent three- and four-strand terms, the missing invariance under the kernel, and, at the level of actions, the coefficient of $\vertex{4}(\Sigma_2, \Sigma_2)$, which would need $\tfrac{\alpha^2}{2} = - \tfrac{1}{4!}$.
The exclusion is a statement within the class of $f$ above; outside it lie redefinitions containing the propagator of $\Sigma_2$ or of $\Psi$, $\hbar$-dependent ones, and families with data added to the $\sprod{m}{n}$.
The identities were checked in a finite model of random graded-symmetric vertices, the components being computed literally in the symmetric algebra, with control runs that remove the obstructions one at a time.

\subsubsection{The coupling required by the Maurer--Cartan form}

\noindent\status{no ledger entry}\record{109-mc-cubic-coupling}
The two mismatches $\gs \ell_2(\Psi_2)$ and $\gs \sprod{1}{2}(\Psi)$ are the two variations of one term, so we add it to the cubic action and ask what it does.
With $\Psi_2 = \alpha \Sigma_2$, the terms of~\eqref{ch11:eq:S-MC} linear in $\Psi_2$ are $\alpha \gs \vertex{3}(\Sigma_2, \Psi) + \tfrac{\alpha \gs^2}{2} \vertex{4}(\Sigma_2, \Psi^2)$, and the mixed term of the cubic action, $\tfrac{2 \gs^2}{4!} \vertex{4}(\Sigma_2, \Psi^2)$, fixes the normalization:
\begin{equation}
  S'
    = \caction[\le 4] + c \, \gs \, \vertex{3}(\Sigma_2, \Psi),
  \qquad
  \alpha
    = c
    = \frac{1}{6}.
\end{equation}
The two mismatches give the same value, since the ratio $1 : \tfrac{1}{2}$ of $\vertex{3}(\Psi_2, \Psi)$ and $\vertex{4}(\Psi_2, \Psi^2)$ comes from the exponential alone; the footnote of the paper on the restricted correspondence states this term with $\Psi_2 = \Sigma_2 / 6$.
The equations of motion of $S'$ are
\begin{equation}
  \mathcal{E}'_\Psi
    = Q_B \Psi + \frac{\gs}{2} \, \ell_2(\Psi^2) + c \, \gs \, \ell_2(\Sigma_2) + \frac{\gs^2}{3!} \, \ell_3(\Psi, \Sigma_2),
  \qquad
  \mathcal{E}'_\Sigma
    = c \, \gs \, \sprod{1}{2}(\Psi) + \frac{\gs^2}{12} \bigl[ \sprod{2}{2}(\Psi^2) - \sprod{2}{2}(\Sigma_2) \bigr].
\end{equation}
At $c = \tfrac{1}{6}$ the two mismatches are supplied and the invariance under the kernel is lost, while the coefficient of the source, the ghost-five sector, the term $\ell_3(\Psi^3)$, the baseline term $\sprod{1}{3}(\Psi)$, the three- and four-strand terms and the action-level coefficient remain.

The added term has one slot filled by $\Psi$ and two by the strands of $\Sigma_2$: it is bilinear in the fields, a mixing of order $\gs$ between $\Psi$ and a field whose propagator, $- 12 \, \sprod{2}{2}^{-1} / \gs^2$, is of order $\gs^{-2}$.
Assuming $\sprod{2}{2}$ invertible, the solution of the equation of $\Sigma_2$ is
\begin{equation}
  \Sigma_2^*
    = \Psi^2 + \frac{12 c}{\gs} \, \sprod{2}{2}^{-1} \sprod{1}{2}(\Psi),
\end{equation}
with a part of order $\gs^{-1}$.
In the full action a term $\gs^{n - 2} \vertex{n}(\Sigma_2^k, \Psi^{n - 2 k})$ then contributes to the eliminated action first at order $\gs^{n - 2 - k}$, so that the order $\gs^p$ receives contributions from every order $n \le 2 p + 4$.
At order $\gs^0$ only $n \le 4$ enters, and the eliminated action is
\begin{equation}
  S'[\Psi]
    = S_2 + 6 c^2 \, W(\Psi) + O(\gs),
  \qquad
  W(\Psi)
    = \pair[2]{\sprod{1}{2}(\Psi)}{\sprod{2}{2}^{-1} \sprod{1}{2}(\Psi)},
\end{equation}
that is $S_2 + \tfrac{1}{6} W$ at $c = \tfrac{1}{6}$: a nonlocal quadratic term, an integral over the momenta of the two strands.
In the quartic truncation alone the orders $\gs$ and $\gs^2$ are $2 S_3 + S_4$, but in the full action the quintic term adds $\tfrac{1}{30} \gs \vertex{5}((\sprod{2}{2}^{-1} \sprod{1}{2} \Psi)^2, \Psi)$ at order $\gs$, and the higher orders add more; only the order $\gs^0$ is a statement about the theory.
Heuristically $W$ is a bubble in which both strands of the auxiliary line join two cubic vertices; since $\sprod{2}{2}^{-1}$ is the inverse of a quartic vertex and not a pair of propagators, the picture is only suggestive.

Assume that $W$ is finite and not identically zero for the closed-string vertices.
This is motivated by its truncation to the tachyon at zero strand momentum, $\vertex{3}(T^3)^2 / \vertex{4}(T^4) \, t^2 \ne 0$, and by the absence of any known structural reason for it to vanish; its finiteness, an integral over the strand momenta, and the invertibility of $\sprod{2}{2}$ on the continuum are not established.
Under this assumption, and with the operator $\sprod{2}{2}$ of the $\Sigma_2$ sector kept, no rescaling of the two coefficients of the $\Sigma_2$ sector, no redefinition of $\Sigma_2$, no redefinition of $\Psi$ that reduces to the identity at $\gs = 0$, and no counterterm in $\Psi$ alone that keeps $Q_B \Psi$ as the order $\gs^0$ part of the equation of $\Psi$ restores closed string field theory at order $\gs^0$ while keeping the two supplied terms.
The counterterm $- 6 c^2 W$ would restore $S_2$, but it puts the same nonlocal term into the zeroth-order equation of $\Psi$, where the Maurer--Cartan equation has only $Q_B \Psi$.
Whether a redefinition of $\Psi$ independent of $\gs$ can absorb $W$ is open for the closed string; in one finite model it cannot, since the free parts differ in rank.
If $\sprod{2}{2}$ has a kernel on which $\vertex{3}$ does not vanish, the coupling also forbids the quotient: the kernel components of $\Sigma_2$ become Lagrange multipliers for $\vertex{3}(N, \Psi) = 0$, and the quotient is available only if $\vertex{3}(\ker \vertex{4}, \cdot) = 0$ on the states of ghost number two.
The cubic action of this write-up does not include this coupling.
The term it would supply, $\gs \sprod{1}{2}(\Psi)$, is a condition that closed string field theory itself does not satisfy: without the coupling it constrains $\Psi$, and with it it is absorbed by $\Sigma_2$ at the price of the term of order $\gs^0$.
The identities of the quartic truncation were checked in finite models; the statements about the full action are at order $\gs^0$ only and assume the property of $W$ stated above.

\subsubsection{Components of the gauge transformation}

\noindent\status{mixed}\record{111-ibl-gauge-components (R110)}
In the output-counting form, with $\mathfrak{m}' = \Psi + \Sigma_2$ and $\lambda' = \Lambda_1 + \Lambda_2$ at ghost numbers $(2, 4; 1, 3)$, the transformation $\delta e^{\mathfrak{m}} = [\hat{\mathfrak{p}}, \lambda] e^{\mathfrak{m}}$ gives
\begin{equation}
  \delta \mathfrak{m}'
    = e^{- \mathfrak{m}'} [\hat{\mathfrak{p}}', \lambda'] e^{\mathfrak{m}'},
\end{equation}
the sum of the connected configurations of~\eqref{ch11:eq:mc-expansion} with one extra factor $\lambda'$ feeding at least one of its strands.
With $a$ copies of $\Psi$ and $b = b_1 + b_2$ copies of $\Sigma_2$, fed singly ($b_1$) or doubly ($b_2$), a term has
\begin{equation}
  \text{coefficient}
    = \frac{1}{a! \, b_1! \, b_2!},
  \qquad
  \hbar^{\ell - 1} \gs^{k + \ell - 2},
  \qquad
  \Ngh
    = 2 L + \ell_p - 1,
\end{equation}
the choice of the fed strand inside a singly fed factor being part of the lift.
There are 26 connected terms on the one- to four-strand components.
The transformation of $\Psi$ has six terms, all at order $\hbar^0$ and ghost number two:
\begin{equation}
  \delta \Psi
    = Q_B \Lambda_1 + \gs \, \ell_2(\Psi, \Lambda_1) + \frac{\gs^2}{2} \, \ell_3(\Psi^2, \Lambda_1) + \gs \, \ell_2(\Lambda_2) + \gs^2 \, \ell_3(\Psi, \Lambda_2) + \gs^2 \, \ell_3(\Sigma_2, \Lambda_1).
\end{equation}
The first three terms are the gauge transformation of closed string field theory, checked against the literal symmetric algebra.
The transformation of $\Sigma_2$ splits by ghost number: seven one-output lifts at ghost number four and order $\hbar^0$, whose term of order $\gs^0$ is $\hat{Q}_B \Lambda_2$, the graded-Leibniz extension of $Q_B$, with coefficient $+ 1$ by the normalization of $\Lambda_2$; and three two-output terms at ghost number five and order $\hbar$, $\gs \sprod{1}{2}(\Lambda_1) + \gs^2 \sprod{2}{2}(\Psi, \Lambda_1) + \gs^2 \sprod{2}{2}(\Lambda_2)$.
Three- and four-strand components are generated, seven and three terms: the two-strand truncation is not preserved.
For $\Lambda_2$ with strands at ghost numbers $(2, 1)$, $\hat{Q}_B \Lambda_2$ has a part with strands at $(3, 1)$, outside the $(2, 2)$ sector of $\Sigma_2$, which is a constraint on $\Lambda_2$ if $\Sigma_2$ is kept in that sector.

The signs of the 26 terms are fixed by the Koszul rule of~\cite[(2.1)]{Cieliebak:2015:HomologicalAlgebraRelated} in the physical parity of the fed and passed letters, by the choice of the family on odd letters, by two normalizations and by one identification: a raised two-output tensor on a mixed sector, such as $\sprod{1}{2}(\Lambda_1)$ with strands at $(2, 3)$, is graded antisymmetric in the physical parity, and the re-signing $(-1)^{\epsilon_{c_1}}$ on its first strand restores the symmetry in the shifted parity that the parity rule requires; the placement of this sign is a convention, not tested by a composition.
No hypothesis enters the list of terms.
What is established is the following.
The list of terms, their orders and the ghost-number rule are verified; the combinatorial coefficients are verified on some rows and derived from the formula above on the others.
The signs and coefficients of the terms at order $\hbar^0$ through order $\gs$, namely $Q_B \Lambda_1$, $\gs \ell_2(\Psi, \Lambda_1)$ and $\gs \ell_2(\Lambda_2)$ in $\delta \Psi$, $\hat{Q}_B \Lambda_2$ and the two lifts of order $\gs$ in $\delta \Sigma_2$, and the three-strand term $\gs \widehat{\ell_2}(\Lambda_2, \Sigma_2)$, are verified on a cyclic differential graded Lie algebra, $\mathfrak{sl}_2 \otimes \Lambda[\theta_1, \theta_2, \theta_3]$, where the transformation preserves the one- and two-strand Maurer--Cartan components at those orders; the sign of $\hat{Q}_B \Lambda_2$ relative to $\gs \ell_2(\Lambda_2)$ is fixed there by $Q_B \ell_2(\Lambda_2) + \ell_2(\hat{Q}_B \Lambda_2) = 0$, and a flipped sign fails.
In that model the gauge parameters of ghost number one are $Q_B$-closed, so the test of $Q_B \Lambda_1$ and of part of $\hat{Q}_B \Lambda_2$ is vacuous.
The signs of the remaining rows are open, and the invariance beyond order $\gs$ at order $\hbar^0$ rests on the operator form of the relations and on the relations with two contractions stated at the head of this section; it has not been run.
